% Leçon 7 — Expression écrite : l'héritage — Chinois Tle A — Cours signé OnBuch+
% Programme officiel MINESEC. Compiler avec : tectonic ZH07-ecrit-l-heritage.tex
\newcommand{\DOCMATIERE}{Chinois}
\newcommand{\DOCNIVEAU}{Tle A}
\newcommand{\DOCMODULE}{Module 1 : Vie familiale et intégration sociale}
\newcommand{\DOCLECON}{ZH07}
\newcommand{\DOCTITRE}{Leçon 7 — Expression écrite : l'héritage}
% OnBuch+ — préambule « cours signés » ENRICHI (compilé avec tectonic / XeLaTeX)
% Système visuel « Pop » d'OnBuch+ : fond crème (jamais blanc), bordures encre
% épaisses, ombres « dures » décalées, titres Archivo Black en capitales.
% Version MATHÉMATIQUES : graphes (pgfplots/tikz), boîtes pédagogiques
% (définition, propriété, méthode, attention, le savais-tu, exercice résolu,
% exercice, TP, bilan), page de garde et signature OnBuch+ ancrée en bas.
%
% Macros attendues dans main.tex AVANT \input{preamble} :
%   \DOCMATIERE  \DOCNIVEAU  \DOCMODULE  \DOCLECON (numéro)  \DOCTITRE
\documentclass[11pt]{article}

\usepackage{fontspec}
\usepackage[french]{babel} % français = langue principale ; le chinois via \zh{..} / textebox (xeCJK)
\usepackage{xeCJK}
\usepackage{pifont} % \ding{51} \ding{55} utilisés par les modèles
\usepackage[a4paper,margin=2cm,top=3.4cm,bottom=2.8cm,headheight=1.4cm,footskip=1.3cm]{geometry}
\usepackage{amsmath,amssymb,amsfonts}
\usepackage{mathtools} % \xmapsto, \coloneqq... souvent utilisés par les modèles
\usepackage{cancel} % simplifications barrées (bilans de forces)
% \abs{z} (module, valeur absolue) et \norm{u} (norme), utilisés par les modèles
\providecommand{\abs}{}\let\abs\relax\DeclarePairedDelimiter{\abs}{\lvert}{\rvert}
\providecommand{\norm}{}\let\norm\relax\DeclarePairedDelimiter{\norm}{\lVert}{\rVert}
\usepackage[table]{xcolor} % [table] : fournit aussi \rowcolors (alternance de lignes)
\usepackage{pagecolor}
\usepackage{tikz}
\usetikzlibrary{arrows.meta,positioning,calc,shapes.geometric,shapes.misc,shapes.arrows,decorations.pathmorphing,decorations.pathreplacing,decorations.markings,patterns,shadows,mindmap,trees,fit,backgrounds,matrix,intersections,angles,quotes}
\usepackage{pgfplots}
\usepackage[version=4]{mhchem} % équations nucléaires \ce{^{235}_{92}U}
\usepackage[european,siunitx]{circuitikz} % schémas électriques
\pgfplotsset{compat=1.18}
\usepgfplotslibrary{fillbetween,groupplots} % aires entre courbes (calcul intégral)
\usepackage{enumitem}
\usepackage{fancyhdr}
% [most] charge toutes les bibliothèques tcolorbox (listings, minted, raster,
% documentation...) et gonfle inutilement la mémoire TeX sur un document long
% (~300 boîtes) : on ne charge que ce qui est réellement utilisé.
\usepackage[skins,breakable]{tcolorbox}
\usepackage{graphicx}
\usepackage{colortbl}
\usepackage{booktabs}
\usepackage{tabularx}
\usepackage{diagbox} % case d'en-tête diagonale des tableaux à double entrée
% Colonnes extensibles centrée / gauche / droite (C, L, R), souvent utilisées par les modèles
\newcolumntype{C}{>{\centering\arraybackslash}X}
\newcolumntype{L}{>{\raggedright\arraybackslash}X}
\newcolumntype{R}{>{\raggedleft\arraybackslash}X}
\usepackage{array}
\usepackage{multicol}
\usepackage{multirow}
\usepackage{siunitx}
% parse-numbers=false : accepte aussi \SI{2,5 \times 10^{-3}}{...}
\sisetup{locale=FR,output-decimal-marker={,},group-minimum-digits=4,parse-numbers=false}
\DeclareSIUnit\ohm{\ensuremath{\symup{\Omega}}} % U+2126 absent de la police
\DeclareSIUnit\division{div} % réglages d'oscilloscope (ms/div, V/div)
\DeclareSIUnit\tr{tr} % tours (vitesse de rotation, ex. \SI{1500}{\tr\per\minute})
\DeclareSIUnit\molar{M} % concentration molaire (ex. \SI{3}{\molar}, \SI{1}{\micro\molar})
\DeclareSIUnit\an{an} % année (le modèle confond parfois avec \year, un compteur TeX) — ex. \SI{5}{\kilo\meter\per\an}
\DeclareSIUnit\FCFA{FCFA} % franc CFA (ex. \SI{500}{\FCFA\per\kilo\gram})
\DeclareSIUnit\degreeBrix{\textdegree Brix} % teneur en sucre d'un sirop/jus (réfractométrie)
\DeclareSIUnit\month{mois} % le modèle utilise parfois \month, un compteur TeX, comme unité de durée
\usepackage{qrcode}
% \degree : commande fréquemment utilisée par les modèles pour un angle ou une
% température en dehors de \SI{}{} (non fournie par les paquets chargés).
\providecommand{\degree}{^{\circ}}
% \qed (fin de démonstration), utilisé par les modèles sans amsthm
\providecommand{\qed}{\hfill$\square$}
% \barre : double barre des valeurs interdites dans un tableau de variations (tkz-tab)
\providecommand{\barre}{\ensuremath{\|}}
% environnement proof (amsthm non chargé) utilisé par les modèles
\ifdefined\proof\else\newenvironment{proof}[1][Démonstration]{\par\noindent\textit{#1.}\ }{\qed\par}\fi
\usepackage{lastpage}
\usepackage{hyperref}
\hypersetup{hidelinks}

% ============================================================
% Palette « Pop » OnBuch+ — mêmes hex que la classe P (plus_ui.dart)
% ============================================================
\definecolor{poporange}{HTML}{F59321}
\definecolor{popdark}{HTML}{E07A0C}
\definecolor{popcream}{HTML}{FFF6E8}
\definecolor{popink}{HTML}{1C1714}
\definecolor{poppurple}{HTML}{7A5AE0}
\definecolor{popgreen}{HTML}{2E9E6A}
\definecolor{poppink}{HTML}{E0457B}
\definecolor{popblue}{HTML}{2F7DE1}
\definecolor{popgold}{HTML}{C98A12}
\definecolor{popmuted}{HTML}{8A7F76}
\definecolor{popline}{HTML}{EADFD2}
\definecolor{popgray}{HTML}{8A7F76}
\colorlet{popgrey}{popgray}
% Alias tolérants (le modèle nomme parfois les couleurs différemment)
\colorlet{popwhite}{white}
\colorlet{popblack}{popink}
\colorlet{popred}{poppink}
\colorlet{popyellow}{popgold}
% Teintes claires (fonds de tableaux/encadrés) — le modèle nomme parfois
% les couleurs avec un suffixe L (ex. popblueL) sans les avoir définies.
\colorlet{poporangeL}{poporange!20}
\colorlet{popdarkL}{popdark!20}
\colorlet{poppurpleL}{poppurple!20}
\colorlet{popgreenL}{popgreen!20}
\colorlet{poppinkL}{poppink!20}
\colorlet{popblueL}{popblue!20}
\colorlet{popgoldL}{popgold!20}
\colorlet{popredL}{popred!20}
\colorlet{popgrayL}{popgray!20}
% Teintes claires pour fonds de tableaux / zones
\colorlet{popblueL}{popblue!10!white}
\colorlet{poporangeL}{poporange!14!white}
\colorlet{popgreenL}{popgreen!12!white}
\colorlet{poppurpleL}{poppurple!12!white}
\colorlet{poppinkL}{poppink!12!white}
\colorlet{popgoldL}{popgold!14!white}

\pagecolor{popcream}

% ============================================================
% Typographie
% ============================================================
\defaultfontfeatures{Ligatures=TeX}
\setmainfont{PlusJakartaSans-Medium.ttf}[Path=../../../fonts/,BoldFont=PlusJakartaSans-Bold.ttf]
% Symboles, API et emojis absents de la police principale : repli sur DejaVu Sans (caractères actifs)
\newfontface\symfont{DejaVuSans.ttf}[Path=../../../fonts/]
\makeatletter
\newcommand\symactive[1]{\catcode#1=13 \begingroup\lccode`\~=#1 \lowercase{\endgroup\def~}{{\symfont\char#1}}}
\newcommand\symblank[1]{\catcode#1=13 \begingroup\lccode`\~=#1 \lowercase{\endgroup\def~}{}}
\newcommand\symhyph[1]{\catcode#1=13 \begingroup\lccode`\~=#1 \lowercase{\endgroup\def~}{\mbox{-}}}
\newcount\symc
\newcommand\symrange[2]{\symc=#1 \@whilenum\symc<#2 \do{\expandafter\symactive\expandafter{\the\symc}\advance\symc by 1 }}
\newcommand\symblankrange[2]{\symc=#1 \@whilenum\symc<#2 \do{\expandafter\symblank\expandafter{\the\symc}\advance\symc by 1 }}
\makeatother
\symrange{"0250}{"02FF}\symactive{"2713}\symactive{"2714}\symactive{"2717}\symactive{"2718}\symactive{"2610}\symactive{"2611}\symactive{"2612}
\symrange{"2600}{"27BF}
\catcode"2705=13 \begingroup\lccode`\~="2705 \lowercase{\endgroup\def~}{{\symfont\char"2713}}\symblank{"26BD}\symblank{"26A1}
\symrange{"2460}{"257F}\symactive{"223C}\symactive{"2248}\symactive{"2260}
\symhyph{"2011}\symhyph{"2E1A}\symblankrange{"1F300}{"1FAFF}\symblank{"FE0F}
% Caractères chinois : WenQuanYi Zen Hei (sous-ensemble GB2312 + pinyin), gras/italique simulés
\setCJKmainfont{WenQuanYiZenHei-subset.ttf}[Path=../../../fonts/,AutoFakeBold=3,AutoFakeSlant=0.2]
\xeCJKsetup{CJKspace=false}
% Pinyin : voyelles à caron (ǎ ǐ ǒ ǔ ǖ ǘ ǚ ǜ) absentes de la police principale -> caractères actifs
\newfontface\pyfont{WenQuanYiZenHei-subset.ttf}[Path=../../../fonts/]
\newcommand\pyactive[1]{\catcode#1=13 \begingroup\lccode`\~=#1 \lowercase{\endgroup\def~}{{\pyfont\char#1}}}
\pyactive{"01CD}\pyactive{"01CE}\pyactive{"01CF}\pyactive{"01D0}\pyactive{"01D1}\pyactive{"01D2}\pyactive{"01D3}\pyactive{"01D4}\pyactive{"01D5}\pyactive{"01D6}\pyactive{"01D7}\pyactive{"01D8}\pyactive{"01D9}\pyactive{"01DA}\pyactive{"01DB}\pyactive{"01DC}
% Maths en sans-serif (Fira Math) avec chiffres et lettres droites de Plus
% Jakarta, pour que les formules se fondent dans le texte.
\usepackage[math-style=ISO,bold-style=ISO]{unicode-math}
\setmathfont{FiraMath-Regular.otf}[Path=../../../fonts/]
\setmathfont{PlusJakartaSans-Medium.ttf}[Path=../../../fonts/,range={up/num,up/{Latin,latin},\mathup}]
\usepackage{icomma} % virgule décimale sans espace en mode math
% Caractères mathématiques Unicode absents des polices de texte (Plus Jakarta,
% Archivo Black) : rendus par leur équivalent mathématique, en texte comme en
% formule — sinon ils s'affichent en carré vide (ex. « Arithmétique dans ℤ »).
\usepackage{newunicodechar}
\newunicodechar{ℝ}{\ensuremath{\mathbb{R}}}
\newunicodechar{ℤ}{\ensuremath{\mathbb{Z}}}
\newunicodechar{ℂ}{\ensuremath{\mathbb{C}}}
\newunicodechar{ℕ}{\ensuremath{\mathbb{N}}}
\newunicodechar{ℚ}{\ensuremath{\mathbb{Q}}}
\newunicodechar{𝔻}{\ensuremath{\mathbb{D}}}
\newunicodechar{α}{\ensuremath{\alpha}}
\newunicodechar{β}{\ensuremath{\beta}}
\newunicodechar{γ}{\ensuremath{\gamma}}
\newunicodechar{δ}{\ensuremath{\delta}}
\newunicodechar{ε}{\ensuremath{\varepsilon}}
\newunicodechar{θ}{\ensuremath{\theta}}
\newunicodechar{λ}{\ensuremath{\lambda}}
\newunicodechar{μ}{\ensuremath{\mu}}
\newunicodechar{σ}{\ensuremath{\sigma}}
\newunicodechar{τ}{\ensuremath{\tau}}
\newunicodechar{φ}{\ensuremath{\varphi}}
\newunicodechar{ψ}{\ensuremath{\psi}}
\newunicodechar{ω}{\ensuremath{\omega}}
\newunicodechar{Σ}{\ensuremath{\Sigma}}
% majuscules produites par \MakeUppercase dans les titres (α-aminés -> Α)
\newunicodechar{Α}{\ensuremath{\alpha}}
\newunicodechar{Β}{\ensuremath{\beta}}
\newunicodechar{Γ}{\ensuremath{\Gamma}}
\newunicodechar{Δ}{\ensuremath{\Delta}}
\newunicodechar{Ε}{\ensuremath{\varepsilon}}
\newunicodechar{Θ}{\ensuremath{\Theta}}
\newunicodechar{Λ}{\ensuremath{\Lambda}}
\newunicodechar{Μ}{\ensuremath{\mu}}
\newunicodechar{Τ}{\ensuremath{\tau}}
\newunicodechar{Φ}{\ensuremath{\Phi}}
\newunicodechar{Ψ}{\ensuremath{\Psi}}
\newunicodechar{Ω}{\ensuremath{\Omega}}
\newunicodechar{↦}{\ensuremath{\mapsto}}
\newunicodechar{∈}{\ensuremath{\in}}
\newunicodechar{∉}{\ensuremath{\notin}}
\newunicodechar{∧}{\ensuremath{\wedge}}
\newunicodechar{⇔}{\ensuremath{\Leftrightarrow}}
\newunicodechar{⟺}{\ensuremath{\Longleftrightarrow}}
\newunicodechar{⇒}{\ensuremath{\Rightarrow}}
\newunicodechar{⟹}{\ensuremath{\Longrightarrow}}
\newunicodechar{∘}{\ensuremath{\circ}}
\newunicodechar{∪}{\ensuremath{\cup}}
\newunicodechar{∩}{\ensuremath{\cap}}
\newunicodechar{≡}{\ensuremath{\equiv}}
\newunicodechar{⊂}{\ensuremath{\subset}}
\newunicodechar{⊥}{\ensuremath{\perp}}
\newunicodechar{∃}{\ensuremath{\exists}}
\newunicodechar{∀}{\ensuremath{\forall}}
\newunicodechar{∛}{\ensuremath{\sqrt[3]{\,}}}
\newunicodechar{∜}{\ensuremath{\sqrt[4]{\,}}}
\newunicodechar{⃗}{\ensuremath{{}^{\to}}}
\newunicodechar{ⁿ}{\ifmmode^{n}\else\textsuperscript{\ensuremath{n}}\fi}
\newunicodechar{ˣ}{\ifmmode^{x}\else\textsuperscript{\ensuremath{x}}\fi}
\newunicodechar{ᵇ}{\ifmmode^{b}\else\textsuperscript{\ensuremath{b}}\fi}
\newunicodechar{ʳ}{\ifmmode^{r}\else\textsuperscript{\ensuremath{r}}\fi}
\newunicodechar{ᵘ}{\ifmmode^{u}\else\textsuperscript{\ensuremath{u}}\fi}
\newunicodechar{ʸ}{\ifmmode^{y}\else\textsuperscript{\ensuremath{y}}\fi}
\newunicodechar{ᵃ}{\ifmmode^{a}\else\textsuperscript{\ensuremath{a}}\fi}
\newunicodechar{ᵉ}{\ifmmode^{e}\else\textsuperscript{\ensuremath{e}}\fi}
\newunicodechar{ᵏ}{\ifmmode^{k}\else\textsuperscript{\ensuremath{k}}\fi}
\newunicodechar{ᵖ}{\ifmmode^{p}\else\textsuperscript{\ensuremath{p}}\fi}
\newunicodechar{ᴺ}{\ifmmode^{N}\else\textsuperscript{\ensuremath{N}}\fi}
\newunicodechar{ᶜ}{\ifmmode^{c}\else\textsuperscript{\ensuremath{c}}\fi}
\newunicodechar{ʰ}{\ifmmode^{h}\else\textsuperscript{\ensuremath{h}}\fi}
\newunicodechar{ᵠ}{\ifmmode^{q}\else\textsuperscript{\ensuremath{q}}\fi}
\newunicodechar{ₙ}{\ifmmode_{n}\else\textsubscript{\ensuremath{n}}\fi}
\newunicodechar{ₐ}{\ifmmode_{a}\else\textsubscript{\ensuremath{a}}\fi}
\newunicodechar{ᵢ}{\ifmmode_{i}\else\textsubscript{\ensuremath{i}}\fi}
\newunicodechar{ᵣ}{\ifmmode_{r}\else\textsubscript{\ensuremath{r}}\fi}
\newunicodechar{ₖ}{\ifmmode_{k}\else\textsubscript{\ensuremath{k}}\fi}
\newunicodechar{ᵦ}{\ifmmode_{\beta}\else\textsubscript{\ensuremath{\beta}}\fi}
\newunicodechar{ₓ}{\ifmmode_{x}\else\textsubscript{\ensuremath{x}}\fi}
\newfontfamily\popdisplay{ArchivoBlack-Regular.ttf}[Path=../../../fonts/]
\newfontfamily\poplogo{SpaceGrotesk-Bold.ttf}[Path=../../../fonts/]
\newfontfamily\popsemi{PlusJakartaSans-SemiBold.ttf}[Path=../../../fonts/]
\newfontfamily\popxbold{PlusJakartaSans-ExtraBold.ttf}[Path=../../../fonts/]
\newfontfamily\popxboldls{PlusJakartaSans-ExtraBold.ttf}[Path=../../../fonts/,LetterSpace=3.0]

\setlist[enumerate]{leftmargin=1.7em,itemsep=5pt,topsep=4pt}
\setlist[itemize]{leftmargin=1.7em,itemsep=4pt,topsep=4pt,label={\color{poporange}\textbullet}}
\setlength{\parindent}{0pt}
\setlength{\parskip}{5pt}
\renewcommand{\arraystretch}{1.25}

% pgfplots : style Pop par défaut
\pgfplotsset{
  every axis/.append style={axis line style={popink,line width=1pt},tick style={popink},
    grid style={popline,line width=0.5pt},label style={font=\small},tick label style={font=\footnotesize}},
  popcurve/.style={poporange,line width=1.8pt,no markers,smooth},
}
% Même style disponible dans un \draw TikZ simple (hors pgfplots)
\tikzset{popcurve/.style={poporange,line width=1.8pt}}
\tikzset{popgrid/.style={popline,very thin}}
\colorlet{popcurve}{poporange} % le nom du style est parfois utilisé comme couleur

% ============================================================
% Pill / squiggle
% ============================================================
\newcommand{\poppill}[3]{%
  \tikz[baseline=-0.55ex]\node[draw=popink,line width=1.2pt,rounded corners=2.6pt,%
    fill=#2,inner xsep=6.5pt,inner ysep=3pt]{%
    \popxboldls\fontsize{7}{7}\selectfont\color{#3}\MakeUppercase{#1}};%
}
\newcommand{\popsquiggle}[1]{%
  \tikz[baseline=-0.4ex]{
    \draw[#1,line width=2.6pt,line cap=round]
      plot [smooth] coordinates {(0,0.09) (0.34,0.01) (0.68,0.21) (1.02,0.01) (1.36,0.10)};
  }%
}
% Mot-clé surligné (marqueur orange sous le texte)
\newcommand{\cle}[1]{{\popxbold\color{popdark}#1}}

% ============================================================
% En-tête / pied
% ============================================================
\pagestyle{fancy}
\fancyhf{}
\renewcommand{\headrulewidth}{0pt}
\renewcommand{\footrulewidth}{0pt}
\lhead{\raisebox{-0.15ex}{\poplogo\fontsize{13}{13}\selectfont\color{popink}OnBuch\color{poporange}+}}
\rhead{\poppill{\DOCMATIERE\ \textbullet\ \DOCNIVEAU\ \textbullet\ Leçon \DOCLECON}{white}{popink}}
\lfoot{\popxbold\fontsize{7.6}{7.6}\selectfont\color{popmuted}\MakeUppercase{Cours signé OnBuch+}\ \textbullet\ {\popsemi\DOCTITRE}}
\rfoot{\tikz[baseline=-0.6ex]\node[rounded corners=8.5pt,fill=popink,minimum height=17pt,minimum width=17pt,inner xsep=5pt,inner sep=0]{\popxbold\fontsize{8}{8}\selectfont\color{popcream}\thepage\,/\,\pageref*{LastPage}};}

% ============================================================
% Titres
% ============================================================
\newcommand{\chaptitre}[2]{%
  \begin{tcolorbox}[enhanced,colback=poporange,colframe=popink,boxrule=2.6pt,arc=11pt,
      drop shadow=popink,left=15pt,right=15pt,top=12pt,bottom=11pt,before skip=3pt,after skip=18pt]
    \poppill{#2}{popink}{popcream}\par\vspace{9pt}
    {\raggedright\hyphenpenalty=10000\exhyphenpenalty=10000\popdisplay\fontsize{22}{24}\selectfont\color{popcream}\MakeUppercase{#1}\par}
  \end{tcolorbox}
}
% Section numérotée : pastille orange avec le numéro + titre Archivo
\newcounter{popsec}
\newcommand{\coursec}[1]{%
  \stepcounter{popsec}\setcounter{popsub}{0}%
  \par\vspace{16pt}\noindent
  \begin{minipage}{\linewidth}
  \tikz[baseline=-0.7ex]\node[draw=popink,line width=1.6pt,fill=poporange,rounded corners=4pt,minimum size=22pt,inner sep=0]{\popdisplay\fontsize{12}{12}\selectfont\color{popcream}\thepopsec};%
  \hspace{8pt}{\popdisplay\fontsize{15}{17}\selectfont\color{popink}\MakeUppercase{#1}}\par
  \vspace{4pt}\noindent{\color{poporange}\rule{36pt}{3.6pt}}
  \end{minipage}\par\nopagebreak\vspace{8pt}
}
\newcounter{popsub}
\newcommand{\courssub}[1]{%
  \stepcounter{popsub}%
  \par\vspace{9pt}\noindent{\popxbold\fontsize{11.5}{13}\selectfont\color{popink}{\color{poporange}\thepopsec.\thepopsub}\hspace{6pt}#1}\par\nopagebreak\vspace{4pt}
}

% ============================================================
% Boîtes pédagogiques — même recette : fond blanc, bordure épaisse colorée,
% ombre dure encre, badge plein. Titre optionnel [#1].
% ============================================================
\tcbset{popbase/.style={breakable,enhanced,colback=white,boxrule=2.3pt,arc=9pt,
  shadow={3pt}{-3pt}{0pt}{popink},left=14pt,right=14pt,top=11pt,bottom=11pt,before skip=11pt,after skip=15pt,
  fontupper=\popsemi\fontsize{10.6}{14.5}\selectfont\color{popink},
  fontlower=\popsemi\fontsize{10.6}{14.5}\selectfont\color{popink}}}
% Coche dessinée (le glyphe ✓ n'existe pas dans les polices du document)
\newcommand{\popcheck}[1]{\tikz[baseline=-0.5ex]\draw[#1,line width=1.6pt,line cap=round,line join=round](-0.13,0.0)--(-0.04,-0.09)--(0.13,0.1);}
\providecommand{\coche}{\popcheck{popgreen}}
\providecommand{\resultat}[1]{\par\smallskip\noindent\fcolorbox{popgreen}{popgreenL}{\parbox{\dimexpr\linewidth-2\fboxsep-2\fboxrule\relax}{\textbf{Résultat.} #1}}\par\smallskip}
\newcommand{\popboxhead}[3]{\poppill{#1}{#2}{white}\ifx\relax#3\relax\else\hspace{6pt}{\popxbold\fontsize{10.6}{12}\selectfont\color{popink}#3}\fi\par\vspace{7pt}}

\newtcolorbox{aretenir}[1][]{popbase,colframe=popgreen,before upper={\popboxhead{À retenir}{popgreen}{#1}}}
% Texte en chinois (document de lecture, dialogue, extrait) : cadre
% d'étude. Un titre optionnel (source, auteur) s'affiche dans l'en-tête.
\newtcolorbox{textebox}[1][]{popbase,colframe=popblue,before upper={\popboxhead{文本 · Texte}{popblue}{#1}},after upper={}}
% Marques ✓ / ✗ (forme correcte / fautive) souvent utilisées par les modèles
\providecommand{\cross}{\textcolor{poppink}{\ensuremath{\times}}}
\providecommand{\xmark}{\cross}\providecommand{\crossmark}{\cross}\providecommand{\cmark}{\ensuremath{\checkmark}}
% Ligne de réponse / justification à compléter (inventée par les modèles)
\providecommand{\justification}[1]{\par\noindent\textit{Justification~:} \dotfill\par}
% \textipa (package tipa non chargé) : la police ne couvre pas l'API -> simple passage
\providecommand{\textipa}[1]{#1}
% Soulignement (ulem non chargé) : \uline, \ul
\providecommand{\uline}[1]{\underline{#1}}\providecommand{\ul}[1]{\underline{#1}}
% \glossaire{\item ...} inventée par les modèles : liste de glossaire
\providecommand{\glossaire}[1]{\begin{itemize}#1\end{itemize}}
% \alert (beamer) inventée par les modèles : texte mis en évidence
\providecommand{\alert}[1]{\textbf{\textcolor{poppink}{#1}}}
% variantes ASCII d'ordinaux écrites en macro
\providecommand{\iere}{\textsuperscript{re}}\providecommand{\ieme}{\textsuperscript{e}}\providecommand{\ere}{\textsuperscript{re}}\providecommand{\eme}{\textsuperscript{e}}\providecommand{\er}{\textsuperscript{er}}\providecommand{\ie}{\textsuperscript{e}}
% Ordinaux français écrits en macro par les modèles : 1\ière, 3\ième
\newcommand{\ière}{\textsuperscript{re}}\newcommand{\ième}{\textsuperscript{e}}
% \exemple (inventée par les modèles) : étiquette « Exemple : »
\providecommand{\exemple}{\textbf{Exemple~:}\ }
% \square utilisable aussi hors mode math (cases à cocher)
\AtBeginDocument{\let\squareMath\square\renewcommand{\square}{\ensuremath{\squareMath}}}
% Mot ou phrase en chinois à l'intérieur d'un paragraphe français
\newcommand{\zh}[1]{\ifmmode\text{#1}\else{#1}\fi}
\let\eng\zh \let\esp\zh \let\all\zh % alias tolérés
% flèches d'intonation inventées par les modèles
\providecommand{\textsearrow}{}\renewcommand{\textsearrow}{\ensuremath{\searrow}}\providecommand{\textnearrow}{}\renewcommand{\textnearrow}{\ensuremath{\nearrow}}
% \fr{..} : mot français (inventé par les modèles)
\providecommand{\fr}[1]{\textit{#1}}
% zhbox : encadré de texte chinois inventé par les modèles
\newenvironment{zhbox}{\begin{quote}}{\end{quote}}
% \vs : « versus » inventé par les modèles
\providecommand{\vs}{\textit{vs}\ }
% \blank : case à compléter inventée par les modèles
\providecommand{\blank}[1][]{\underline{\hspace{1.6em}}}
\newtcolorbox{exemplebox}[1][]{popbase,colframe=poporange,before upper={\popboxhead{Exemple}{poporange}{#1}}}
\newtcolorbox{definition}[1][]{popbase,colframe=popblue,before upper={\popboxhead{Définition}{popblue}{#1}}}
\newtcolorbox{propriete}[1][]{popbase,colframe=poppurple,before upper={\popboxhead{Propriété}{poppurple}{#1}}}
\newtcolorbox{methode}[1][]{popbase,colframe=popgold,before upper={\popboxhead{Méthode}{popgold}{#1}}}
\newtcolorbox{attention}[1][]{popbase,colframe=poppink,before upper={\popboxhead{Attention}{poppink}{#1}}}
\newtcolorbox{experience}[1][]{popbase,colframe=popblue,colback=popblueL,before upper={\popboxhead{Expérience}{popblue}{#1}}}
% « Le savais-tu ? » : carte violette pleine (contexte, culture, Cameroun)
\newtcolorbox{savaistu}[1][]{popbase,colback=poppurple,colframe=popink,
  fontupper=\popsemi\fontsize{10.4}{14.2}\selectfont\color{white},
  before upper={\poppill{Le savais-tu\,?}{white}{poppurple}\ifx\relax#1\relax\else\hspace{6pt}{\popxbold\color{white}#1}\fi\par\vspace{7pt}}}
% Exercice résolu : énoncé en haut, \tcblower puis la solution sur fond crème
\newcounter{popexo}\newcounter{popexr}
\newtcolorbox{exoresolu}[1][]{popbase,colframe=poporange,colbacklower=popcream,
  segmentation style={popink,line width=1.2pt,dashed},
  before upper={\stepcounter{popexr}\popboxhead{Exercice résolu \thepopexr}{poporange}{#1}},
  before lower={\poppill{Solution}{popink}{popcream}\par\vspace{6pt}}}
% Exercice (non résolu en place) — la correction est dans la partie Corrigés
\newtcolorbox{exercice}[1][]{popbase,colframe=popink,
  before upper={\stepcounter{popexo}\popboxhead{Exercice \thepopexo}{popink}{#1}}}
% Corrigé
\newtcolorbox{corrige}[1][]{popbase,colframe=popgreen,colback=popgreenL,
  before upper={\popboxhead{Corrigé}{popgreen}{#1}}}
% Objectifs / prérequis / bilan
\newtcolorbox{objectifs}{popbase,colframe=popink,colback=poporangeL,
  before upper={\popboxhead{Objectifs de la leçon}{popink}{}}}
\newtcolorbox{prerequis}{popbase,colframe=popink,colback=popblueL,
  before upper={\popboxhead{Prérequis}{popblue}{}}}
\newtcolorbox{bilan}{popbase,colframe=popink,colback=white,boxrule=2.8pt,
  before upper={\popboxhead{Fiche bilan}{poporange}{L'essentiel à connaître pour l'examen}}}
% Figure encadrée avec légende
\newtcolorbox{popfigure}[1][]{enhanced,breakable=false,colback=white,colframe=popline,boxrule=1.4pt,arc=8pt,
  left=10pt,right=10pt,top=10pt,bottom=8pt,before skip=10pt,after skip=12pt,halign=center,
  lower separated=false,halign lower=center,
  fontlower=\popsemi\fontsize{9}{11.5}\selectfont\color{popmuted},
  #1}
\newcommand{\legende}[1]{\par\vspace{6pt}{\popsemi\fontsize{9}{11.5}\selectfont\color{popmuted}#1}}

% ============================================================
% Signature OnBuch+
% ============================================================
\newcommand{\poplogoinline}[1]{{\poplogo\fontsize{#1}{#1}\selectfont\color{popink}OnBuch\color{poporange}+}}
\newcommand{\signatureonbuch}{%
  \begin{tcolorbox}[enhanced,colback=popink,colframe=popink,boxrule=2.2pt,arc=10pt,
      drop shadow=poporange,left=14pt,right=14pt,top=10pt,bottom=10pt,before skip=0pt,after skip=0pt]
    \tikz[baseline=-0.7ex]\node[circle,fill=popgreen,draw=popcream,line width=1.3pt,minimum size=22pt,inner sep=0]{\popcheck{white}};%
    \hspace{9pt}\begin{minipage}[c]{0.62\linewidth}
      {\popxbold\fontsize{12}{13}\selectfont\color{popcream}Signé\ \textbullet\ {\poplogo OnBuch\color{poporange}+}}\par\vspace{2pt}
      {\raggedright\popsemi\fontsize{8}{10.5}\selectfont\color{popline}\MakeUppercase{Équipe pédagogique OnBuch+}\\\MakeUppercase{Conforme au programme officiel MINESEC}\par}
    \end{minipage}\hfill
    {\poplogo\fontsize{22}{22}\selectfont\color{popcream}OnBuch\color{poporange}+}
  \end{tcolorbox}%
}

% ============================================================
% Page de garde
% #1 titre · #2 matière · #3 niveau · #4 module · #5 n° leçon · #6 durée ·
% #7 description / objectifs (texte ou itemize)
% ============================================================
\newcommand{\pagegarde}[7]{%
  \thispagestyle{empty}
  \newgeometry{a4paper,margin=1.7cm,top=1.6cm,bottom=1.4cm}%
  \begin{tikzpicture}[remember picture,overlay]
    \fill[poporange!18] ([xshift=-3.2cm,yshift=-3.6cm]current page.north east) circle (4.6cm);
    \fill[poppurple!14] ([xshift=2.4cm,yshift=7.6cm]current page.south west) circle (3.8cm);
  \end{tikzpicture}%
  {\poplogo\fontsize{30}{30}\selectfont\color{popink}OnBuch\color{poporange}+}\hfill
  \raisebox{0.6ex}{\poppill{Cours signé \textbullet\ Premium}{popink}{popcream}}\par
  \vspace{1pt}\popsquiggle{poppurple}\par
  \vspace{8pt}
  \begin{tcolorbox}[enhanced,colback=poporange,colframe=popink,boxrule=2.8pt,arc=13pt,
      drop shadow=popink,left=18pt,right=18pt,top=12pt,bottom=13pt,after skip=10pt]
    \poppill{#2 \textbullet\ #3}{popink}{popcream}\hspace{5pt}\poppill{#4}{white}{popink}\par\vspace{8pt}
    {\popdisplay\fontsize{14}{14}\selectfont\color{popink}LEÇON #5}\par\vspace{4pt}
    {\raggedright\hyphenpenalty=10000\exhyphenpenalty=10000\popdisplay\fontsize{25}{27}\selectfont\color{popcream}\MakeUppercase{#1}\par}
    \vspace{8pt}
    {\popxbold\fontsize{9.5}{9.5}\selectfont\color{popink}\MakeUppercase{Durée conseillée : #6}}
  \end{tcolorbox}
  \begin{tcolorbox}[enhanced,colback=white,colframe=popink,boxrule=2.2pt,arc=10pt,
      drop shadow=popink,left=16pt,right=16pt,top=10pt,bottom=10pt,after skip=8pt]
    \poppill{Dans cette leçon}{poppurple}{white}\par\vspace{5pt}
    \popsemi\fontsize{9.3}{12.6}\selectfont\color{popink}#7
  \end{tcolorbox}
  \noindent\begin{minipage}[t]{0.487\linewidth}
    \begin{tcolorbox}[enhanced,colback=popgreen,colframe=popink,boxrule=2.2pt,arc=10pt,
        drop shadow=popink,left=11pt,right=11pt,top=10pt,bottom=10pt,height=3.1cm,valign=center]
      \tcbox[colback=white,colframe=popink,boxrule=1.3pt,arc=5pt,boxsep=0pt,nobeforeafter,
        left=3pt,right=3pt,top=3pt,bottom=3pt,valign=center]{\qrcode[height=1.9cm]{https://play.google.com/store/apps/details?id=com.luvvix.onbuch&hl=fr}}%
      \hspace{8pt}\begin{minipage}[c]{0.5\linewidth}
        {\popdisplay\fontsize{11}{12.5}\selectfont\color{white}\MakeUppercase{Télécharge OnBuch}}\par\vspace{4pt}
        {\raggedright\popsemi\fontsize{8.4}{11}\selectfont\color{white}Gratuit \textbullet\ Google Play\\Tuteur IA, annales, concours, résultats\par}
      \end{minipage}
    \end{tcolorbox}
  \end{minipage}\hfill
  \begin{minipage}[t]{0.487\linewidth}
    \begin{tcolorbox}[enhanced,colback=poppurple,colframe=popink,boxrule=2.2pt,arc=10pt,
        drop shadow=popink,left=11pt,right=11pt,top=10pt,bottom=10pt,height=3.1cm,valign=center]
      \tikz[baseline=-0.7ex]\node[circle,fill=white,draw=popink,line width=1.3pt,minimum size=1.6cm,inner sep=0]{\popdisplay\fontsize{24}{24}\selectfont\color{poppurple}+};%
      \hspace{8pt}\begin{minipage}[c]{0.6\linewidth}
        {\popdisplay\fontsize{11}{12.5}\selectfont\color{white}\MakeUppercase{Passe à OnBuch+}}\par\vspace{4pt}
        {\raggedright\popsemi\fontsize{8.4}{11}\selectfont\color{white}Tous les cours signés, Tuteur IA illimité, fiches PDF — onglet OnBuch+ de l'app\par}
      \end{minipage}
    \end{tcolorbox}
  \end{minipage}
  \vspace{8pt}
  \popcontenu
  \vfill
  \signatureonbuch
  \restoregeometry
  \newpage
}

% Rangée « Ce cours contient » (page de garde)
\newcommand{\poptile}[3]{% #1 couleur #2 gros texte #3 légende
  \begin{tcolorbox}[enhanced,colback=#1,colframe=popink,boxrule=2pt,arc=9pt,shadow={2.5pt}{-2.5pt}{0pt}{popink},
      left=6pt,right=6pt,top=9pt,bottom=9pt,halign=center,valign=center,height=1.75cm,width=\linewidth,nobeforeafter]
    {\popdisplay\fontsize{12.5}{13}\selectfont\color{white}\MakeUppercase{#2}}\par\vspace{3pt}
    {\centering\popsemi\fontsize{7.8}{9.5}\selectfont\color{white}#3\par}
  \end{tcolorbox}}
\newcommand{\popcontenu}{%
  \par\noindent{\popxboldls\fontsize{7.5}{7.5}\selectfont\color{popmuted}CE COURS SIGNÉ CONTIENT}\par\vspace{7pt}
  \noindent\begin{minipage}[t]{0.235\linewidth}\poptile{popblue}{Cours}{complet, illustré, pas à pas}\end{minipage}\hfill
  \begin{minipage}[t]{0.235\linewidth}\poptile{poporange}{Méthodes}{et exercices résolus}\end{minipage}\hfill
  \begin{minipage}[t]{0.235\linewidth}\poptile{poppink}{Exercices}{type examen + corrigés}\end{minipage}\hfill
  \begin{minipage}[t]{0.235\linewidth}\poptile{popgreen}{Fiche bilan}{l'essentiel à retenir}\end{minipage}\par}

% Sommaire
\newtcolorbox{sommaire}{enhanced,colback=white,colframe=popink,boxrule=2.2pt,arc=10pt,
  drop shadow=popink,left=18pt,right=18pt,top=13pt,bottom=13pt,before skip=6pt,after skip=20pt,
  before upper={\poppill{Sommaire}{poppurple}{white}\par\vspace{9pt}\popsemi\fontsize{10.6}{14.5}\selectfont\color{popink}}}

% Signature de fin de cours — poussée tout en bas de la dernière page
\newcommand{\findecours}{%
  \par\vfill\nopagebreak
  \begin{center}\popsquiggle{poporange}\end{center}\vspace{4pt}
  \signatureonbuch
}
\AtBeginDocument{\def\llbracket{\mathopen{[\mkern-3.5mu[}}\def\rrbracket{\mathclose{]\mkern-3.5mu]}}} % intervalles d'entiers (glyphe ⟦ absent de la police)
% Glyphes absents de Fira Math : redéfinis à partir de glyphes disponibles
\AtBeginDocument{%
  \def\setminus{\mathbin{\backslash}}\let\smallsetminus\setminus
  \def\vdots{\mathord{\vbox{\baselineskip3.2pt\lineskiplimit0pt\kern4pt\hbox{.}\hbox{.}\hbox{.}}}}%
  \def\ddots{\mathinner{\mkern1mu\raise6.5pt\hbox{.}\mkern2mu\raise3.6pt\hbox{.}\mkern2mu\raise0.7pt\hbox{.}\mkern1mu}}%
  \def\triangle{\mathord{\scalebox{0.9}{$\symup{\Delta}$}}}%
  \def\nabla{\mathord{\raisebox{\depth}{\scalebox{1}[-1]{$\symup{\Delta}$}}}}%
  \def\checkmark{\popcheck{popdark}}%
  \def\star{\mathbin{\ast}}\def\bigstar{\mathord{\ast}}%
  \def\diamond{\mathbin{\scalebox{0.6}{$\Diamond$}}}%
  \def\bigcup{\mathop{\mathchoice{\raisebox{-0.3em}{\scalebox{1.7}{$\cup$}}}{\scalebox{1.25}{$\cup$}}{\cup}{\cup}}\displaylimits}%
  \def\bigcap{\mathop{\mathchoice{\raisebox{-0.3em}{\scalebox{1.7}{$\cap$}}}{\scalebox{1.25}{$\cap$}}{\cap}{\cap}}\displaylimits}%
  \def\lesseqqgtr{\mathrel{\lessgtr}}\def\bigoplus{\mathop{\oplus}\displaylimits}%
}
\providecommand{\ssi}{si et seulement si} % abréviation fréquente
\AtBeginDocument{\providecommand{\wideparen}{\overparen}} % arc de cercle

%% FIGFIX-BEGIN v4 — correctifs globaux des figures (docs/onbuch-generation/tools/figfix)
\usepackage{adjustbox}\usepackage{etoolbox}
% toute figure TikZ/pgfplots plus large que la ligne est réduite à la largeur disponible
\BeforeBeginEnvironment{tikzpicture}{\begin{adjustbox}{max width=\linewidth}}
\AfterEndEnvironment{tikzpicture}{\end{adjustbox}}
% légendes pgfplots lisibles par-dessus les courbes
\pgfplotsset{every axis/.append style={legend style={fill=white,fill opacity=0.92,text opacity=1,draw=black!15,inner sep=3pt}}}
% étiquettes d'arêtes (syntaxe edge["..."]) avec fond blanc
\tikzset{every edge quotes/.append style={fill=white,inner sep=1.2pt}}
%% FIGFIX-END
\begin{document}

\pagegarde{Leçon 7 — Expression écrite : l'héritage}{Chinois}{Tle A}{Module 1 : Vie familiale et intégration sociale}{ZH07}{2 h}{Cette leçon d'expression écrite entraîne les élèves de Terminale A à produire un texte simple en chinois sur le thème de l'héritage familial. Elle couvre l'écriture correcte des caractères clés, un modèle de texte commenté avec connecteurs, ainsi que des exercices de thème et de version.
\vspace{4pt}
\begin{multicols}{2}\small
\begin{itemize}
\item À la fin de cette leçon, je sais écrire correctement les caractères du thème de l'héritage (遗, 产, 承, 嘱, 财, 法...) en respectant radicaux et ordre des traits
\item je sais distinguer les caractères proches (遗/遣, 嘱/属, 财/材) pour éviter les confusions à l'écrit
\item je sais lire et comprendre un texte simple en chinois sur une succession familiale
\item je sais utiliser les connecteurs utiles (首先, 然后, 因此, 但是, 所以) pour structurer un texte écrit
\item je sais rédiger un texte simple de 6 à 10 phrases sur le thème de l'héritage
\item je sais traduire du français vers le chinois (thème) des phrases simples sur la succession
\end{itemize}\end{multicols}}

\begin{sommaire}
\begin{multicols}{2}
\begin{enumerate}
\item Modèle de texte commenté
\item Lexique et glossaire du thème
\item Écriture correcte des caractères
\item Structure et connecteurs du texte
\item Thème : français → chinois
\item Version : chinois → français
\item Grille d'évaluation et entraînement Bac
\item Activité d'intégration
\item Méthodes et exercices résolus
\item Exercices
\item Corrigés des exercices
\item Fiche bilan
\end{enumerate}
\end{multicols}
\end{sommaire}

\begin{objectifs}
\begin{itemize}
\item À la fin de cette leçon, je sais écrire correctement les caractères du thème de l'héritage (遗, 产, 承, 嘱, 财, 法...) en respectant radicaux et ordre des traits
\item je sais distinguer les caractères proches (遗/遣, 嘱/属, 财/材) pour éviter les confusions à l'écrit
\item je sais lire et comprendre un texte simple en chinois sur une succession familiale
\item je sais utiliser les connecteurs utiles (首先, 然后, 因此, 但是, 所以) pour structurer un texte écrit
\item je sais rédiger un texte simple de 6 à 10 phrases sur le thème de l'héritage
\item je sais traduire du français vers le chinois (thème) des phrases simples sur la succession
\item je sais traduire du chinois vers le français (version) un court texte sur l'héritage
\item je sais auto-évaluer ma production écrite à l'aide d'une grille (caractères, grammaire, connecteurs, sens)
\end{itemize}
\end{objectifs}

\begin{prerequis}
\begin{itemize}
\item Vocabulaire de la famille et de la parenté (父亲, 母亲, 儿子, 女儿, 家) vu en Première
\item Structures de base : phrase sujet-verbe-complément, possessif avec 的
\item Verbes fréquents : 有, 给, 是, 把 (structure 把 vue en début d'année)
\item Nombres, classificateurs et expressions de temps (以后, 以前, 的时候)
\item Radicaux courants déjà étudiés (讠, 宀, 贝, 辶)
\end{itemize}
\end{prerequis}

\newpage

\coursec{Modèle de texte commenté}

\courssub{Situation d'accroche et problématique}

À Douala comme à Yaoundé, la disparition d'un chef de famille soulève toujours la même question : qui hérite de la maison, des terres, du magasin ? En Chine aussi, la succession (\zh{遗产}, \emph{yíchǎn}, « héritage ») est un sujet qui touche chaque famille, et la loi chinoise sur l'héritage a été modernisée en 2021 avec l'entrée en vigueur du nouveau Code civil. Au Cameroun, les traditions coutumières se mêlent au droit moderne, exactement comme dans de nombreux pays africains. Savoir rédiger un texte simple en chinois sur ce thème — raconter une succession, décrire un testament, traduire des phrases — est une compétence d'examen attendue en Terminale. Cette leçon vous donne les caractères, les connecteurs et la méthode pour y parvenir.

\textbf{Problématique :} comment raconter en chinois, de façon claire et correcte, le décès d'un proche, les biens qu'il laisse, les personnes qui héritent et la réaction de la famille ?

\courssub{Lecture du texte modèle}

Lisons d'abord le texte modèle, sans dictionnaire, pour en saisir le sens global. Soulignez mentalement les mots qui appartiennent au thème de l'héritage.

\begin{textebox}[我的爷爷去世了 — Wǒ de yéye qùshì le — « Mon grand-père est décédé »]
我的爷爷去年去世了。\\
\emph{Wǒ de yéye qùnián qùshì le.}\\
« Mon grand-père est décédé l'année dernière. »\\[4pt]
他留下了一所房子和一些钱。\\
\emph{Tā liúxià le yì suǒ fángzi hé yìxiē qián.}\\
« Il a laissé une maison et un peu d'argent. »\\[4pt]
按照法律，他的财产由我的爸爸和叔叔继承。\\
\emph{Ànzhào fǎlǜ, tā de cáichǎn yóu wǒ de bàba hé shūshu jìchéng.}\\
« Conformément à la loi, ses biens ont été hérités par mon père et mon oncle. »\\[4pt]
爸爸把房子留给了奶奶住。\\
\emph{Bàba bǎ fángzi liú gěi le nǎinai zhù.}\\
« Papa a laissé la maison à grand-mère pour qu'elle y habite. »\\[4pt]
我们一家人虽然很伤心，但是很团结。\\
\emph{Wǒmen yì jiā rén suīrán hěn shāngxīn, dànshì hěn tuánjié.}\\
« Notre famille, bien que très triste, est très unie. »
\end{textebox}

\courssub{Traduction mot à mot puis naturelle}

Pour bien comprendre la mécanique de chaque phrase, comparons la traduction littérale (mot à mot) et la traduction naturelle en français.

\begin{center}
\begin{tabularx}{\linewidth}{|X|X|X|}
\hline
\rowcolor{popblueL} Phrase chinoise & Mot à mot & Traduction naturelle \\
\hline
我的爷爷去年去世了。 & mon + grand-père + année dernière + décéder + \zh{了} & Mon grand-père est décédé l'année dernière. \\
\hline
他留下了一所房子和一些钱。 & il + laisser + \zh{了} + une (classif. \zh{所}) maison + et + un peu d'argent & Il a laissé une maison et un peu d'argent. \\
\hline
按照法律，他的财产由我的爸爸和叔叔继承。 & selon + loi, ses + biens + par (\zh{由}) + mon père et oncle + hériter & Conformément à la loi, ses biens sont hérités par mon père et mon oncle. \\
\hline
爸爸把房子留给了奶奶住。 & papa + \zh{把} + maison + laisser à + \zh{了} + grand-mère + habiter & Papa a laissé la maison à grand-mère pour qu'elle y vive. \\
\hline
我们一家人虽然很伤心，但是很团结。 & nous + toute la famille + bien que + très triste, mais + très uni & Notre famille, bien que très triste, reste très unie. \\
\hline
\end{tabularx}
\end{center}

\begin{attention}[Le passif chinois avec 由]
La phrase \zh{他的财产由我的爸爸和叔叔继承} est un \cle{passif} : le français dit « ses biens \emph{sont hérités par} mon père ». En chinois, la structure est : \textbf{chose + 由 + personne + verbe}. Ne traduisez jamais \zh{由} par « depuis » ou « de » ici : il introduit l'\emph{agent} de l'action. Comparez : \zh{这封信由老师写。} (\emph{Zhè fēng xìn yóu lǎoshī xiě.}) « Cette lettre est écrite par le professeur. » Remarque : ce passif avec \zh{由} insiste sur \emph{qui fait} l'action ; il est fréquent dans les textes administratifs et juridiques, donc parfaitement à sa place dans un récit de succession.
\end{attention}

\courssub{Repérage des éléments du récit}

Tout récit de succession répond à quatre questions. Entraînons-nous à y répondre à partir du texte modèle.

\begin{itemize}
\item \textbf{Qui décède ?} \zh{我的爷爷} (\emph{wǒ de yéye}) : mon grand-père (paternel). Le verbe est \zh{去世} (\emph{qùshì}), littéralement « quitter le monde », terme respectueux pour « mourir ».
\item \textbf{Que laisse-t-il ?} \zh{一所房子和一些钱} : une maison et un peu d'argent. Verbe-clé : \zh{留下} (\emph{liúxià}), « laisser (derrière soi) ».
\item \textbf{Qui hérite ?} \zh{我的爸爸和叔叔} : mon père et mon oncle (frère cadet du père). Verbe-clé : \zh{继承} (\emph{jìchéng}), « hériter » ; le cadre juridique est indiqué par \zh{按照法律} (\emph{ànzhào fǎlǜ}), « selon la loi ».
\item \textbf{Comment la famille réagit-elle ?} \zh{虽然很伤心，但是很团结} : bien que triste, elle reste unie. Le père laisse la maison à la grand-mère : geste de solidarité familiale.
\end{itemize}

\begin{definition}[Les cinq mots-clés du thème]
\zh{去世} (\emph{qùshì}, v.) : décéder (registre respectueux) ; \zh{留下} (\emph{liúxià}, v.) : laisser, léguer ; \zh{财产} (\emph{cáichǎn}, n.) : biens, patrimoine ; \zh{继承} (\emph{jìchéng}, v.) : hériter ; \zh{法律} (\emph{fǎlǜ}, n.) : la loi, le droit. Ces cinq mots doivent être \cle{soulignés} dans tout texte d'examen sur l'héritage : ce sont eux qui portent le sens du thème. Ajoutez-y \zh{遗产} (\emph{yíchǎn}, n.) : l'héritage, ce qui est transmis — à ne pas confondre avec \zh{财产}, qui désigne les biens en général, même du vivant du propriétaire.
\end{definition}

\courssub{Analyse de la structure du texte}

Le texte modèle suit le plan classique en \cle{trois parties}, qu'il faudra reproduire dans vos propres productions :

\begin{enumerate}
\item \textbf{Introduction — le décès} (phrase 1) : on annonce qui est mort et quand (\zh{去年}).
\item \textbf{Développement — les biens et les héritiers} (phrases 2 à 4) : ce qui est laissé, qui hérite selon la loi, et comment le partage se concrétise (\zh{爸爸把房子留给了奶奶住}).
\item \textbf{Conclusion — le sentiment familial} (phrase 5) : la réaction de la famille, avec le connecteur d'opposition \zh{虽然……但是……} « bien que... mais... ».
\end{enumerate}

\begin{popfigure}
\begin{tikzpicture}[
  node distance=0.55cm and 1.1cm,
  box/.style={draw=popblue, fill=popblueL, rounded corners=3pt, align=center, minimum width=3.1cm, minimum height=1.05cm, font=\small},
  arr/.style={-{Stealth[length=3mm]}, thick, popdark}
]
\node[box, align=center] (a){Décès\\去世 \emph{qùshì}};
\node[box, right=of a, align=center] (b){Biens laissés\\留下的财产};
\node[box, right=of b, align=center] (c){Héritiers\\继承人};
\node[box, below=of c, align=center] (d){Partage\\分配};
\node[box, left=of d, align=center] (e){Conclusion familiale\\团结 \emph{tuánjié}};
\draw[arr] (a) -- (b);
\draw[arr] (b) -- (c);
\draw[arr] (c) -- (d);
\draw[arr] (d) -- (e);
\end{tikzpicture}
\legende{Organigramme du récit de succession : du décès à la conclusion familiale.}
\end{popfigure}

\begin{methode}[Lire un texte d'examen en quatre étapes]
\begin{enumerate}
\item \textbf{Première lecture globale} : lisez tout le texte sans vous arrêter sur les caractères inconnus ; dégagez le sujet (ici : une succession familiale).
\item \textbf{Soulignez les mots du thème} : repérez les mots-clés comme \zh{去世}, \zh{留下}, \zh{财产}, \zh{继承}, \zh{法律}. Ils donnent le squelette du sens.
\item \textbf{Repérez les connecteurs} : \zh{按照} (selon), \zh{虽然……但是……} (bien que... mais...), \zh{和} (et). Ils indiquent les relations logiques entre les phrases.
\item \textbf{Identifiez la structure en trois parties} : introduction (décès) → développement (biens, héritiers) → conclusion (sentiments). Cette grille vous servira pour la compréhension \emph{et} pour votre propre rédaction.
\end{enumerate}
\end{methode}

\courssub{Exemples traduits et observation des structures}

\begin{exemplebox}[Deux phrases modèles à réutiliser]
\zh{他留下了一所房子。}\\
\emph{Tā liúxià le yì suǒ fángzi.}\\
« Il a laissé une maison. » — Structure : Sujet + \zh{留下} + \zh{了} + complément. Notez le classificateur \zh{所} pour les bâtiments.\\[6pt]
\zh{财产由儿子继承。}\\
\emph{Cáichǎn yóu érzi jìchéng.}\\
« Les biens sont hérités par le fils. » — Structure passive : chose + \zh{由} + agent + verbe.
\end{exemplebox}

\begin{propriete}[La construction avec 把 dans les phrases de transfert]
Pour dire « A laisse/donne l'objet X à B », le chinois utilise souvent la construction : \textbf{A + 把 + objet + verbe + 给 + B}. Exemple du texte : \zh{爸爸把房子留给了奶奶住} « Papa a laissé la maison à grand-mère ». Le français place l'objet après le verbe ; le chinois l'avance avec \zh{把} : ne dites jamais \zh{爸爸留给房子奶奶}, c'est un ordre des mots impossible. La particule \zh{了} se place après le verbe (ou après \zh{给}) : \zh{把房子留给了奶奶}.
\end{propriete}

\begin{attention}[Le ton de 一 devant un quatrième ton]
Dans \zh{一所房子}, le caractère \zh{一} se prononce \emph{yì} (deuxième ton) parce qu'il précède un quatrième ton (\zh{所}, \emph{suǒ}, est au troisième ton : on dit alors \emph{yí} ? Non — vérifions : \zh{所} est au troisième ton, donc \zh{一} reste au quatrième ton de base devant un ton non-quatrième : \emph{yì suǒ}). Règle : \zh{一} se prononce \emph{yí} devant un quatrième ton (\zh{一份遗嘱}, \emph{yí fèn yízhǔ}), et \emph{yì} devant les premier, deuxième et troisième tons (\zh{一所房子}, \emph{yì suǒ fángzi} ; \zh{一些钱}, \emph{yìxiē qián}). À l'écrit, on note le ton réellement prononcé.
\end{attention}

\begin{popfigure}
\begin{tikzpicture}[mindmap, grow cyclic,
  every node/.style={concept, align=center, font=\small},
  root concept/.append style={concept color=popblue, text=white, minimum size=2.6cm},
  level 1/.append style={level distance=3.4cm, sibling angle=90},
  level 1 concept/.append style={concept color=poporangeL, text=popdark, minimum size=2.2cm}
]
\node[root concept, align=center]{遗产\\ \emph{yíchǎn}\\ héritage}
  child { node[align=center]{Personnes\\ 继承人 héritiers\\ 爸爸 叔叔 儿子} }
  child { node[align=center]{Biens\\ 房子 maison\\ 钱 argent\\ 财产 patrimoine} }
  child { node[align=center]{Loi\\ 法律 loi\\ 按照 selon\\ 继承法 successions} }
  child { node[align=center]{Sentiments\\ 伤心 triste\\ 团结 uni} };
\end{tikzpicture}
\legende{Carte mentale lexicale centrée sur 遗产 : les quatre familles de mots du thème.}
\end{popfigure}

\courssub{Texte d'application : lecture guidée}

\begin{textebox}[叔叔的遗嘱 — Shūshu de yízhǔ — « Le testament de mon oncle »]
我的叔叔是杜阿拉的一个商人。今年三月，他生病去世了。\\
\emph{Wǒ de shūshu shì Dù'ālā de yí ge shāngrén. Jīnnián sānyuè, tā shēngbìng qùshì le.}\\
« Mon oncle était un commerçant de Douala. En mars de cette année, il est mort de maladie. »\\[4pt]
他生前写了一份遗嘱。按照遗嘱，他的商店由他的儿子继承，\\
\emph{Tā shēngqián xiě le yí fèn yízhǔ. Ànzhào yízhǔ, tā de shāngdiàn yóu tā de érzi jìchéng,}\\
« De son vivant, il avait écrit un testament. Selon le testament, son magasin est hérité par son fils, »\\[4pt]
房子留给他的妻子，一部分钱捐给了村里的学校。\\
\emph{fángzi liú gěi tā de qīzi, yí bùfen qián juān gěi le cūn lǐ de xuéxiào.}\\
« la maison est laissée à sa femme, et une partie de l'argent a été donnée à l'école du village. »\\[4pt]
开始的时候，家里有些人不高兴，因为他们也想继承商店。\\
\emph{Kāishǐ de shíhou, jiā lǐ yǒuxiē rén bù gāoxìng, yīnwèi tāmen yě xiǎng jìchéng shāngdiàn.}\\
« Au début, certaines personnes de la famille n'étaient pas contentes, car elles voulaient aussi hériter du magasin. »\\[4pt]
后来，大家一起读了遗嘱，明白了叔叔的想法。\\
\emph{Hòulái, dàjiā yìqǐ dú le yízhǔ, míngbai le shūshu de xiǎngfǎ.}\\
« Ensuite, tous ont lu le testament ensemble et ont compris la pensée de l'oncle. »\\[4pt]
现在，我们一家人又团结又和睦。\\
\emph{Xiànzài, wǒmen yì jiā rén yòu tuánjié yòu hémù.}\\
« Maintenant, notre famille est à la fois unie et harmonieuse. »
\end{textebox}

\begin{center}
\begin{tabularx}{\linewidth}{|X|X|X|X|}
\hline
\rowcolor{popblueL} Caractères & Pinyin & Nature & Traduction \\
\hline
遗产 & yíchǎn & n. & héritage, patrimoine transmis \\
\hline
遗嘱 & yízhǔ & n. & testament \\
\hline
生前 & shēngqián & loc. temp. & de son vivant \\
\hline
商人 & shāngrén & n. & commerçant \\
\hline
商店 & shāngdiàn & n. & magasin, boutique \\
\hline
捐 & juān & v. & donner (faire un don) \\
\hline
明白 & míngbai & v. & comprendre \\
\hline
想法 & xiǎngfǎ & n. & idée, pensée \\
\hline
和睦 & hémù & adj. & harmonieux \\
\hline
\end{tabularx}
\end{center}

\textbf{Questions de compréhension.} 1) Qui est décédé et quand ? 2) Quels sont les trois biens mentionnés et qui les reçoit ? 3) Pourquoi certaines personnes n'étaient-elles pas contentes ? 4) Quel connecteur marque l'opposition dans la dernière phrase ?

\begin{exoresolu}[Repérage dans le texte « Le testament de mon oncle »]
Énoncé. Relevez dans le texte : (a) le verbe qui signifie « décéder » ; (b) les deux structures passives avec \zh{由} ou \zh{给} ; (c) le connecteur de cause ; (d) la conclusion familiale. Traduisez chaque élément.
\tcblower
Solution. (a) \zh{去世} (\emph{qùshì}) : « décéder », dans \zh{他生病去世了} « il est mort de maladie ». (b) \zh{他的商店由他的儿子继承} « son magasin est hérité par son fils » (passif avec \zh{由}) et \zh{一部分钱捐给了村里的学校} « une partie de l'argent a été donnée à l'école du village » (\zh{捐给} + destinataire). (c) \zh{因为} (\emph{yīnwèi}) : « parce que », dans \zh{因为他们也想继承商店} « parce qu'ils voulaient aussi hériter du magasin ». (d) \zh{我们一家人又团结又和睦} « notre famille est à la fois unie et harmonieuse » : conclusion positive, avec la structure \zh{又……又……} « à la fois... et... », comme dans le texte modèle.
\end{exoresolu}

\begin{exercice}[À vous de jouer : mini-production guidée]
En vous inspirant du texte modèle, écrivez trois phrases en chinois : (1) annoncez le décès de votre grand-mère l'année dernière ; (2) dites qu'elle a laissé une maison, héritée par votre mère (utilisez \zh{由}) ; (3) concluez sur l'union de la famille avec \zh{虽然……但是……}. Vous utiliserez obligatoirement les mots \zh{去世}, \zh{留下}, \zh{继承}.
\end{exercice}

\begin{corrige}[Mini-production guidée]
Modèle de réponse : \zh{我的奶奶去年去世了。她留下了一所房子，房子由我的妈妈继承。我们一家人虽然很伤心，但是很团结。} (\emph{Wǒ de nǎinai qùnián qùshì le. Tā liúxià le yì suǒ fángzi, fángzi yóu wǒ de māma jìchéng. Wǒmen yì jiā rén suīrán hěn shāngxīn, dànshì hěn tuánjié.}) « Ma grand-mère est décédée l'année dernière. Elle a laissé une maison, héritée par ma mère. Notre famille, bien que très triste, est très unie. » Barème indicatif : 1 point par mot-clé correctement employé, 1 point pour le passif avec \zh{由}, 1 point pour la paire \zh{虽然……但是……}.
\end{corrige}

\begin{savaistu}[La loi chinoise sur les successions]
En Chine, la loi sur les successions (\zh{继承法}, \emph{jìchéngfǎ}), intégrée en 2021 au nouveau Code civil, prévoit que les enfants, le conjoint et les parents du défunt sont héritiers au premier rang, hommes et femmes à \cle{égalité}. C'est un point de comparaison intéressant avec le Cameroun, où le droit moderne coexiste avec les coutumes : dans beaucoup de familles camerounaises, la succession est réglée par la réunion de famille, parfois au détriment des filles ou des veuves. En classe de chinois, ce sujet permet un débat riche : \zh{中国和喀麦隆的继承习惯一样吗？} (\emph{Zhōngguó hé Kāmàilóng de jìchéng xíguàn yíyàng ma?}) « Les coutumes successorales chinoises et camerounaises sont-elles les mêmes ? »
\end{savaistu}

\begin{attention}[Pièges fréquents des francophones]
\begin{itemize}
\item \textbf{Oublier le classificateur} : on dit \zh{一所房子} (une maison), jamais \zh{一房子}. Pour un testament, c'est \zh{一份遗嘱} ; pour une lettre, \zh{一封信}.
\item \textbf{Confondre} \zh{去世} et \zh{死} (\emph{sǐ}) : \zh{死} est direct, voire brutal ; pour un proche, utilisez toujours le respectueux \zh{去世}.
\item \textbf{Mal placer} \zh{了} : il suit le verbe (\zh{留下了}, \zh{去世了}), pas le complément ; on ne dit pas \zh{他留下一所房子了} dans ce récit.
\item \textbf{Traduire « bien que » par un seul connecteur} : en chinois, \zh{虽然} et \zh{但是} fonctionnent en paire dans la même phrase, contrairement au français qui n'emploie qu'un seul connecteur. N'omettez pas \zh{但是} dans la seconde partie.
\item \textbf{Confondre} \zh{遗产} et \zh{财产} : \zh{财产} désigne les biens en général ; \zh{遗产} désigne précisément ce qui est transmis après un décès.
\end{itemize}
\end{attention}

\begin{aretenir}[À retenir de cette section]
Un récit de succession en chinois se construit en trois temps : \textbf{décès} (\zh{去世}) → \textbf{biens et héritiers} (\zh{留下}, \zh{财产}, \zh{由……继承}) → \textbf{conclusion familiale} (\zh{虽然……但是……}, \zh{团结}). Maîtriser les cinq mots-clés, le passif avec \zh{由} et la construction avec \zh{把} suffit à produire un texte simple, correct et bien structuré — exactement ce qu'attend l'examinateur.
\end{aretenir}

\coursec{Lexique et glossaire du thème}

Dans la section précédente, nous avons lu et commenté un texte modèle sur l'héritage : la famille de M. Mbarga, à Yaoundé, se réunit après le décès du grand-père pour partager ses biens dans le respect de son testament. Ce texte reposait sur un petit nombre de mots-clés — \zh{遗产}, \zh{继承}, \zh{遗嘱}, \zh{分配}… — que nous allons maintenant étudier un par un, en profondeur. Un texte d'expression écrite ne peut être réussi que si son vocabulaire est \cle{exact}, \cle{précis} et \cle{correctement écrit} : c'est l'objectif de cette section.

\courssub{Le glossaire du thème : observation d'abord}

Avant tout tableau, lisons ces deux phrases tirées de situations réelles de la vie camerounaise :

\zh{爷爷去世后，留下了一大笔遗产。}

\emph{Yéye qùshì hòu, liúxià le yí dà bǐ yíchǎn.}

« Après le décès de grand-père, il a laissé un gros héritage. »

\zh{他们按照法律分配遗产，一家人虽然很伤心，但是很团结。}

\emph{Tāmen ànzhào fǎlǜ fēnpèi yíchǎn, yì jiā rén suīrán hěn shāngxīn, dànshì hěn tuánjié.}

« Ils partagent l'héritage selon la loi ; toute la famille, bien que triste, reste unie. »

Observons : le mot \zh{遗产} apparaît deux fois — c'est le \cle{mot pivot} du thème. Autour de lui gravitent un verbe d'action (\zh{留下}, « laisser »), un verbe de réception (\zh{继承}, « hériter »), un cadre juridique (\zh{遗嘱}, \zh{法律}) et des sentiments (\zh{伤心}, \zh{团结}). Tout bon texte sur l'héritage organise exactement ces quatre pôles : le bien, les actions, la loi, les émotions.

\begin{definition}[遗产 yíchǎn — le mot central du thème]
\zh{遗产} (yíchǎn), nom, désigne l'\cle{ensemble des biens laissés par une personne décédée} : maison, terrain, argent, objets. Le caractère \zh{遗} signifie « laisser, léguer » et \zh{产} signifie « bien, propriété » : littéralement, « les biens laissés ». On l'emploie aussi au sens figuré : \zh{文化遗产} (wénhuà yíchǎn), « le patrimoine culturel ».
\end{definition}

\courssub{Tableau glossaire du thème}

\begin{center}
\begin{tabularx}{\linewidth}{|X|X|X|X|}
\hline
\rowcolor{popblueL} Caractères & Pinyin & Nature & Traduction \\
\hline
遗产 & yíchǎn & n. & héritage \\
\hline
继承 & jìchéng & v. & hériter \\
\hline
财产 & cáichǎn & n. & biens, patrimoine \\
\hline
遗嘱 & yízhǔ & n. & testament \\
\hline
去世 & qùshì & v. & décéder \\
\hline
留下 & liúxià & v. & laisser (en héritage) \\
\hline
法律 & fǎlǜ & n. & loi \\
\hline
分配 & fēnpèi & v. & répartir, partager \\
\hline
团结 & tuánjié & adj./v. & uni ; s'unir \\
\hline
伤心 & shāngxīn & adj. & triste, affligé \\
\hline
\end{tabularx}
\end{center}

\begin{attention}[Caractères proches : ne pas confondre]
Trois paires de caractères du glossaire se ressemblent dangereusement :
\begin{itemize}
\item \zh{遗} (yí, laisser) contient la clé du mouvement \zh{辶} ; ne pas l'écrire comme \zh{遣} (qiǎn, envoyer).
\item \zh{产} (chǎn, bien) se distingue de \zh{严} (yán, sévère) par le trait supérieur.
\item \zh{财} (cái, richesse) a la clé \zh{贝} à gauche ; \zh{材} (cái, matériau) a la clé du bois \zh{木}. Même pinyin, sens différents !
\end{itemize}
\end{attention}

\courssub{Carte mentale du champ lexical}

Pour mémoriser durablement, organisons le vocabulaire en quatre familles autour du mot pivot \zh{遗产}.

\begin{popfigure}
\begin{tikzpicture}[mindmap, grow cyclic, every node/.style={concept, align=center, font=\small},
root concept/.append style={concept color=popblue, font=\small\bfseries},
level 1/.append style={level distance=3.2cm, sibling angle=90},
level 1 concept/.append style={concept color=poporangeL}]
\node[root concept, align=center]{遗产\\ yíchǎn\\ héritage}
child { node[align=center]{Actions\\ 留下 liúxià\\ 继承 jìchéng\\ 分配 fēnpèi} }
child { node[align=center]{Biens\\ 财产 cáichǎn\\ 房子 fángzi\\ 钱 qián} }
child { node[align=center]{Droit\\ 遗嘱 yízhǔ\\ 法律 fǎlǜ\\ 去世 qùshì} }
child { node[align=center]{Sentiments\\ 伤心 shāngxīn\\ 团结 tuánjié\\ 一家人 yì jiā rén} };
\end{tikzpicture}
\legende{Carte mentale du champ lexical de l'héritage : quatre familles autour du mot pivot 遗产.}
\end{popfigure}

\courssub{Les collocations : les mots ne vivent pas seuls}

En chinois comme en français, un mot s'emploie avec des partenaires habituels. Apprendre \zh{遗产} seul ne suffit pas : il faut connaître ses \cle{collocations}, c'est-à-dire les verbes qui l'accompagnent naturellement.

\begin{center}
\begin{tabularx}{\linewidth}{|X|X|X|}
\hline
\rowcolor{popblueL} Collocation & Pinyin & Traduction \\
\hline
留下遗产 & liúxià yíchǎn & laisser un héritage \\
\hline
继承财产 & jìchéng cáichǎn & hériter des biens \\
\hline
立遗嘱 & lì yízhǔ & rédiger un testament \\
\hline
按照法律 & ànzhào fǎlǜ & selon la loi, conformément à la loi \\
\hline
分配遗产 & fēnpèi yíchǎn & partager l'héritage \\
\hline
\end{tabularx}
\end{center}

\begin{propriete}[Structure des collocations verbe + objet]
En chinois, le verbe précède toujours son objet : \cle{Verbe + Objet}. Contrairement au français qui utilise souvent une préposition (« hériter \emph{de} la maison », « rédiger \emph{un} testament »), le chinois accroche directement l'objet au verbe : \zh{继承财产} (litt. « hériter biens »), \zh{立遗嘱} (litt. « dresser testament »). Le verbe \zh{立} (lì, « dresser, établir ») est le verbe consacré pour le testament : on dit \zh{立遗嘱} et non \zh{写遗嘱} dans un registre soigné, même si \zh{写} (écrire) se comprend.
\end{propriete}

\begin{exemplebox}[Les collocations en contexte]
\zh{爷爷留下了一大笔遗产。}\\
\emph{Yéye liúxià le yí dà bǐ yíchǎn.} — « Grand-père a laissé un gros héritage. » (\zh{大笔} = « une grosse somme de », classificateur \zh{笔} pour les sommes d'argent.)

\zh{他们按照法律分配遗产。}\\
\emph{Tāmen ànzhào fǎlǜ fēnpèi yíchǎn.} — « Ils partagent l'héritage selon la loi. »

\zh{父亲去世前立了遗嘱。}\\
\emph{Fùqin qùshì qián lì le yízhǔ.} — « Avant de décéder, père avait rédigé un testament. »

\zh{大儿子继承了房子和土地。}\\
\emph{Dà érzi jìchéng le fángzi hé tǔdì.} — « Le fils aîné a hérité de la maison et du terrain. »
\end{exemplebox}

\courssub{Le registre soutenu : 去世 ou 死 ?}

\begin{propriete}[Nuances de registre : dire « mourir » avec respect]
Le chinois distingue soigneusement les registres de langue, surtout pour les sujets sensibles :
\begin{itemize}
\item \zh{死} (sǐ) : « mourir », registre neutre voire brutal ; à éviter pour parler d'un proche dans un texte écrit.
\item \zh{去世} (qùshì) : « décéder », litt. « quitter le monde » ; registre \cle{soutenu et respectueux}, à privilégier dans toute rédaction.
\item \zh{逝世} (shìshì) : registre encore plus solennel, réservé aux personnalités.
\end{itemize}
Dans une copie de Bac, écrivez toujours \zh{爷爷去世了} (« grand-père est décédé ») plutôt que \zh{爷爷死了}.
\end{propriete}

\begin{attention}[Piège de traduction : « hériter » n'est pas 收到]
Les francophones traduisent souvent « hériter » par \zh{收到} (shōudào, « recevoir »). C'est une erreur : \zh{收到} signifie recevoir un objet, une lettre, un colis — sans aucune idée de succession. Le verbe exact est \zh{继承} (jìchéng).\\
« Il hérite de la maison » = \zh{他继承了房子。} (\emph{Tā jìchéng le fángzi.}) et non \zh{他收到了房子} (qui voudrait dire « il a reçu la maison », comme un paquet !). De même, « recevoir un héritage » au sens juridique se dit \zh{继承遗产}, jamais \zh{收到遗产}.
\end{attention}

\courssub{Expressions figées à réutiliser dans vos rédactions}

\begin{definition}[Deux expressions toutes faites pour le texte d'héritage]
\begin{itemize}
\item \zh{一家人} (yì jiā rén) : « toute la famille », « la famille (comme un tout) ». S'emploie comme sujet : \zh{一家人都很伤心。} « Toute la famille est très triste. »
\item \zh{虽然很伤心，但是很团结} (suīrán hěn shāngxīn, dànshì hěn tuánjié) : « bien que tristes, mais unis ». Structure \zh{虽然……但是……} (« bien que… mais… ») : en chinois, contrairement au français, les deux termes de la concession se disent. Cette phrase est une excellente conclusion de texte sur l'héritage.
\end{itemize}
\end{definition}

\begin{savaistu}[Pourquoi le caractère 财 contient-il un coquillage ?]
Le caractère \zh{财} (cái, « richesse, bien ») contient à gauche le radical \zh{贝} (bèi, « coquillage »). Dans la Chine ancienne, bien avant l'invention des pièces de monnaie, les coquillages servaient de monnaie d'échange. C'est pourquoi presque tous les caractères liés à l'argent contiennent \zh{贝} : \zh{贵} (guì, cher), \zh{贫} (pín, pauvre), \zh{购} (gòu, acheter), \zh{货} (huò, marchandise). Retenir la clé \zh{贝}, c'est retenir toute une famille de mots de l'argent — très utile pour le thème de l'héritage !
\end{savaistu}

\courssub{Texte de réinvestissement lexical}

\begin{textebox}[Le partage de la famille Essomba]
埃松巴爷爷上个月去世了。他留下了一大笔遗产：一所房子、一块地和一些钱。去世以前，他立了遗嘱。\\
\emph{Àisōngbā yéye shàng ge yuè qùshì le. Tā liúxià le yí dà bǐ yíchǎn : yì suǒ fángzi, yí kuài dì hé yìxiē qián. Qùshì yǐqián, tā lì le yízhǔ.}\\
« Le grand-père Essomba est décédé le mois dernier. Il a laissé un gros héritage : une maison, un terrain et un peu d'argent. Avant de décéder, il avait rédigé un testament. »\\
一家人按照法律分配遗产。大儿子继承了房子，女儿继承了地。\\
\emph{Yì jiā rén ànzhào fǎlǜ fēnpèi yíchǎn. Dà érzi jìchéng le fángzi, nǚ'ér jìchéng le dì.}\\
« Toute la famille partage l'héritage selon la loi. Le fils aîné hérite de la maison, la fille hérite du terrain. »\\
虽然很伤心，但是他们很团结。\\
\emph{Suīrán hěn shāngxīn, dànshì tāmen hěn tuánjié.}\\
« Bien que très tristes, ils restent très unis. »
\end{textebox}

\textbf{Glossaire du texte.} \zh{上个月} (shàng ge yuè) : le mois dernier ; \zh{一所房子} (yì suǒ fángzi) : une maison (classificateur \zh{所} pour les bâtiments) ; \zh{一块地} (yí kuài dì) : un terrain ; \zh{以前} (yǐqián) : avant ; \zh{女儿} (nǚ'ér) : fille.

\textbf{Questions de compréhension.}
\begin{enumerate}
\item Quand le grand-père Essomba est-il décédé ?
\item Quels biens a-t-il laissés ?
\item Qui hérite de la maison ? Qui hérite du terrain ?
\item Relevez dans le texte trois collocations étudiées dans cette section.
\end{enumerate}

\courssub{Exercice résolu et entraînement}

\begin{exoresolu}[Complétez avec le mot du glossaire qui convient]
Complétez : 1. 奶奶去世了，留下了很多……。 2. 孩子们……了父母的财产。 3. 他们……法律分配遗产。 4. 虽然很……，但是一家人很……。
\tcblower
\textbf{Solution détaillée.}
\begin{enumerate}
\item \zh{奶奶去世了，留下了很多遗产。} (\emph{Nǎinai qùshì le, liúxià le hěn duō yíchǎn.}) « Grand-mère est décédée et a laissé beaucoup de biens. » Le verbe \zh{留下} appelle l'objet \zh{遗产}.
\item \zh{孩子们继承了父母的财产。} (\emph{Háizimen jìchéng le fùmǔ de cáichǎn.}) « Les enfants ont hérité des biens de leurs parents. » Le verbe exact pour « hériter » est \zh{继承}, pas \zh{收到}.
\item \zh{他们按照法律分配遗产。} (\emph{Tāmen ànzhào fǎlǜ fēnpèi yíchǎn.}) « Ils partagent l'héritage selon la loi. » Collocation figée \zh{按照法律}.
\item \zh{虽然很伤心，但是一家人很团结。} (\emph{Suīrán hěn shāngxīn, dànshì yì jiā rén hěn tuánjié.}) « Bien que très tristes, toute la famille reste unie. » Structure de concession \zh{虽然……但是……}.
\end{enumerate}
\end{exoresolu}

\begin{exercice}[À vous de jouer]
\begin{enumerate}
\item Traduisez en chinois : « Avant de décéder, mon grand-père a rédigé un testament. »
\item Traduisez en chinois : « Les enfants partagent l'héritage selon la loi. »
\item Trouvez l'erreur et corrigez : \zh{他收到了父亲的房子。} (sens voulu : « Il a hérité de la maison de son père. »)
\item Classez ces mots dans les quatre familles de la carte mentale : \zh{遗嘱}, \zh{伤心}, \zh{财产}, \zh{分配}.
\end{enumerate}
\end{exercice}

\begin{corrige}[Exercice « À vous de jouer »]
\begin{enumerate}
\item \zh{爷爷去世以前立了遗嘱。} (\emph{Yéye qùshì yǐqián lì le yízhǔ.}) — Attention à la place de \zh{以前} : après le groupe verbal.
\item \zh{孩子们按照法律分配遗产。} (\emph{Háizimen ànzhào fǎlǜ fēnpèi yíchǎn.})
\item Il faut remplacer \zh{收到} par \zh{继承} : \zh{他继承了父亲的房子。} (\emph{Tā jìchéng le fùqin de fángzi.})
\item \zh{遗嘱} : famille du droit ; \zh{伤心} : famille des sentiments ; \zh{财产} : famille des biens ; \zh{分配} : famille des actions.
\end{enumerate}
\end{corrige}

\begin{aretenir}[L'essentiel du lexique de l'héritage]
\begin{itemize}
\item Le mot pivot est \zh{遗产} (yíchǎn) ; il s'entoure de quatre familles : actions (\zh{留下}, \zh{继承}, \zh{分配}), biens (\zh{财产}), droit (\zh{遗嘱}, \zh{法律}, \zh{去世}), sentiments (\zh{伤心}, \zh{团结}).
\item Apprenez les mots avec leurs collocations : \zh{留下遗产}, \zh{继承财产}, \zh{立遗嘱}, \zh{按照法律}, \zh{分配遗产}.
\item Registre soutenu : \zh{去世} et non \zh{死} ; « hériter » se dit \zh{继承} et jamais \zh{收到}.
\item La clé \zh{贝} (coquillage = monnaie ancienne) rassemble les caractères de la richesse : \zh{财}, \zh{贵}, \zh{贫}.
\end{itemize}
\end{aretenir}

Ce lexique solidement acquis, nous pouvons passer à la section suivante : l'écriture correcte des caractères du thème, trait par trait.

\coursec{Écriture correcte des caractères}

Dans la section précédente, nous avons constitué le lexique du thème de l'héritage : \zh{遗产} (yíchǎn, l'héritage), \zh{遗嘱} (yízhǔ, le testament), \zh{继承} (jìchéng, hériter), \zh{财产} (cáichǎn, les biens), \zh{法律} (fǎlǜ, la loi). Mais connaître un mot ne suffit pas au baccalauréat : il faut aussi savoir l'\emph{écrire} correctement, trait par trait. Un caractère mal tracé — clé oubliée, trait manquant, composant confondu — peut changer le sens du mot ou le rendre illisible. Cette section vous donne les règles d'ordre des traits et les astuces de mémorisation pour les \textbf{sept caractères essentiels} du thème (les six du vocabulaire de base plus \zh{继}).

\courssub{Observer avant d'écrire : la logique de construction des caractères}

\begin{exemplebox}[Observer les caractères du thème]
Regardez attentivement ces mots avant de lire les explications :\\
\zh{遗产} yíchǎn « l'héritage » \quad \zh{遗嘱} yízhǔ « le testament » \quad \zh{继承} jìchéng « hériter »\\
\zh{财产} cáichǎn « les biens » \quad \zh{法律} fǎlǜ « la loi »\\[0.3em]
Que remarquez-vous ? Chaque caractère est \cle{composé} de parties plus petites : \zh{遗} contient \zh{贵} (guì, précieux) et la clé du mouvement \zh{辶} ; \zh{财} contient \zh{贝} (bèi, le coquillage, ancienne monnaie) ; \zh{法} contient la clé de l'eau \zh{氵} ; \zh{继} contient la clé de la soie \zh{纟}. Un caractère chinois n'est presque jamais un dessin arbitraire : c'est un assemblage logique de composants (clé sémantique + composant phonétique ou sémantique).
\end{exemplebox}

\begin{popfigure}
\begin{tikzpicture}[font=\large, node distance=1.2cm and 1.5cm]
% 遗 = 贵 + 辶
\node[draw=popblue, fill=popblueL, rounded corners, minimum width=1.6cm, minimum height=1.6cm] (yi) {\zh{遗}};
\node[draw=poporange, fill=poporangeL, rounded corners, minimum width=1.1cm, minimum height=1.1cm] (gui) [above left=of yi] {\zh{贵}};
\node[draw=popgreen, fill=popgreenL, rounded corners, minimum width=1.1cm, minimum height=1.1cm] (chuo) [below left=of yi] {\zh{辶}};
\draw[-{Stealth}, thick, poporange] (gui) -- (yi);
\draw[-{Stealth}, thick, popgreen] (chuo) -- (yi);
\node[popmuted, align=center, anchor=north] at (gui.south) {\small guì : précieux};
\node[popmuted, align=center, anchor=north] at (chuo.south) {\small clé du mouvement};
% 财 = 贝 + 才
\node[draw=popblue, fill=popblueL, rounded corners, minimum width=1.6cm, minimum height=1.6cm] (cai) [right=3.5cm of yi] {\zh{财}};
\node[draw=poporange, fill=poporangeL, rounded corners, minimum width=1.1cm, minimum height=1.1cm] (bei) [above left=of cai] {\zh{贝}};
\node[draw=poppurple, fill=poppurpleL, rounded corners, minimum width=1.1cm, minimum height=1.1cm] (cai2) [above right=of cai] {\zh{才}};
\draw[-{Stealth}, thick, poporange] (bei) -- (cai);
\draw[-{Stealth}, thick, poppurple] (cai2) -- (cai);
\node[popmuted, anchor=north] at (bei.south) {\small bèi : coquillage};
\node[popmuted, anchor=north] at (cai2.south) {\small cái : talent};
% 继 = 纟 + 米/戊 (partie droite)
\node[draw=popblue, fill=popblueL, rounded corners, minimum width=1.6cm, minimum height=1.6cm] (ji) [right=3.5cm of cai] {\zh{继}};
\node[draw=poporange, fill=poporangeL, rounded corners, minimum width=1.1cm, minimum height=1.1cm] (si) [above left=of ji] {\zh{纟}};
\node[draw=poppurple, fill=poppurpleL, rounded corners, minimum width=1.1cm, minimum height=1.1cm] (ji_right) [above right=of ji] {\zh{米}/\zh{戊}};
\draw[-{Stealth}, thick, poporange] (si) -- (ji);
\draw[-{Stealth}, thick, poppurple] (ji_right) -- (ji);
\node[popmuted, anchor=north] at (si.south) {\small sī : soie};
\node[popmuted, anchor=north] at (ji_right.south) {\small partie droite (phonétique)};
\end{tikzpicture}
\legende{Décomposition des caractères 遗, 财 et 继 en composants connus.}
\end{popfigure}

\courssub{Les règles générales d'ordre des traits}

\begin{propriete}[Règles fondamentales d'écriture]
L'écriture des caractères chinois obéit à des règles stables, à appliquer dans l'ordre de priorité suivant :
\begin{enumerate}
\item \cle{De haut en bas} : \zh{产} se commence par le point du haut ; \zh{承} commence par le trait horizontal supérieur.
\item \cle{De gauche à droite} : \zh{财} s'écrit d'abord \zh{贝} (à gauche), puis \zh{才} (à droite) ; \zh{继} s'écrit \zh{纟} puis la partie droite.
\item \cle{Horizontal avant vertical} quand ils se croisent (ex: \zh{十}, \zh{干}, le croisement dans \zh{承}).
\item \cle{Extérieur avant l'intérieur} pour les caractères encadrés (ex: \zh{国}, \zh{园}), mais on \cle{referme en dernier} (le trait de fermeture horizontal du bas).
\item \cle{Centre avant les côtés} pour les caractères à symétrie verticale marquée : dans \zh{承}, le crochet vertical central (trait 5) précède les traits latéraux du bas (traits 6 et 7).
\item \cle{Règle spéciale de la clé 辶} : pour les caractères avec la clé du mouvement \zh{辶} comme \zh{遗}, \zh{遣}, \zh{这}, \zh{进}, on écrit \textbf{toujours la partie intérieure (le corps du caractère) d'abord} et la clé \zh{辶} \textbf{en dernier} (3 traits : point, virgule brisée \zh{㇆}, puis la longue vague finale \zh{㇏}).
\end{enumerate}
\end{propriete}

\begin{attention}[La clé 辶 en dernier, sans exception]
Beaucoup d'élèves francophones, habitués à écrire de gauche à droite, tracent d'abord la clé \zh{辶} qui « enveloppe » le caractère par la gauche et le bas. C'est une erreur majeure. Dans \zh{遗}, on écrit d'abord \zh{贵} en entier (9 traits au comptage pédagogique), \emph{puis seulement} la clé \zh{辶} (3 traits). Même règle pour \zh{这} (zhè, ceci), \zh{进} (jìn, entrer), \zh{遣} (qiǎn, envoyer), \zh{近} (jìn, près).
\end{attention}

\courssub{Étude détaillée des sept caractères du thème}

\textbf{遗 (yí) — laisser, léguer ; 12 traits (comptage pédagogique : 贵 9 + 辶 3).}\par
Composants : \zh{贵} (guì, précieux) + clé \zh{辶} (chuò, mouvement). 
Ordre : d'abord \zh{贵} de haut en bas (le haut \zh{覀}, puis \zh{冖}, puis \zh{贝} et les deux points), ensuite la clé \zh{辶} : 1. point, 2. virgule brisée (horizontal + crochet vers le bas), 3. longue vague finale qui « porte » tout le caractère. 
Mots : \zh{遗产} (yíchǎn), \zh{遗嘱} (yízhǔ), \zh{遗忘} (yíwàng, oublier).

\textbf{产 (chǎn) — produire, naître ; 6 traits.}\par
Structure : haut \zh{立} (simplifié en traits) + bas. 
Ordre : 1. point (sommet), 2. horizontal court, 3. point gauche, 4. point droit, 5. horizontal long (traversant), 6. grand trait tombant \zh{丿} (piě) final qui équilibre le caractère. 
Mots : \zh{遗产}, \zh{财产}, \zh{产生} (chǎnshēng, naître/apparaître).

\textbf{继 (jì) — continuer, succéder, hériter ; 10 traits.}\par
Structure : gauche-droite. Radical \zh{纟} (sī, soie, 3 traits) à gauche ; partie droite (7 traits) en haut à droite. 
Ordre : 1--3. Clé \zh{纟} (trois traits diagonaux montants : \zh{㇀}, \zh{㇀}, \zh{㇀} avec crochet final). 4--10. Partie droite de haut en bas : point, horizontal, trait tombant, point, horizontal, vertical, point final. 
Astuce : la soie (\zh{纟}) qui se transmet de génération en génération. 
Mots : \zh{继承} (jìchéng), \zh{继续} (jìxù, continuer).

\textbf{承 (chéng) — supporter, recevoir, hériter ; 8 traits.}\par
Structure : haut \zh{受} (recevoir) + bas \zh{寸} (pouce/règle), mais écrit en un bloc unique à symétrie centrale apparente. 
Ordre exact (respectant « haut→bas » et « centre→côtés » pour le bas) :
\begin{enumerate}
\item Horizontal haut (sommet du \zh{受})
\item Vertical gauche (côté gauche du \zh{受})
\item Horizontal milieu (barre du \zh{受})
\item Horizontal bas (fermeture du haut \zh{受})
\item \textbf{Crochet vertical central} (trait 5, axe de symétrie, commence au milieu du trait 4)
\item Horizontal gauche (bras gauche du \zh{寸})
\item Vertical droit + crochet (bras droit du \zh{寸})
\item Point final (en bas à droite)
\end{enumerate}
\emph{Attention :} on n'écrit pas le crochet central (trait 5) en tout premier ; on achève le cadre supérieur (traits 1--4) avant de descendre au centre.
Mots : \zh{继承}, \zh{承认} (chéngrèn, admettre), \zh{承担} (chéngdān, assumer).

\textbf{嘱 (zhǔ) — ordonner, recommander, confier ; 12 traits (écriture simplifiée courante).}\par
Structure : gauche-droite. Radical \zh{口} (kǒu, bouche, 3 traits) à gauche ; \zh{属} (shǔ, 9 traits simplifiés) à droite. 
Ordre : 1--3. \zh{口} (vertical, horizontal+crochet, horizontal). 4--12. \zh{属} : d'abord \zh{尸} (shī, corps, 3 traits : horizontal+crochet, horizontal, trait tombant), puis la partie basse \zh{朱} (zhū, 6 traits : horizontal, vertical, horizontal+crochet, horizontal, horizontal, vertical). 
C'est le caractère le plus complexe du thème : décomposez-le impérativement en \zh{口} + \zh{尸} + \zh{朱}. 
Mots : \zh{遗嘱} (yízhǔ), \zh{嘱咐} (zhǔfu, recommander).

\textbf{财 (cái) — richesse, argent ; 7 traits.}\par
Structure : gauche-droite. Radical \zh{贝} (bèi, coquillage/monnaie, 4 traits) à gauche ; \zh{才} (cái, talent, 3 traits) à droite. 
Ordre : 1--4. \zh{贝} (horizontal, vertical, horizontal+crochet, horizontal — \emph{pas de trait tombant final pour 贝 seul}). 5--7. \zh{才} (horizontal, vertical+crochet, trait tombant). 
Sens : le coquillage-monnaie (\zh{贝}) + le talent (\zh{才}) = la richesse. 
Mots : \zh{财产}, \zh{财富} (cáifù), \zh{金融} (jīnróng, finance — bien que sans 财, même radical).

\textbf{法 (fǎ) — loi, méthode, modèle ; 8 traits.}\par
Structure : gauche-droite. Radical \zh{氵} (shuǐ, eau, 3 traits) à gauche ; \zh{去} (qù, aller, 5 traits) à droite. 
Ordre : 1--3. \zh{氵} (deux points puis trait remontant \zh{㇀}). 4--8. \zh{去} (point, horizontal, trait tombant, horizontal, vertical crochu). 
Étymologie : l'eau qui coule selon une règle, un modèle. 
Mots : \zh{法律} (fǎlǜ), \zh{办法} (bànfǎ, solution/moyen), \zh{法语} (fǎyǔ, français).

\begin{center}
\begin{tabularx}{\linewidth}{|X|X|X|X|}
\hline
\rowcolor{popblueL} Caractère & Pinyin & Traits (simplifié) & Ordre clé et astuce \\
\hline
\zh{遗} & yí & 12 (9+3) & \zh{贵} d'abord (haut→bas), clé \zh{辶} en dernier (point, zhé, nà) \\
\hline
\zh{产} & chǎn & 6 & Point, horizontal, 2 points, horizontal long, \zh{丿} final \\
\hline
\zh{继} & jì & 10 (3+7) & \zh{纟} d'abord (g→d), puis partie droite (haut→bas) \\
\hline
\zh{承} & chéng & 8 & Cadre haut (1-4), crochet central (5), côtés bas (6-7), point (8) \\
\hline
\zh{嘱} & zhǔ & 12 (3+9) & \zh{口} (g), puis \zh{属} : \zh{尸} (haut) puis \zh{朱} (bas) \\
\hline
\zh{财} & cái & 7 (4+3) & \zh{贝} (g), puis \zh{才} (d) \\
\hline
\zh{法} & fǎ & 8 (3+5) & \zh{氵} (g : 2 points + ti), puis \zh{去} (d) \\
\hline
\end{tabularx}
\end{center}

\courssub{S'entraîner sur la grille 米字格}

Pour écrire des caractères équilibrés, les élèves chinois s'entraînent sur la \cle{grille 米字格} (mǐzìgé, « grille du caractère 米 ») : un carré partagé par les médianes et les diagonales, qui aide à placer chaque trait au bon endroit. Le centre du caractère doit occuper le centre de la grille ; les composants gauche et droit se partagent les moitiés gauche et droite.

\begin{popfigure}
\begin{tikzpicture}[scale=1.1]
% Grille 1 : 遗
\draw[popline, thick] (0,0) rectangle (3,3);
\draw[popline, dashed] (0,1.5) -- (3,1.5);
\draw[popline, dashed] (1.5,0) -- (1.5,3);
\draw[popline, dashed] (0,0) -- (3,3);
\draw[popline, dashed] (0,3) -- (3,0);
\node at (1.5,1.5) {\Huge\zh{遗}};
\node[below, popmuted] at (1.5,0) {\small 1--9 : \zh{贵} \quad 10--12 : \zh{辶}};
% Grille 2 : 承
\begin{scope}[xshift=4cm]
\draw[popline, thick] (0,0) rectangle (3,3);
\draw[popline, dashed] (0,1.5) -- (3,1.5);
\draw[popline, dashed] (1.5,0) -- (1.5,3);
\draw[popline, dashed] (0,0) -- (3,3);
\draw[popline, dashed] (0,3) -- (3,0);
\node at (1.5,1.5) {\Huge\zh{承}};
\node[below, popmuted] at (1.5,0) {\small 1-4 haut, 5 centre, 6-7 bas, 8 point};
\end{scope}
% Grille 3 : 嘱
\begin{scope}[xshift=8cm]
\draw[popline, thick] (0,0) rectangle (3,3);
\draw[popline, dashed] (0,1.5) -- (3,1.5);
\draw[popline, dashed] (1.5,0) -- (1.5,3);
\draw[popline, dashed] (0,0) -- (3,3);
\draw[popline, dashed] (0,3) -- (3,0);
\node at (1.5,1.5) {\Huge\zh{嘱}};
\node[below, popmuted] at (1.5,0) {\small 1-3 : \zh{口} \quad 4-12 : \zh{属}};
\end{scope}
% Grille 4 : 继
\begin{scope}[xshift=12cm]
\draw[popline, thick] (0,0) rectangle (3,3);
\draw[popline, dashed] (0,1.5) -- (3,1.5);
\draw[popline, dashed] (1.5,0) -- (1.5,3);
\draw[popline, dashed] (0,0) -- (3,3);
\draw[popline, dashed] (0,3) -- (3,0);
\node at (1.5,1.5) {\Huge\zh{继}};
\node[below, popmuted] at (1.5,0) {\small 1-3 : \zh{纟} \quad 4-10 : partie dr.};
\end{scope}
\end{tikzpicture}
\legende{Grilles d'écriture 米字格 pour 遗, 承, 嘱 et 继, avec rappel de l'ordre des traits.}
\end{popfigure}

\begin{methode}[Mémoriser un caractère difficile en cinq étapes]
\begin{enumerate}
\item \cle{Décomposer} le caractère en composants connus : \zh{嘱} = \zh{口} + \zh{尸} + \zh{朱} ; \zh{继} = \zh{纟} + \text{(partie droite)}. Chaque partie a un sens ou un son qui aide la mémoire.
\item \cle{Compter les traits} et apprendre leur ordre exact (ex: 12 traits pour \zh{嘱}).
\item \cle{Écrire 5 fois} le caractère sur une grille 米字格, en comptant les traits à voix haute (\"un, deux, trois...\") pour ancrer l'ordre moteur.
\item \cle{Réécrire de mémoire}, puis comparer avec le modèle et corriger les proportions (largeur gauche/droite, hauteur haut/bas).
\item \cle{Réutiliser} le caractère dans un mot et une phrase complète : \zh{爷爷立了遗嘱。} (Yéye lì le yízhǔ. — « Grand-père a rédigé un testament. »)
\end{enumerate}
\end{methode}

\courssub{Ne pas confondre les caractères proches (pièges d'examen)}

Certains caractères se ressemblent beaucoup mais ont des sens très différents. Au baccalauréat, les confondre coûte des points en écriture comme en traduction.

\begin{center}
\begin{tabularx}{\linewidth}{|X|X|X|X|}
\hline
\rowcolor{popblueL} Paire & Pinyin & Différence graphique & Sens \\
\hline
\zh{遗} / \zh{遣} & yí / qiǎn & Intérieur : \zh{贵} (précieux) vs \zh{遣} (intérieur diff.) & léguer / envoyer, dépêcher \\
\hline
\zh{嘱} / \zh{属} & zhǔ / shǔ & \zh{嘱} a \zh{口} (bouche) à g. ; \zh{属} n'en a pas & recommander / appartenir, catégorie \\
\hline
\zh{财} / \zh{材} & cái / cái & \zh{财} : \zh{贝} (coquillage) ; \zh{材} : \zh{木} (bois) & richesse, argent / matériau, bois, talent \\
\hline
\zh{承} / \zh{永} & chéng / yǒng & \zh{承} a 3 horizontaux au centre + bas \zh{寸} ; \zh{永} a 1 horizontal + traits calligraphiques & hériter, supporter / éternel, toujours \\
\hline
\zh{继} / \zh{续} & jì / xù & \zh{继} : \zh{纟} (soie) ; \zh{续} : \zh{纟} + \zh{卖} (vendre) & succéder, hériter / continuer, prolonger \\
\hline
\end{tabularx}
\end{center}

\begin{attention}[嘱 ou 属 ? La bouche fait toute la différence]
L'erreur la plus fréquente des francophones est d'écrire \zh{遗属} au lieu de \zh{遗嘱}. Retenez l'étymologie : \zh{嘱} s'écrit avec la \cle{bouche} \zh{口} car le testament est d'abord une \emph{parole confiée}, une recommandation orale du défunt (\zh{嘱咐}). \zh{属} (shǔ) signifie « appartenir » : \zh{属于} (shǔyú, appartenir à) ou « catégorie » : \zh{属于}. 

De même, \zh{财} (richesse, avec le coquillage-monnaie \zh{贝}) ne doit pas être confondu avec \zh{材} (matériau, bois, talent, avec la clé du bois \zh{木}) : les deux se prononcent \emph{cái} (homophones parfaits), mais \zh{财产} (les biens, le patrimoine) s'écrit \textbf{toujours} avec \zh{贝}. \zh{材料} (matériaux), \zh{人才} (talent) s'écrivent avec \zh{木}.
\end{attention}

\begin{exemplebox}[Les caractères en contexte : phrases modèles]
\begin{itemize}
\item \zh{立遗嘱} — \emph{lì yízhǔ} — « rédiger un testament » (verbe \zh{立} : établir, dresser).
\item \zh{财产} — \emph{cáichǎn} — « les biens, le patrimoine ».
\item \zh{法律} — \emph{fǎlǜ} — « la loi ».
\item \zh{继承遗产} — \emph{jìchéng yíchǎn} — « hériter d'un patrimoine ».
\item \zh{老人把财产留给了孩子。} — \emph{Lǎorén bǎ cáichǎn liú gěi le háizi.} — « Le vieil homme a laissé ses biens à ses enfants. » (Structure \zh{把}).
\item \zh{根据法律，子女都可以继承遗产。} — \emph{Gēnjù fǎlǜ, zǐnǚ dōu kěyǐ jìchéng yíchǎn.} — « Selon la loi, les enfants peuvent tous hériter du patrimoine. »
\item \zh{父亲嘱咐儿子要努力工作。} — \emph{Fùqīn zhǔfu érzi yào nǔlì gōngzuò.} — « Le père a recommandé à son fils de travailler dur. »
\end{itemize}
\end{exemplebox}

\begin{exoresolu}[Identifier et corriger les erreurs d'écriture]
Énoncé : Un élève de Terminale a écrit la phrase suivante. Deux caractères sont mal choisis ; retrouvez-les, corrigez-les et justifiez. \\
\zh{爷爷立了遗属，把材产留给了孙子。}
\tcblower
Solution détaillée :
\begin{enumerate}
\item \zh{遗属} est incorrect : il faut écrire \zh{遗嘱} (yízhǔ, testament), avec la bouche \zh{口}, car le testament est une parole confiée (\zh{嘱咐}). \zh{属} (shǔ) signifie « appartenir » (ex: \zh{属于}).
\item \zh{材产} est incorrect : il faut écrire \zh{财产} (cáichǎn, les biens), avec \zh{贝} (le coquillage-monnaie). \zh{材} (cái), avec la clé du bois \zh{木}, signifie « matériau, bois, talent ».
\item Phrase corrigée : \zh{爷爷立了遗嘱，把财产留给了孙子。} — \emph{Yéye lì le yízhǔ, bǎ cáichǎn liú gěi le sūnzi.} — « Grand-père a rédigé un testament et a laissé ses biens à son petit-fils. »
\end{enumerate}
\end{exoresolu}

\begin{exercice}[À vous d'écrire]
\begin{enumerate}
\item Écrivez 5 fois chacun des caractères \zh{遗}, \zh{承}, \zh{嘱} et \zh{继} sur une grille 米字格, en comptant les traits à voix haute.
\item Recopiez la phrase \zh{根据法律，子女都可以继承遗产。} en veillant scrupuleusement à l'ordre des traits de \zh{遗}, \zh{继} et \zh{承}.
\item Complétez avec le bon caractère (\zh{嘱} ou \zh{属} ; \zh{财} ou \zh{材}) :\\
  a) \zh{遗\quad} (testament) \quad b) \zh{\quad 于} (appartenir à) \quad c) \zh{\quad 产} (les biens) \quad d) \zh{材\quad} (le matériau).
\item Dictée de caractères (le professeur dicte, l'élève écrit) : \zh{遗嘱}, \zh{继承}, \zh{财产}, \zh{法律}, \zh{嘱咐}, \zh{属于}.
\end{enumerate}
\end{exercice}

\begin{corrige}[Exercice : À vous d'écrire]
1. Vérification des traits : \zh{遗} = 12 traits (\zh{贵} puis \zh{辶}) ; \zh{承} = 8 traits (cadre haut, crochet central, côtés bas, point) ; \zh{嘱} = 12 traits (\zh{口} puis \zh{属}) ; \zh{继} = 10 traits (\zh{纟} puis partie droite).\\
2. Phrase recopiée : \zh{根据法律，子女都可以继承遗产。} — « Selon la loi, les enfants peuvent tous hériter du patrimoine. »\\
3. a) \zh{遗嘱} (testament) ; b) \zh{属于} (appartenir à) ; c) \zh{财产} (les biens) ; d) \zh{材料} (le matériau).\\
4. Dictée : vérifier l'absence de confusion \zh{嘱}/\zh{属} et \zh{财}/\zh{材}, et l'ordre des traits de \zh{遗} (\zh{辶} dernier) et \zh{承} (centre avant côtés bas).
\end{corrige}

\begin{aretenir}[L'essentiel de la section — Fiche de révision Bac]
\begin{itemize}
\item \cle{Clé 辶 (mouvement)} : \textbf{toujours en dernier} (\zh{遗}, \zh{这}, \zh{进}, \zh{近}, \zh{遣}). Ordre : intérieur → point → zhé → nà.
\item \cle{Structure gauche-droite} : \textbf{gauche avant droite} (\zh{财}, \zh{嘱}, \zh{法}, \zh{继}, \zh{属}). Radical gauche souvent sémantique (\zh{贝}, \zh{口}, \zh{氵}, \zh{纟}).
\item \cle{Structure haut-bas / symétrique} : \textbf{haut avant bas} ; pour \zh{承}, cadre haut (1-4) → \textbf{crochet central (5)} → côtés bas (6-7) → point (8).
\item \cle{Pièges homophones / graphiques} : 
      \begin{itemize}
      \item \zh{遗嘱} (testament, bouche \zh{口}) $\neq$ \zh{遗属} (incorrect, \zh{属} = appartenir).
      \item \zh{财产} (biens, coquillage \zh{贝}) $\neq$ \zh{材产} (incorrect, \zh{材} = bois/matériau).
      \item \zh{继承} (hériter, soie \zh{纟}) $\neq$ \zh{继续} (continuer).
      \end{itemize}
\item \cle{Méthode 5 étapes} : Décomposer → Compter/Ordonner → Écrire 5× (grille + oral) → Mémoire + Correction → Contextualiser (mot + phrase).
\end{itemize}
\end{aretenir}

\coursec{Structure et connecteurs du texte}

Dans la section précédente, nous avons appris à écrire correctement les caractères du thème de l'héritage : \zh{遗产} (yíchǎn), \zh{遗嘱} (yízhǔ), \zh{继承} (jìchéng), \zh{分配} (fēnpèi), \zh{法律} (fǎlǜ), \zh{传统} (chuántǒng), en respectant l'ordre des traits (haut → bas, gauche → droite, horizontal avant vertical, extérieur avant intérieur) et en identifiant les radicaux (\zh{辶} pour 遗, \zh{立} pour 产, \zh{讠} pour 嘱, \zh{纟} pour 继/统, \zh{扌} pour 承, \zh{氵} pour 法, \zh{彳} pour 律, \zh{亻} pour 传). Mais un bon texte n'est pas seulement une suite de caractères bien tracés : c'est un ensemble de phrases \cle{organisées} et \cle{reliées}. Dans cette section, nous allons apprendre le plan type d'un texte sur l'héritage et les connecteurs qui donnent du rythme et de la logique à votre production écrite. C'est exactement ce que le correcteur du Bac attend : un texte clair, structuré, avec des liaisons visibles.

\begin{aretenir}[Rappel : Écriture des caractères clés du thème]
\begin{itemize}
\item \zh{遗} (yí) : radical \zh{辶} (chuò, marche), 13 traits. Ordre : haut (艹) → milieu (贵) → bas (辶).
\item \zh{产} (chǎn) : radical \zh{立} (lì, debout), 11 traits. Ordre : haut (生) → bas (立).
\item \zh{嘱} (zhǔ) : radical \zh{讠} (yán, parole), 15 traits. Gauche (讠) → droite (主 + 口).
\item \zh{继} (jì) : radical \zh{纟} (sī, fil), 10 traits. Gauche (纟) → droite (米 + 幂).
\item \zh{承} (chéng) : radical \zh{扌} (shǒu, main), 8 traits. Gauche (扌) → droite (丰 + 寸).
\item \zh{法} (fǎ) : radical \zh{氵} (shuǐ, eau), 8 traits. Gauche (氵) → droite (去).
\item \zh{律} (lǜ) : radical \zh{彳} (chì, pas), 9 traits. Gauche (彳) → milieu (聿) → droite (彳).
\item \zh{传} (chuán) : radical \zh{亻} (rén, homme), 6 traits. Gauche (亻) → droite (专).
\item \zh{统} (tǒng) : radical \zh{纟} (sī), 7 traits. Gauche (纟) → droite (充).
\end{itemize}
\end{aretenir}

\courssub{Observer d'abord : un texte modèle}

Lisons d'abord un court texte original. Soulignez mentalement les mots qui relient les phrases entre elles.

\begin{textebox}[爷爷的遗产 — L'héritage de grand-père]
我的爷爷去年去世了。\\(Wǒ de yéye qùnián qùshì le. — Mon grand-père est décédé l'année dernière.)\\[4pt]
首先，他没有立遗嘱，所以家人不知道怎么办。\\(Shǒuxiān, tā méiyǒu lì yízhǔ, suǒyǐ jiārén bù zhīdào zěnme bàn. — D'abord, il n'avait pas fait de testament, donc la famille ne savait pas quoi faire.)\\[4pt]
然后，我们开了一个家庭会议。按照法律，房子由大儿子继承，钱由三个孩子平分。\\(Ránhòu, wǒmen kāi le yí ge jiātíng huìyì. Ànzhào fǎlǜ, fángzi yóu dà érzi jìchéng, qián yóu sān ge háizi píngfēn. — Ensuite, nous avons tenu une réunion de famille. Selon la loi, la maison est héritée par le fils aîné, l'argent est partagé également entre les trois enfants.)\\[4pt]
爷爷虽然不在了，但是他给我们留下了很多美好的回忆。\\(Yéye suīrán bú zài le, dànshì tā gěi wǒmen liúxià le hěn duō měihǎo de huíyì. — Bien que grand-père ne soit plus là, il nous a laissé beaucoup de beaux souvenirs.)\\[4pt]
最后，爸爸把爷爷的手表送给了我，因为我最喜欢爷爷。\\(Zuìhòu, bàba bǎ yéye de shǒubiǎo sòng gěi le wǒ, yīnwèi wǒ zuì xǐhuan yéye. — Enfin, papa m'a offert la montre de grand-père, parce que j'aimais le plus grand-père.)\\[4pt]
因此，我明白了：遗产不只是钱和房子，还有爱。\\(Yīncǐ, wǒ míngbai le: yíchǎn bù zhǐ shì qián hé fángzi, hái yǒu ài. — C'est pourquoi j'ai compris : l'héritage, ce n'est pas seulement l'argent et la maison, il y a aussi l'amour.)
\end{textebox}

\textbf{Questions de compréhension :}
\begin{enumerate}
\item 爷爷立遗嘱了吗？(Yéye lì yízhǔ le ma? — Grand-père a-t-il fait un testament ?)
\item 房子由谁继承？(Fángzi yóu shéi jìchéng? — Par qui la maison est-elle héritée ?)
\item 作者明白了什么？(Zuòzhě míngbai le shénme? — Qu'a compris l'auteur ?)
\end{enumerate}

\textbf{Glossaire :}
\begin{center}
\begin{tabularx}{\linewidth}{|X|X|X|X|}
\hline
\rowcolor{popblueL} Caractères & Pinyin & Nature & Traduction \\
\hline
去世 & qùshì & verbe & décéder (terme respectueux) \\
\hline
立遗嘱 & lì yízhǔ & loc. verbale & faire un testament \\
\hline
家庭会议 & jiātíng huìyì & nom & réunion de famille \\
\hline
平分 & píngfēn & verbe & partager en parts égales \\
\hline
回忆 & huíyì & nom & souvenir \\
\hline
手表 & shǒubiǎo & nom & montre \\
\hline
明白 & míngbai & verbe & comprendre \\
\hline
\end{tabularx}
\end{center}

Observez : le texte suit toujours le même mouvement en \cle{trois parties}, et chaque phrase est introduite ou reliée par un connecteur. C'est cette architecture que nous allons maintenant étudier en détail.

\courssub{Le plan type d'un texte sur l'héritage}

\begin{propriete}[Le plan en trois parties]
Tout texte narratif ou explicatif sur l'héritage se construit en trois parties :
\begin{enumerate}
\item \textbf{开头 (kāitóu) — l'introduction} : présenter la situation (le décès, la famille, le contexte). Exemple : \zh{我的爷爷去年去世了。}
\item \textbf{中间 (zhōngjiān) — le développement} : décrire les biens, les héritiers, le partage, les discussions. C'est la partie la plus longue.
\item \textbf{结尾 (jiéwěi) — la conclusion} : exprimer les sentiments, tirer une leçon, ouvrir sur une réflexion. Exemple : \zh{遗产不只是钱，还有爱。}
\end{enumerate}
\end{propriete}

\begin{popfigure}
\begin{tikzpicture}[
  bloc/.style={rectangle, rounded corners=6pt, draw=popblue, fill=popblueL, thick, minimum width=3.4cm, minimum height=1.6cm, align=center, font=\small},
  conn/.style={rectangle, rounded corners=4pt, draw=poporange, fill=poporangeL, minimum width=3.2cm, align=center, font=\footnotesize},
  fleche/.style={-{Stealth[length=3mm]}, thick, popdark}
]
\node[bloc, align=center] (intro) at (0,0){\textbf{开头}\\Introduction\\situation, décès};
\node[bloc, fill=poppurpleL, draw=poppurple, align=center] (dev) at (5.2,0){\textbf{中间}\\Développement\\biens, héritiers};
\node[bloc, fill=popgreenL, draw=popgreen, align=center] (conc) at (10.4,0){\textbf{结尾}\\Conclusion\\sentiments, leçon};
\node[conn, align=center] (c1) at (0,-2.2){首先\\shǒuxiān};
\node[conn, align=center] (c2) at (5.2,-2.2){然后 / 按照 / 由\\ránhòu / ànzhào / yóu};
\node[conn, align=center] (c3) at (10.4,-2.2){最后 / 因此\\zuìhòu / yīncǐ};
\draw[fleche] (intro) -- (dev);
\draw[fleche] (dev) -- (conc);
\draw[fleche, poporange] (intro) -- (c1);
\draw[fleche, poporange] (dev) -- (c2);
\draw[fleche, poporange] (conc) -- (c3);
\end{tikzpicture}
\legende{La structure en trois parties et les connecteurs associés à chaque étape.}
\end{popfigure}

\courssub{Les connecteurs chronologiques}

Ces connecteurs ordonnent le récit dans le temps. Ils se placent \cle{en début de phrase} (ou juste après le sujet).

\begin{center}
\begin{tabularx}{\linewidth}{|X|X|X|}
\hline
\rowcolor{popblueL} Connecteur & Exemple (caractères + pinyin) & Traduction \\
\hline
首先 shǒuxiān (d'abord) & 首先，爷爷去世了。Shǒuxiān, yéye qùshì le. & D'abord, grand-père est décédé. \\
\hline
然后 ránhòu (ensuite) & 然后，家人开了会。Ránhòu, jiārén kāi le huì. & Ensuite, la famille a tenu une réunion. \\
\hline
以后 yǐhòu (après) & 爷爷去世以后，我们很难过。Yéye qùshì yǐhòu, wǒmen hěn nánguò. & Après le décès de grand-père, nous étions très tristes. \\
\hline
最后 zuìhòu (enfin) & 最后，我们分了遗产。Zuìhòu, wǒmen fēn le yíchǎn. & Enfin, nous avons partagé l'héritage. \\
\hline
\end{tabularx}
\end{center}

\begin{attention}[以后 : place particulière]
Contrairement à \zh{首先} et \zh{然后} qui ouvrent la phrase, \zh{以后} se place \textbf{après} le groupe verbal ou la proposition qu'il suit : \zh{爷爷去世以后} (après le décès de grand-père), littéralement « grand-père décéder après ». Ne dites jamais \zh{以后爷爷去世} pour « après le décès de grand-père » : cela signifierait « plus tard, grand-père décédera ».
\end{attention}

\courssub{Les connecteurs logiques : cause et concession}

\begin{propriete}[因为...所以... : la cause double]
En chinois, la cause et la conséquence peuvent s'exprimer \textbf{ensemble} dans la même phrase, contrairement au français qui choisit « parce que » \emph{ou} « donc » :\\
\cle{Sujet + 因为 + cause，所以 + conséquence}\\
\zh{因为爷爷没有立遗嘱，所以家人按照法律分配遗产。}\\
\emph{Yīnwèi yéye méiyǒu lì yízhǔ, suǒyǐ jiārén ànzhào fǎlǜ fēnpèi yíchǎn.}\\
« Parce que grand-père n'a pas fait de testament, la famille partage l'héritage selon la loi. »\\
On peut omettre \zh{所以} à l'oral, mais garder les deux est plus clair à l'écrit et très apprécié au Bac.
\end{propriete}

\begin{propriete}[虽然...但是... : la concession double]
En français, on dit « \emph{bien qu'}il soit pauvre, il est heureux » : un seul mot de liaison. En chinois, \textbf{les deux parties s'expriment obligatoirement} :\\
\cle{虽然 + concession，但是 + constat}\\
\zh{虽然他很穷，但是他很快乐。}\\
\emph{Suīrán tā hěn qióng, dànshì tā hěn kuàilè.}\\
« Bien qu'il soit pauvre, il est heureux. »\\
\zh{虽然爷爷不在了，但是我们常常想起他。}\\
\emph{Suīrán yéye bú zài le, dànshì wǒmen chángcháng xiǎngqǐ tā.}\\
« Bien que grand-père ne soit plus là, nous pensons souvent à lui. »\\
\textbf{Ne jamais omettre 但是} après \zh{虽然} : c'est l'erreur la plus fréquente des francophones.
\end{propriete}

\begin{exemplebox}[因此 : la conclusion logique]
\zh{因此} (yīncǐ, « c'est pourquoi ») introduit la conséquence finale, souvent en conclusion de texte :\\
\zh{爷爷没有立遗嘱，因此家人按照法律分配遗产。}\\
\emph{Yéye méiyǒu lì yízhǔ, yīncǐ jiārén ànzhào fǎlǜ fēnpèi yíchǎn.}\\
« Grand-père n'a pas fait de testament, c'est pourquoi la famille partage l'héritage selon la loi. »
\end{exemplebox}

\courssub{Les connecteurs d'addition}

\begin{center}
\begin{tabularx}{\linewidth}{|X|X|X|}
\hline
\rowcolor{popblueL} Connecteur & Exemple (caractères + pinyin) & Traduction \\
\hline
还有 hái yǒu (il y a aussi) & 遗产里有房子，还有钱。Yíchǎn lǐ yǒu fángzi, hái yǒu qián. & Dans l'héritage, il y a une maison, et il y a aussi de l'argent. \\
\hline
而且 érqiě (de plus) & 房子很大，而且很漂亮。Fángzi hěn dà, érqiě hěn piàoliang. & La maison est grande, et de plus elle est très belle. \\
\hline
也 yě (aussi) & 哥哥也继承了遗产。Gēge yě jìchéng le yíchǎn. & Le grand frère aussi a hérité. \\
\hline
\end{tabularx}
\end{center}

\begin{attention}[Place de 也]
\zh{也} se place \textbf{avant le verbe}, jamais en fin de phrase comme le français « aussi » : on dit \zh{我也喜欢爷爷} (moi aussi j'aime grand-père), jamais \zh{我喜欢爷爷也}. De même, \zh{而且} relie deux adjectifs ou deux verbes, tandis que \zh{还有} ajoute un nom ou une phrase entière.
\end{attention}

\courssub{La structure 按照 + nom : « selon »}

\begin{propriete}[按照 : « selon, conformément à »]
\cle{按照 + nom + verbe} exprime la conformité à une règle, une loi, un document :\\
\zh{按照法律分配遗产。} \emph{Ànzhào fǎlǜ fēnpèi yíchǎn.} « Partager l'héritage selon la loi. »\\
\zh{按照遗嘱，房子由大儿子继承。} \emph{Ànzhào yízhǔ, fángzi yóu dà érzi jìchéng.} « Selon le testament, la maison est héritée par le fils aîné. »\\
\zh{按照中国的传统，儿子继承父亲的姓。} \emph{Ànzhào Zhōngguó de chuántǒng, érzi jìchéng fùqin de xìng.} « Selon la tradition chinoise, le fils hérite du nom du père. »
\end{propriete}

\begin{attention}[La place de 按照]
\zh{按照} se place \textbf{avant le nom} auquel il se rapporte, en tête de phrase ou de proposition, et \textbf{jamais en fin de phrase}. On dit \zh{按照法律分配} et non \zh{分配法律按照}. Le francophone est tenté de calquer « partager l'héritage \emph{selon la loi} » en plaçant le complément à la fin : en chinois, ce complément de conformité vient \emph{avant} le verbe.
\end{attention}

\courssub{La tournure passive 由 + personne + verbe}

\begin{propriete}[由 : « par » dans la tournure passive]
Pour mettre en avant le bien hérité plutôt que la personne, le chinois écrit utilise :\\
\cle{Chose + 由 + personne + verbe}\\
\zh{财产由儿子继承。} \emph{Cáichǎn yóu érzi jìchéng.} « Les biens sont hérités par le fils. »\\
\zh{房子由奶奶继承。} \emph{Fángzi yóu nǎinai jìchéng.} « La maison est héritée par grand-mère. »\\
Cette tournure est élégante à l'écrit et très utile dans un texte sur l'héritage, car elle permet de varier : « le fils hérite de la maison » / « la maison est héritée par le fils ».
\end{propriete}

\courssub{Rappel : la structure 把}

Vous avez déjà rencontré la structure \zh{把}. Elle est précieuse pour dire qu'on « laisse » ou « transmet » un bien à quelqu'un :

\begin{propriete}[Sujet + 把 + objet + verbe + complément]
\cle{Sujet + 把 + objet + verbe + 给 + bénéficiaire}\\
\zh{爸爸把房子留给了奶奶。}\\
\emph{Bàba bǎ fángzi liú gěi le nǎinai.}\\
« Papa a laissé la maison à grand-mère. »\\
\zh{爷爷把土地分给了三个孩子。}\\
\emph{Yéye bǎ tǔdì fēn gěi le sān ge háizi.}\\
« Grand-père a partagé la terre entre les trois enfants. »\\
Le verbe ne peut jamais rester seul après \zh{把} : il faut un complément (\zh{给了}, \zh{留给}, \zh{分给}…).
\end{propriete}

\begin{methode}[Rédiger un texte sur l'héritage en 4 étapes]
\begin{enumerate}
\item \textbf{Écrire le plan} en trois parties : \zh{开头} (situation), \zh{中间} (biens et héritiers), \zh{结尾} (sentiments, leçon).
\item \textbf{Choisir trois connecteurs différents}, par exemple \zh{首先} / \zh{然后} / \zh{最后}, plus un connecteur logique (\zh{因为...所以...} ou \zh{虽然...但是...}).
\item \textbf{Rédiger phrase par phrase}, en plaçant chaque connecteur à la bonne place (\zh{也} avant le verbe, \zh{以后} après le groupe, \zh{按照} avant le nom).
\item \textbf{Relire en vérifiant chaque caractère} (traits, radicaux) et chaque structure (\zh{把} suivi d'un complément, \zh{但是} présent après \zh{虽然}).
\end{enumerate}
\end{methode}

\begin{exoresolu}[Relier les phrases avec le bon connecteur]
Énoncé : complétez avec \zh{首先}, \zh{然后}, \zh{最后}, \zh{因为...所以...}, \zh{虽然...但是...}, \zh{按照}, \zh{由}.\\
1. \zh{＿＿＿爷爷去世了，＿＿＿家人很伤心。}\\
2. \zh{＿＿＿，我们开了一个家庭会议。}\\
3. \zh{＿＿＿法律，房子＿＿＿大儿子继承。}\\
4. \zh{＿＿＿爷爷不在了，＿＿＿他的爱还在。}
\tcblower
Solution détaillée :\\
1. \zh{因为爷爷去世了，所以家人很伤心。} (Yīnwèi yéye qùshì le, suǒyǐ jiārén hěn shāngxīn. — Parce que grand-père est décédé, la famille est très triste.) Cause + conséquence : \zh{因为...所以...}\\
2. \zh{然后，我们开了一个家庭会议。} (Ránhòu, wǒmen kāi le yí ge jiātíng huìyì. — Ensuite, nous avons tenu une réunion de famille.) Étape intermédiaire du récit : \zh{然后}.\\
3. \zh{按照法律，房子由大儿子继承。} (Ànzhào fǎlǜ, fángzi yóu dà érzi jìchéng. — Selon la loi, la maison est héritée par le fils aîné.) \zh{按照} avant le nom \zh{法律} ; tournure passive avec \zh{由}.\\
4. \zh{虽然爷爷不在了，但是他的爱还在。} (Suīrán yéye bú zài le, dànshì tā de ài hái zài. — Bien que grand-père ne soit plus là, son amour est toujours là.) Concession : les deux mots \zh{虽然} et \zh{但是} sont obligatoires.
\end{exoresolu}

\courssub{Thème : Français → Chinois}

\begin{exercice}[Thème — L'héritage]
Traduisez les phrases suivantes en chinois (caractères + pinyin). Utilisez le vocabulaire et les structures de la leçon.
\begin{enumerate}
\item Mon père a fait un testament l'année dernière.
\item Selon la loi, la maison est héritée par ma grande sœur.
\item Bien qu'il soit riche, il n'est pas heureux.
\item Grand-père a laissé sa montre à mon petit frère.
\item C'est pourquoi nous devons respecter la tradition.
\end{enumerate}
\end{exercice}

\begin{corrige}[Exercice 1 — Thème]
\begin{enumerate}
\item \zh{我爸爸去年立了遗嘱。} (Wǒ bàba qùnián lì le yízhǔ.)
\item \zh{按照法律，房子由大姐继承。} (Ànzhào fǎlǜ, fángzi yóu dàjiě jìchéng.)
\item \zh{虽然他很有钱，但是他不快乐。} (Suīrán tā hěn yǒuqián, dànshì tā bù kuàilè.)
\item \zh{爷爷把手表留给了弟弟。} (Yéye bǎ shǒubiǎo liú gěi le dìdi.)
\item \zh{因此，我们必须尊重传统。} (Yīncǐ, wǒmen bìxū zūnzhòng chuántǒng.)
\end{enumerate}
\end{corrige}

\courssub{Version : Chinois → Français}

\begin{exercice}[Version — L'héritage]
Traduisez les phrases suivantes en français.
\begin{enumerate}
\item 爷爷去世以后，家人很伤心。
\item 首先，我们要开家庭会议。
\item 遗产里有钱，还有房子。
\item 爸爸把土地分给了三个儿子。
\item 虽然没有遗嘱，但是大家都很公平。
\end{enumerate}
\end{exercice}

\begin{corrige}[Exercice 2 — Version]
\begin{enumerate}
\item Après le décès de grand-père, la famille était très triste.
\item D'abord, nous devons tenir une réunion de famille.
\item Dans l'héritage, il y a de l'argent, et il y a aussi une maison.
\item Papa a partagé la terre entre les trois fils.
\item Bien qu'il n'y ait pas de testament, tout le monde est très juste / équitable.
\end{enumerate}
\end{corrige}

\courssub{Grille d'évaluation et entraînement Bac}

\begin{methode}[Grille d'auto-évaluation pour l'expression écrite (Bac)]
Notez votre production sur 20 points en vérifiant chaque critère :
\begin{center}
\begin{tabularx}{\linewidth}{|X|X|X|}
\hline
\rowcolor{popblueL} Critère & Indicateurs & Points \\
\hline
\cle{Respect du plan} & 3 parties visibles : \zh{开头} / \zh{中间} / \zh{结尾} & 4 \\
\hline
\cle{Connecteurs} & Au moins 3 connecteurs chronologiques (\zh{首先/然后/最后}) + 1 logique (\zh{因为...所以...} ou \zh{虽然...但是...}) bien placés & 4 \\
\hline
\cle{Structures grammaticales} & \zh{按照} + nom avant verbe ; \zh{由} (passif) ; \zh{把} + O + V + \zh{给} ; \zh{也} avant verbe & 4 \\
\hline
\cle{Vocabulaire thème} & Emploi correct de : \zh{遗产, 遗嘱, 继承, 分配, 法律, 传统, 平分, 回忆} & 3 \\
\hline
\cle{Caractères \& Ponctuation} & Caractères corrects (traits, radicaux), ponctuation chinoise pleine chasse (，。；) & 3 \\
\hline
\cle{Cohérence \& Longueur} & Texte cohérent, 80--100 caractères environ, ton respectueux & 2 \\
\hline
\end{tabularx}
\end{center}
\textbf{Conseil Bac :} Entraînez-vous à écrire un texte complet en 20 minutes chrono. Relisez-vous systématiquement avec cette grille.
\end{methode}

\begin{aretenir}[L'essentiel de la section]
\begin{itemize}
\item Un texte sur l'héritage suit le plan \cle{开头 → 中间 → 结尾} : situation, partage, leçon.
\item Chronologie : \zh{首先} / \zh{然后} / \zh{最后} en tête de phrase ; \zh{以后} après le groupe verbal.
\item Logique : \zh{因为...所以...} et \zh{虽然...但是...} s'emploient \textbf{par paires} en chinois ; \zh{因此} conclut.
\item Addition : \zh{还有}, \zh{而且}, \zh{也} (avant le verbe).
\item \zh{按照 + nom} avant le verbe ; tournure passive \cle{chose + 由 + personne + verbe} ; \zh{把 + objet + verbe + 给 + bénéficiaire}.
\item Écriture : maîtriser traits, radicaux et ordre des traits pour \zh{遗产, 遗嘱, 继承, 分配, 法律, 传统}.
\end{itemize}
\end{aretenir}

\coursec{Thème : français → chinois}

Dans la section précédente, nous avons étudié la structure du texte sur l'héritage et ses connecteurs (\zh{首先}、\zh{然后}、\zh{虽然……但是……}、\zh{所以}…). Nous savons donc désormais \emph{organiser} un texte. Il reste une compétence décisive pour le baccalauréat : la \cle{traduction du français vers le chinois}, appelée \cle{thème}. C'est l'exercice le plus exigeant, car il mobilise à la fois le vocabulaire, la grammaire et l'écriture des caractères. Cette section vous donne une méthode en quatre étapes, puis l'applique à cinq phrases modèles, de la plus simple à la plus difficile.

\courssub{Observer avant de traduire : deux phrases à comparer}

Avant toute règle, observons. Lisez attentivement les deux phrases suivantes et repérez ce qui change de place entre le français et le chinois.

\begin{exemplebox}[Observation : où se place le complément de temps ?]
\zh{我的爷爷去年去世了。}\\
\emph{Wǒ de yéye qùnián qùshì le.}\\
« Mon grand-père est décédé l'année dernière. »\\
\medskip
\zh{他留下了一所房子和一些钱。}\\
\emph{Tā liúxià le yī suǒ fángzi hé yīxiē qián.}\\
« Il a laissé une maison et de l'argent. »
\end{exemplebox}

Que remarque-t-on ?

\begin{itemize}
\item En français, le complément de temps « l'année dernière » se place \emph{après} le verbe ; en chinois, \zh{去年} se place \emph{avant} le verbe \zh{去世}.
\item Le français dit « une maison » ; le chinois intercale un \cle{classificateur} entre le nombre et le nom : \zh{一所房子}.
\item « De l'argent » (quantité indéfinie) se rend par \zh{一些钱}, littéralement « quelque-argent ».
\item Le passé est marqué par \zh{了} placé juste après le verbe : \zh{去世了}, \zh{留下了}.
\end{itemize}

Toute la méthode du thème découle de ces quatre observations.

\courssub{La méthode de traduction en quatre étapes}

\begin{methode}[Traduire une phrase du français vers le chinois]
\begin{enumerate}
\item \textbf{Repérer S–V–C.} Soulignez dans la phrase française le \cle{sujet}, le \cle{verbe} et le(s) \cle{complément}(s). Identifiez aussi les compléments de temps et de manière.
\item \textbf{Choisir le mot exact du glossaire.} Pour chaque élément, prenez le mot chinois appris dans le lexique du thème (ne traduisez jamais « au feeling » : \zh{去世} « décéder » n'est pas \zh{死} dans un texte soutenu).
\item \textbf{Ordonner à la chinoise.} Placez les éléments selon l'ordre chinois : Sujet + complément de temps + verbe (+ \zh{了}) + complément d'objet avec son classificateur.
\item \textbf{Vérifier les caractères.} Relisez chaque caractère : traits, radicaux, homophones (\zh{试} / \zh{式} / \zh{事}), et vérifiez le classificateur.
\end{enumerate}
\end{methode}

\begin{popfigure}
\begin{tikzpicture}[
  node distance=0.55cm,
  box/.style={draw=popblue, fill=popblueL, rounded corners=3pt, align=center, text width=4.6cm, inner sep=5pt, font=\small},
  arr/.style={-{Stealth[length=2.5mm]}, thick, popdark}
]
\node[box, align=center] (p1){\textbf{1. Phrase française}\\ lire et comprendre le sens};
\node[box, below=of p1, align=center] (p2){\textbf{2. Repérage S--V--C}\\ sujet / verbe / compléments\\ (temps, objet, manière)};
\node[box, below=of p2, align=center] (p3){\textbf{3. Choix du vocabulaire}\\ mot exact du glossaire\\ + classificateur adapté};
\node[box, below=of p3, align=center] (p4){\textbf{4. Ordre des mots chinois}\\ S + temps + V(了) + C};
\node[box, below=of p4, align=center] (p5){\textbf{5. Vérification}\\ caractères, tons, classificateur};
\draw[arr] (p1) -- (p2);
\draw[arr] (p2) -- (p3);
\draw[arr] (p3) -- (p4);
\draw[arr] (p4) -- (p5);
\end{tikzpicture}
\legende{Organigramme de la méthode de traduction (thème) : de la phrase française à la phrase chinoise vérifiée.}
\end{popfigure}

\courssub{La règle d'or : le complément de temps avant le verbe}

\begin{propriete}[Place du complément de temps]
En chinois, le complément de temps se place \textbf{avant le verbe} (juste après le sujet), contrairement au français où il se place souvent après.\\
\medskip
\textbf{Structure :} Sujet + complément de temps + verbe (+ 了) + complément\\
\medskip
\zh{我的爷爷去年去世了。} — \emph{Wǒ de yéye qùnián qùshì le.} — « Mon grand-père est décédé l'année dernière. »\\
\zh{他两年前去世了。} — \emph{Tā liǎng nián qián qùshì le.} — « Il est décédé il y a deux ans. »\\
\zh{我们明天分财产。} — \emph{Wǒmen míngtiān fēn cáichǎn.} — « Nous partageons les biens demain. »
\end{propriete}

\begin{attention}[Erreur n° 1 des francophones]
Ne dites \textbf{jamais} \zh{我的爷爷去世了去年} : c'est un calque du français « est décédé l'année dernière ». En chinois, \zh{去年} doit précéder le verbe : \zh{去年去世}. De même, « il y a deux ans » se dit \zh{两年前}, placé avant le verbe : \zh{他两年前去世了}, et non \zh{他去世了两年前}.
\end{attention}

\courssub{Classificateurs et quantités indéfinies}

\begin{propriete}[Le classificateur 所 et l'indéfini 一些]
\begin{itemize}
\item « Une maison » ne se traduit pas par \zh{一房子} : le nom \zh{房子} exige le classificateur \zh{所} (réservé aux bâtiments : maisons, écoles, hôpitaux) : \zh{一所房子} \emph{yì suǒ fángzi}.
\item « De l'argent », quantité indéfinie, se rend par \zh{一些钱} \emph{yìxiē qián} (« une certaine quantité d'argent »). \zh{些} est le classificateur des quantités indéfinies au pluriel.
\item Rappel de la sandhi tonale sur \zh{一} : devant un 4\textsuperscript{e} ton, \zh{一} se prononce \emph{yí} (2\textsuperscript{e} ton) ; devant les 1\textsuperscript{e}, 2\textsuperscript{e} et 3\textsuperscript{e} tons, il se prononce \emph{yì} (4\textsuperscript{e} ton). Comme \zh{所} (suǒ) et \zh{些} (xiē) sont au 3\textsuperscript{e} ton, on prononce \emph{yì suǒ} et \emph{yìxiē}.
\end{itemize}
\end{propriete}

\begin{center}
\begin{tabularx}{\linewidth}{|X|X|X|X|}
\hline
\rowcolor{popblueL} Caractères & Pinyin & Nature & Traduction \\
\hline
去世 & qùshì & verbe & décéder (registre soutenu) \\
\hline
留下 & liúxià & verbe & laisser (en héritage) \\
\hline
房子 & fángzi & nom & maison \\
\hline
一所房子 & yì suǒ fángzi & expression & une maison \\
\hline
一些钱 & yìxiē qián & expression & de l'argent, une somme \\
\hline
按照 & ànzhào & préposition & selon, conformément à \\
\hline
法律 & fǎlǜ & nom & loi, droit \\
\hline
继承 & jìchéng & verbe & hériter \\
\hline
财产 & cáichǎn & nom & biens, patrimoine \\
\hline
伤心 & shāngxīn & adjectif & triste, affligé \\
\hline
团结 & tuánjié & adjectif/verbe & uni, s'unir \\
\hline
立遗嘱 & lì yízhǔ & expression & rédiger un testament \\
\hline
把 & bǎ & préposition & (construction de disposition) \\
\hline
\end{tabularx}
\end{center}

\courssub{Analyse guidée des cinq phrases modèles}

Appliquons maintenant la méthode, phrase par phrase.

\textbf{Phrase 1.} « Mon grand-père est décédé l'année dernière. »

Repérage : S = « mon grand-père » ; V = « est décédé » ; C de temps = « l'année dernière ». Vocabulaire : \zh{我的爷爷} / \zh{去世} / \zh{去年}. Ordre chinois : S + temps + V + \zh{了}.

\zh{我的爷爷去年去世了。} — \emph{Wǒ de yéye qùnián qùshì le.}

Notez le possessif : « mon grand-père (paternel) » = \zh{我的爷爷}. Le \zh{了} final marque l'accompli : le décès est un fait réalisé.

\textbf{Phrase 2.} « Il a laissé une maison et de l'argent. »

Repérage : S = « il » ; V = « a laissé » ; C = double objet « une maison et de l'argent ». Vocabulaire : \zh{他} / \zh{留下} / \zh{一所房子和一些钱}. Le verbe \zh{留下} est suivi de \zh{了} (action accomplie) ; les deux objets sont reliés par \zh{和}.

\zh{他留下了一所房子和一些钱。} — \emph{Tā liúxià le yì suǒ fángzi hé yìxiē qián.}

\textbf{Phrase 3.} « Selon la loi, les enfants héritent des biens. »

Repérage : complément de manière « selon la loi » ; S = « les enfants » ; V = « héritent » ; C = « des biens ». En chinois, le groupe prépositionnel \zh{按照法律} se place en tête de phrase, comme en français — c'est un des rares cas où l'ordre coïncide. « Hériter de » se dit simplement \zh{继承} + nom, sans préposition.

\zh{按照法律，孩子们继承财产。} — \emph{Ànzhào fǎlǜ, háizimen jìchéng cáichǎn.}

\begin{attention}[Pas de préposition après 继承]
Le français dit « hériter \emph{des} biens » ; le chinois dit directement \zh{继承财产}, sans équivalent de « de ». N'écrivez pas \zh{继承的财产}. En revanche, dans la tournure nominale « hériter des biens \emph{du père} », le possessif garde son \zh{的} : \zh{继承父亲的财产}.
\end{attention}

\textbf{Phrase 4.} « Bien que la famille soit triste, elle reste unie. »

Repérage : proposition subordonnée de concession + proposition principale. Le connecteur appris dans la section 4 est \zh{虽然……但是……}. Attention : « la famille » au sens de « les membres de la famille » se dit \zh{家人} ; les deux adjectifs sont précédés de \zh{很} (liaison obligatoire entre sujet et adjectif).

\zh{虽然家人很伤心，但是他们很团结。} — \emph{Suīrán jiārén hěn shāngxīn, dànshì tāmen hěn tuánjié.}

\begin{propriete}[虽然……但是…… : les deux moitiés sont obligatoires]
Contrairement au français, où l'on dit « bien que…, (mais)… » avec un seul connecteur, le chinois emploie \textbf{les deux} : \zh{虽然} en tête de la subordonnée et \zh{但是} en tête de la principale. Omettre \zh{但是} est une erreur fréquente de francophone.
\end{propriete}

\textbf{Phrase 5 (défi).} « Papa a rédigé un testament et a laissé la maison à ma mère. »

Cette phrase exige la construction en \zh{把}, qui « déplace » l'objet avant le verbe pour insister sur ce qu'on en fait. Structure : Sujet + \zh{把} + objet + verbe + \zh{给} + bénéficiaire.

\zh{爸爸立了遗嘱，把房子留给了妈妈。} — \emph{Bàba lì le yízhǔ, bǎ fángzi liúgěi le māma.}

Analyse : \zh{立了遗嘱} « a rédigé un testament » (expression figée : on « dresse » un testament, \zh{立}) ; puis \zh{把房子留给了妈妈} « il a disposé de la maison en la laissant à maman ». Le verbe \zh{留给} « laisser à » est suivi du bénéficiaire.

\begin{aretenir}[La construction en 把]
\textbf{Structure :} Sujet + 把 + objet + verbe + complément (résultat, bénéficiaire, direction).\\
Emploi : quand l'objet subit une action qui le transforme ou le déplace. « Laisser la maison à quelqu'un » = la maison change de mains, d'où \zh{把房子留给了妈妈}. Le français n'a pas d'équivalent : il faut penser « quant à la maison, il l'a laissée à maman ».
\end{aretenir}

\courssub{Tableau comparatif des cinq traductions}

\begin{center}
\begin{tabularx}{\linewidth}{|X|X|X|}
\hline
\rowcolor{popblueL} Français & Chinois + pinyin & Point de grammaire \\
\hline
Mon grand-père est décédé l'année dernière. & \zh{我的爷爷去年去世了。} \emph{Wǒ de yéye qùnián qùshì le.} & temps avant le verbe ; 了 \\
\hline
Il a laissé une maison et de l'argent. & \zh{他留下了一所房子和一些钱。} \emph{Tā liúxià le yì suǒ fángzi hé yìxiē qián.} & classificateur 所 ; 一些 \\
\hline
Selon la loi, les enfants héritent des biens. & \zh{按照法律，孩子们继承财产。} \emph{Ànzhào fǎlǜ, háizimen jìchéng cáichǎn.} & 按照 en tête ; pas de « de » \\
\hline
Bien que la famille soit triste, elle reste unie. & \zh{虽然家人很伤心，但是他们很团结。} \emph{Suīrán jiārén hěn shāngxīn, dànshì tāmen hěn tuánjié.} & double connecteur ; 很 + adj. \\
\hline
Papa a rédigé un testament et a laissé la maison à ma mère. & \zh{爸爸立了遗嘱，把房子留给了妈妈。} \emph{Bàba lì le yízhǔ, bǎ fángzi liúgěi le māma.} & construction en 把 \\
\hline
\end{tabularx}
\end{center}

\courssub{Exercice résolu et entraînement}

\begin{exoresolu}[Thème guidé : « Ma grand-mère a laissé sa maison à mes parents. »]
Énoncé : traduisez en chinois la phrase « Ma grand-mère a laissé sa maison à mes parents. »
\tcblower
\textbf{Étape 1 — Repérage.} S = « ma grand-mère » ; V = « a laissé » ; C = « sa maison » + bénéficiaire « à mes parents ».\\
\textbf{Étape 2 — Vocabulaire.} grand-mère (paternelle) = \zh{奶奶} (nǎinai) ; laisser à = \zh{留给} (liúgěi) ; maison = \zh{房子} ; mes parents = \zh{我的父母} (wǒ de fùmǔ).\\
\textbf{Étape 3 — Ordre.} L'objet « sa maison » subit un transfert : construction en \zh{把}. S + 把 + objet + V + bénéficiaire.\\
\textbf{Étape 4 — Rédaction et vérification.}\\
\zh{奶奶把她的房子留给了我的父母。}\\
\emph{Nǎinai bǎ tā de fángzi liúgěi le wǒ de fùmǔ.}\\
Vérification : \zh{把} avant l'objet, \zh{了} après le verbe, possessif \zh{她的}, caractères \zh{留} (radicaux 卯 + 田) et \zh{遗} non confondus.
\end{exoresolu}

\begin{exercice}[À vous de traduire]
Traduisez en chinois, en appliquant les quatre étapes de la méthode :
\begin{enumerate}
\item « Mon oncle est décédé il y a trois ans. »
\item « Il a laissé deux maisons et de l'argent. »
\item « Selon la tradition, le fils aîné hérite de la maison. »
\item « Bien que nous soyons tristes, nous restons unis. »
\item « Ma mère a rédigé un testament et a laissé l'argent à mes sœurs. »
\end{enumerate}
\end{exercice}

\begin{corrige}[Exercice « À vous de traduire »]
\begin{enumerate}
\item \zh{我的叔叔三年前去世了。} \emph{Wǒ de shūshu sān nián qián qùshì le.} — temps \zh{三年前} avant le verbe. (Note : \zh{叔叔} = oncle paternel cadet ; \zh{舅舅} = oncle maternel).
\item \zh{他留下了两所房子和一些钱。} \emph{Tā liúxià le liǎng suǒ fángzi hé yìxiē qián.} — « deux » devant un classificateur = \zh{两} (liǎng), jamais \zh{二} (èr).
\item \zh{按照传统，大儿子继承房子。} \emph{Ànzhào chuántǒng, dà érzi jìchéng fángzi.}
\item \zh{虽然我们很伤心，但是我们很团结。} \emph{Suīrán wǒmen hěn shāngxīn, dànshì wǒmen hěn tuánjié.} — double connecteur obligatoire.
\item \zh{妈妈立了遗嘱，把钱留给了我的妹妹们。} \emph{Māma lì le yízhǔ, bǎ qián liúgěi le wǒ de mèimeimen.} — construction en \zh{把} ; \zh{妹妹们} = sœurs cadettes.
\end{enumerate}
\end{corrige}

\begin{attention}[Bilan des pièges du thème]
\begin{itemize}
\item Complément de temps \textbf{avant} le verbe : \zh{去年去世}, jamais \zh{去世了去年}.
\item Classificateur obligatoire : \zh{一所房子}, \zh{两所房子} (avec \zh{两}).
\item Pas de préposition après \zh{继承} : \zh{继承财产}.
\item Double connecteur : \zh{虽然……但是……}.
\item Adjectif précédé de \zh{很} : \zh{很伤心}, \zh{很团结}.
\item Transfert d'un bien : penser à la construction en \zh{把}.
\item Sandhi de \zh{一} : \emph{yì} devant 3\textsuperscript{e} ton (\zh{一所} \emph{yì suǒ}, \zh{一些} \emph{yìxiē}).
\end{itemize}
\end{attention}

Vous disposez maintenant d'une méthode complète et de cinq modèles de référence. La section suivante vous entraînera à l'exercice inverse, la \cle{version} : traduire du chinois vers le français, où les difficultés changent de nature.

\coursec{Leçon 7 — Expression écrite : l'héritage}

\courssub{Modèle de texte commenté}

\begin{textebox}[Héritage familial : le testament de grand-père]
我的爷爷去年去世了。\\
\emph{Wǒ de yéye qùnián qùshìle.}\\
« Mon grand-père paternel est décédé l'an dernier. »\\[4pt]
他留下了一份遗嘱，把房子和土地留给了我的父亲和叔叔。\\
\emph{Tā liúxiàle yī fèn yízhǔ, bǎ fángzi hé tǔdì liú gěile wǒ de fùqin hé shūshu.}\\
« Il a laissé un testament, léguant la maison et la terre à mon père et à mon oncle. »\\[4pt]
虽然遗产不多，但家人没有吵架，大家决定一起种咖啡。\\
\emph{Suīrán yíchǎn bù duō, dàn jiārén méiyǒu chǎojià, dàjiā juédìng yìqǐ zhòng kāfēi.}\\
« Bien que l'héritage soit modeste, la famille ne s'est pas disputée ; tous ont décidé de cultiver du café ensemble. »
\end{textebox}

\begin{aretenir}[Analyse du texte modèle]
\begin{itemize}
\item \textbf{Situation de communication} : annonce d'un décès, lecture d'un testament, partage pacifique, projet commun.
\item \textbf{Structure narrative} : 1. Fait déclencheur (\zh{去世}) + temps (\zh{去年}) ; 2. Disposition légale (\zh{遗嘱}, \zh{把} structure) ; 3. Conséquence positive (\zh{虽然……但……}, projet \zh{决定一起……}).
\item \textbf{Registre} : langue standard, vocabulaire HSK 3-4, adapté à une rédaction scolaire ou un récit familial.
\end{itemize}
\end{aretenir}

\courssub{Lexique et glossaire du thème}

\begin{center}
\begin{tabularx}{\linewidth}{|X|X|X|X|}
\hline
\rowcolor{popblueL} Caractères & Pinyin & Nature & Traduction \\
\hline
继承 & jìchéng & verbe & hériter, succéder à \\
\hline
遗产 & yíchǎn & nom & héritage, succession, biens laissés \\
\hline
遗嘱 & yízhǔ & nom & testament \\
\hline
去世 & qùshì & verbe & décéder, mourir (terme respectueux) \\
\hline
留下 & liúxià & verbe & laisser, léguer, conserver \\
\hline
土地 & tǔdì & nom & terre, terrain, foncier \\
\hline
房子 & fángzi & nom & maison, bâtiment \\
\hline
财产 & cáichǎn & nom & biens, patrimoine, fortune \\
\hline
叔叔 & shūshu & nom & oncle (frère cadet du père) \\
\hline
阿姨 & āyí & nom & tante (sœur de la mère) ; dame d'âge mûr \\
\hline
外公 & wàigōng & nom & grand-père maternel \\
\hline
外婆 & wàipó & nom & grand-mère maternelle \\
\hline
爷爷 & yéye & nom & grand-père paternel \\
\hline
奶奶 & nǎinai & nom & grand-mère paternelle \\
\hline
哥哥 & gēge & nom & frère aîné \\
\hline
吵架 & chǎojià & verbe & se disputer, se quereller \\
\hline
按照 & ànzhào & préposition & selon, conformément à, d'après \\
\hline
由 & yóu & préposition & par (agent du passif) \\
\hline
决定 & juédìng & verbe & décider \\
\hline
种 & zhòng & verbe & planter, cultiver \\
\hline
咖啡 & kāfēi & nom & café (plante ou boisson) \\
\hline
一起 & yìqǐ & adv. & ensemble \\
\hline
虽然……但/但是…… & suīrán… dàn/dànshì… & conj. & bien que… (mais)… \\
\hline
以后 & yǐhòu & nom/particule & après ; plus tard (selon position) \\
\hline
份 & fèn & classif. & classif. pour documents (testament, contrat) \\
\hline
\end{tabularx}
\end{center}

\courssub{Écriture correcte des caractères : traits, radicaux, caractères proches}

\begin{definition}[Rappel des règles d'ordre des traits]
\begin{enumerate}
\item Horizontal avant vertical (\zh{十}).
\item Gauche avant droite (\zh{人}).
\item Haut avant bas (\zh{三}).
\item Extérieur avant intérieur (\zh{同}), puis trait de fermeture final.
\item Trait central avant les ailes (\zh{小}).
\item Point en haut à gauche en premier ; point en haut à droite ou intérieur en dernier.
\end{enumerate}
\end{definition}

\begin{center}
\begin{tabularx}{\linewidth}{|X|X|X|X|X|}
\hline
\rowcolor{popblueL} Caractère & Pinyin & Radical (clé) & Nb traits & Caractères proches / Pièges \\
\hline
继 & jì & 纟 (sī, soie) & 10 & \zh{续} (xù, continuer) — même radical, haut différent \\
\hline
承 & chéng & 手 (shǒu, main) → 扌 & 8 & \zh{担} (dān, porter) — même structure, clé main à gauche \\
\hline
遗 & yí & 辶 (chuò, marche) & 12 & \zh{迟} (chí, en retard) — même radical, phonétique différent \\
\hline
产 & chǎn & 立 (lì, debout) & 6 & \zh{生} (shēng, naître) — composant commun en haut \\
\hline
嘱 & zhǔ & 口 (kǒu, bouche) & 15 & \zh{属} (shǔ, appartenir) — même phonétique \zh{蜀}, clé différente \\
\hline
叔 & shū & 又 (yòu, main droite) & 8 & \zh{婶} (shěn, tante paternelle) — même phonétique, clé femme \\
\hline
姨 & yí & 女 (nǚ, femme) & 9 & \zh{姑} (gū, tante paternelle) — même clé, phonétique différent \\
\hline
外 & wài & 夕 (xī, crépuscule) & 5 & \zh{处} (chù, lieu) — haut similaire \\
\hline
咖 & kā & 口 (kǒu, bouche) & 8 & \zh{加} (jiā, ajouter) — phonétique à droite \\
\hline
啡 & fēi & 口 (kǒu, bouche) & 11 & \zh{非} (fēi, non) — phonétique à droite \\
\hline
\end{tabularx}
\end{center}

\begin{methode}[Écrire \zh{继承} (hériter) correctement]
\begin{enumerate}
\item \zh{继} (10 traits) : clé \zh{纟} (3 traits : point, pli, point) + \zh{东} (5 traits : horizontales, verticale, point) + trait final horizontal long. \textbf{Attention} : ne pas confondre avec \zh{续} (haut = \zh{舛}).
\item \zh{承} (8 traits) : clé \zh{扌} (3 traits) + \zh{丙} (5 traits : horizontale, verticale, deux horizontales, verticale finale). La verticale centrale traverse la première horizontale.
\end{enumerate}
\end{methode}

\begin{experience}[Atelier écriture : « Arbre généalogique chinois »]
En binôme, dessinez un arbre généalogique simplifié (3 générations). Étiquetez chaque membre en caractères chinois (\zh{爷爷}, \zh{奶奶}, \zh{外公}, \zh{外婆}, \zh{父亲}, \zh{母亲}, \zh{叔叔}, \zh{阿姨}, \zh{哥哥}, \zh{姐姐}, \zh{我}). Écrivez les caractères au tableau ; la classe vérifie l'ordre des traits et le radical. Durée : 15 min.
\end{experience}

\courssub{Structure et connecteurs du texte}

\begin{propriete}[Connecteurs logiques pour l'expression écrite (niveau Terminale)]
\begin{center}
\begin{tabularx}{\linewidth}{|X|X|X|}
\hline
\rowcolor{popblueL} Fonction & Structure chinoise & Équivalent français \\
\hline
Temporel (antériorité) & \zh{……以后} / \zh{……之后} & après que… / dès que… / une fois que… \\
\hline
Cause & \zh{因为……所以……} / \zh{由于……} & parce que… donc… / en raison de… \\
\hline
Conséquence & \zh{……所以……} / \zh{结果……} & … donc… / … par conséquent… \\
\hline
Concession & \zh{虽然……但/但是……} & bien que… (mais)… — \textbf{un seul connecteur en français} \\
\hline
Condition & \zh{如果……就……} / \zh{要是……就……} & si… alors… / dans le cas où… \\
\hline
But & \zh{为了……} / \zh{以便……} & pour… / afin de… / de façon à… \\
\hline
Passif (agent) & \zh{被……} (subi) / \zh{由……} (formel, neutre/positif) & par… (passif) \\
\hline
Disposition (把) & \zh{把 + OBJ + V + 结果/方向} & (pas de mot) — antépose l'objet pour insister sur le résultat \\
\hline
Énumération / Ajout & \zh{不但……而且……} / \zh{另外} / \zh{此外} & non seulement… mais encore… / de plus / en outre \\
\hline
Conclusion & \zh{总之} / \zh{综上所述} & en résumé / en conclusion \\
\hline
\end{tabularx}
\end{center}
\end{propriete}

\begin{attention}[Pièges de traduction des connecteurs]
\begin{itemize}
\item \zh{虽然……但是……} $\rightarrow$ Français : \textbf{« Bien que… »} (virgule), \textbf{jamais} « Bien que… mais… ».
\item \zh{因为……所以……} $\rightarrow$ Français : \textbf{« Parce que…, … »} ou \textbf{« Comme…, … »}, \textbf{jamais} « Parce que… donc… »** (redondant).
\item \zh{把} et \zh{被}/\zh{由} ne se traduisent \textbf{jamais} par un mot ; ils dictent la structure de la phrase française (ordre des mots, voix passive/active).
\item \zh{以后} après un verbe/événement = « après que… » (passé) ; en tête de phrase = « plus tard / à l'avenir ».
\end{itemize}
\end{attention}

\courssub{Version : chinois → français}

\begin{methode}[Méthode de la version en quatre étapes (rappel)]
\begin{enumerate}
\item \textbf{Lire le texte entier} sans écrire, pour saisir la situation, les acteurs, la chronologie.
\item \textbf{Repérer le squelette} de chaque phrase : Sujet — Verbe — Compléments. Identifier les mots-outils (\zh{了, 过, 把, 被, 由, 的, 得, 地}).
\item \textbf{Traduire le sens par blocs} (sujet + verbe, puis compléments), en cherchant l'équivalent français \emph{naturel}, pas le calque mot à mot.
\item \textbf{Rédiger et relire à voix haute} la version française. Si ça sonne « chinois », reformulez.
\end{enumerate}
\end{methode}

\begin{popfigure}
\begin{tikzpicture}[node distance=0.7cm,
  etape/.style={rectangle, rounded corners=4pt, draw=popblue, fill=popblueL, align=center, text width=3.2cm, minimum height=1.2cm, font=\small},
  fleche/.style={-{Stealth[length=3mm]}, thick, poporange}]
\node[etape, align=center] (a){Texte chinois\\ \zh{爷爷去世以后，房子由父亲继承。}};
\node[etape, right=1.2cm of a, fill=poporangeL, draw=poporange, align=center] (b){Sens global\\ « Grand-père mort après, maison par père hériter »};
\node[etape, right=1.2cm of b, fill=popgreenL, draw=popgreen, align=center] (c){Français correct\\ « Après le décès du grand-père, la maison a été héritée par le père. »};
\draw[fleche] (a) -- node[above, font=\footnotesize, popdark]{1. Comprendre} (b);
\draw[fleche] (b) -- node[above, font=\footnotesize, popdark]{2. Reformuler} (c);
\end{tikzpicture}
\legende{Processus mental de la version : du chinois au sens, du sens au français.}
\end{popfigure}

\courssub{Phrases modèles commentées (Version)}

\textbf{Phrase 1.}
\zh{奶奶去世以后，我们继承了她的小店。}
\emph{Nǎinai qùshì yǐhòu, wǒmen jìchéngle tā de xiǎodiàn.}
« Après le décès de grand-mère, nous avons hérité de sa petite boutique. »
\begin{itemize}
\item \zh{奶奶} : grand-mère paternelle. \zh{去世} : décéder (soutenu). \zh{以后} placé \textbf{après} l'événement = « après que… ».
\item \zh{继承了} : \zh{了} = accompli $\rightarrow$ passé composé.
\item \zh{她的小店} : possessif $\rightarrow$ « sa boutique » (pas « la boutique d'elle »).
\end{itemize}

\textbf{Phrase 2.}
\zh{叔叔没有孩子，所以他把财产留给了我。}
\emph{Shūshu méiyǒu háizi, suǒyǐ tā bǎ cáichǎn liú gěi le wǒ.}
« Mon oncle n'a pas d'enfant, donc il m'a laissé ses biens. »
\begin{itemize}
\item \zh{叔叔} : oncle (frère cadet du père). Français : « mon oncle ».
\item \zh{把财产留给了我} : structure \zh{把} (objet antéposé). \textbf{Ne pas traduire \zh{把}}. Verbe composé \zh{留给} (laisser à qqn). \zh{了} = accompli.
\end{itemize}

\textbf{Phrase 3.}
\zh{虽然遗产不多，但是家人没有吵架。}
\emph{Suīrán yíchǎn bù duō, dànshì jiārén méiyǒu chǎojià.}
« Bien que l'héritage soit modeste, la famille ne s'est pas disputée. »
\begin{itemize}
\item \zh{虽然……但是……} : concession. Français : \textbf{un seul mot} « Bien que » + indicatif (ou subjonctif si subjectif). « Mais » est optionnel en français (souvent omis à l'écrit soigné).
\item \zh{不多} : « pas beaucoup » $\rightarrow$ « modeste / peu important ».
\item \zh{吵架} : verbe intransitif « se disputer ».
\end{itemize}

\begin{exemplebox}[Voix passive formelle avec \zh{由}]
\zh{土地由我的妈妈继承。}
\emph{Tǔdì yóu wǒ de māma jìchéng.}
« La terre est héritée par ma mère. » $\rightarrow$ Plus naturel : « C'est ma mère qui hérite de la terre. » (Active, mise en relief).
\zh{由} (yóu) introduit l'agent dans un contexte écrit, administratif, juridique. \zh{被} (bèi) marque souvent une nuance subie/négative.
\end{exemplebox}

\courssub{Mini-texte de version guidée}

\begin{textebox}[L'héritage de mon grand-père maternel]
我的外公留下了一些土地。\\
\emph{Wǒ de wàigōng liúxiàle yìxiē tǔdì.}\\[4pt]
按照他的遗嘱，土地由我的妈妈和阿姨继承。\\
\emph{Ànzhào tā de yízhǔ, tǔdì yóu wǒ de māma hé āyí jìchéng.}\\[4pt]
她们决定一起种咖啡。\\
\emph{Tāmen juédìng yìqǐ zhòng kāfēi.}
\end{textebox}

\begin{enumerate}
\item \textbf{Phrase 1.} Sujet \zh{我的外公} (grand-père maternel). Verbe \zh{留下} + \zh{了} (accompli). Objet \zh{一些土地} (quelques terrains). $\rightarrow$ « Mon grand-père maternel a laissé quelques terrains. »
\item \textbf{Phrase 2.} \zh{按照} (conformément à) + \zh{遗嘱} (testament). Passif \zh{由} + agents (\zh{妈妈和阿姨}) + verbe \zh{继承}. $\rightarrow$ « Conformément à son testament, les terrains reviennent à ma mère et à ma tante. » (Actif naturel) ou « … sont hérités par… »
\item \textbf{Phrase 3.} \zh{她们} (elles — radical \zh{女} visible : femmes). \zh{决定} + \zh{一起} + \zh{种咖啡}. $\rightarrow$ « Elles ont décidé de cultiver du café ensemble. »
\end{enumerate}

\courssub{Les pièges de la version (Focus Terminale)}

\begin{attention}[Erreurs fréquentes des francophones — Version]
\begin{itemize}
\item \textbf{Calque de \zh{把}} : Écrire « Il \emph{a pris} les biens et me les a laissés ». \zh{把} n'est pas un verbe. Traduction : « Il m'a laissé les biens. »
\item \textbf{Calque de \zh{虽然……但是……}} : Écrire « Bien que l'héritage soit peu, \emph{mais} la famille… ». \textbf{Interdit} en français standard.
\item \textbf{Confusion grands-parents} : \zh{外公/外婆} (maternel) vs \zh{爷爷/奶奶} (paternel). Préciser « maternel/paternel » si pertinent.
\item \textbf{\zh{阿姨} $\neq$ grand-mère} : \zh{阿姨} = tante (sœur de la mère) ou terme poli pour femme d'âge mûr. Jamais « grand-mère ».
\item \textbf{Temps verbaux} : \zh{了} (accompli) $\rightarrow$ Passé composé / Imparfait selon aspect. \zh{过} (expérience) $\rightarrow$ « a déjà… ». Pas de \zh{了} + verbe d'état $\rightarrow$ Présent / Imparfait.
\item \textbf{Classificateurs oubliés} : \zh{一份遗嘱} (un testament), \zh{一些土地} (des terrains). Ne pas traduire « un testament » par « testament » seul si le chinois marque l'unité.
\end{itemize}
\end{attention}

\begin{exoresolu}[Version notée façon Bac]
Énoncé : Traduisez en français soigné :
\zh{爷爷去世以后，他的房子由我的父亲继承，但是父亲把房子留给了我的哥哥。}

\tcblower

\textbf{Solution détaillée.}

\textbf{Étape 1 — Lecture globale :} Succession : grand-père $\rightarrow$ père $\rightarrow$ frère aîné.

\textbf{Étape 2 — Squelette :}
\begin{itemize}
\item Proposition 1 : \zh{爷爷} (Suj) \zh{去世} (V) \zh{以后} (Temps).
\item Proposition 2 : \zh{他的房子} (Suj passif) \zh{由} \zh{我的父亲} (Agent) \zh{继承} (V).
\item Proposition 3 : \zh{但是} (Concession/opposition) \zh{父亲} (Suj) \zh{把房子留给了} (Structure \zh{把} + V composé + \zh{了}) \zh{我的哥哥} (Destinataire).
\end{itemize}

\textbf{Étape 3 — Traduction du sens :}
\begin{itemize}
\item \zh{爷爷去世以后} $\rightarrow$ « Après le décès de mon grand-père paternel »
\item \zh{他的房子由我的父亲继承} $\rightarrow$ « sa maison a été héritée par mon père » ou « c'est mon père qui a hérité de sa maison »
\item \zh{但是父亲把房子留给了我的哥哥} $\rightarrow$ « mais mon père a laissé la maison à mon frère aîné »
\end{itemize}

\textbf{Étape 4 — Rédaction finale :}
\begin{quote}
« Après le décès de mon grand-père paternel, c'est mon père qui a hérité de sa maison, mais il l'a laissée à mon frère aîné. »
\end{quote}

\textbf{Erreurs à éviter (zéro point si présentes) :}
\begin{itemize}
\item « Après mon grand-père est mort… » (anglicisme/calque structure \zh{以后}).
\item « Son maison a été héritée par mon père » (accord \emph{sa/son}).
\item « Mais mon père \emph{a pris} la maison et l'a laissée… » (calque \zh{把} = prendre).
\item « Mon frère » au lieu de « \emph{mon frère aîné} » (\zh{哥哥} vs \zh{弟弟}).
\end{itemize}
\end{exoresolu}

\begin{exercice}[Entraînement Version — À vous de traduire]
Traduisez en français correct et naturel :
\begin{enumerate}
\item \zh{外婆去世以后，我们继承了她的小店。}
\item \zh{虽然遗产不多，但是家人都很高兴。}
\item \zh{按照爷爷的遗嘱，土地由我的叔叔继承。}
\item \zh{父亲决定把咖啡园留给哥哥，因为哥哥喜欢种地。}
\item \zh{由于没有遗嘱，家人为遗产吵了很久的架。}
\end{enumerate}
\end{exercice}

\begin{corrige}[Exercice « Entraînement Version »]
\begin{enumerate}
\item « Après le décès de ma grand-mère maternelle, nous avons hérité de sa petite boutique. »
\item « Bien que l'héritage soit modeste, toute la famille est très contente. » (Un seul connecteur ; \zh{都} = tous/toute).
\item « Conformément au testament de grand-père (paternel), la terre a été héritée par mon oncle » ou « …, c'est mon oncle qui a hérité de la terre. »
\item « Mon père a décidé de laisser la plantation de café à mon frère aîné, car celui-ci aime cultiver la terre. » (\zh{把} $\rightarrow$ actif ; \zh{因为} $\rightarrow$ car/parce que ; \zh{种地} = cultiver la terre).
\item « Faute de testament / En l'absence de testament, la famille s'est disputée longtemps pour l'héritage. » (\zh{由于} = en raison de ; \zh{为……吵架} = se disputer pour… ; \zh{很久} = longtemps).
\end{enumerate}
\end{corrige}

\courssub{Thème : français → chinois}

\begin{methode}[Méthode du thème en quatre étapes]
\begin{enumerate}
\item \textbf{Analyser la phrase française} : identifier le sujet, le temps, la voix (active/passive), les compléments, la logique (cause, but, concession).
\item \textbf{Choisir la structure chinoise adaptée} : ordre SOV/Sujet-Temps-Lieu-Verbe-Objet ; \zh{把} si action sur objet précis avec résultat ; \zh{被}/\zh{由} si passif ; place des compléments (Temps > Lieu > Manière).
\item \textbf{Assembler le vocabulaire et la grammaire} : mots-outils (\zh{了, 的, 得, 地}), classificateurs, connecteurs.
\item \textbf{Vérifier et lire en pinyin} : tons, ordre des mots, accord \zh{的/得/地}, ponctuation chinoise.
\end{enumerate}
\end{methode}

\begin{propriete}[Correspondances clés Français $\rightarrow$ Chinois (Thème Héritage)]
\begin{center}
\begin{tabularx}{\linewidth}{|X|X|X|}
\hline
\rowcolor{popblueL} Français & Chinois (Structure) & Exemple \\
\hline
Après que… (événement passé) & \zh{Subject + Verbe + 以后 / 之后} & \zh{爷爷去世以后} \\
\hline
Hériter (action accomplie) & \zh{Subject + 继承 + 了 + Objet} & \zh{我们继承了房子} \\
\hline
Laisser / léguer (à qqn) & \zh{Subject + 把 + Objet + 留给 + Personne + 了} & \zh{他把房子留给了我} \\
\hline
Être hérité par (passif formel) & \zh{Objet + 由 + Agent + 继承} & \zh{房子由我继承} \\
\hline
Conformément à / Selon & \zh{按照 + Nom + …} & \zh{按照遗嘱} \\
\hline
Bien que… (concession) & \zh{虽然 + Proposition, 但/但是 + Proposition} & \zh{虽然不多，但没吵架} \\
\hline
Décider de faire & \zh{Subject + 决定 + Verbe} & \zh{决定一起种咖啡} \\
\hline
Se disputer pour… & \zh{为 + Objet + 吵架} & \zh{为遗产吵架} \\
\hline
\end{tabularx}
\end{center}
\end{propriete}

\begin{exemplebox}[Thème guidé : phrases types]
\begin{enumerate}
\item « Après le décès de mon grand-père, mon père a hérité de la maison. »
$\rightarrow$ \zh{爷爷去世以后，我父亲继承了房子。} (Active, simple)
\item « La maison a été laissée à mon frère aîné par mon père. »
$\rightarrow$ \zh{房子由我父亲留给了哥哥。} (Passif \zh{由} + \zh{留给}) OU \zh{我父亲把房子留给了哥哥。} (Active \zh{把} — plus naturel à l'oral).
\item « Bien que l'héritage soit petit, nous sommes contents. »
$\rightarrow$ \zh{虽然遗产不多，但我们很高兴。} (\zh{但} suffit, pas de \zh{但是} obligatoire).
\item « Conformément au testament, la terre revient à ma tante. »
$\rightarrow$ \zh{按照遗嘱，土地由阿姨继承。}
\end{enumerate}
\end{exemplebox}

\begin{exercice}[Entraînement Thème — Traduisez en chinois (caractères + pinyin)]
\begin{enumerate}
\item Après le décès de ma grand-mère maternelle, nous avons hérité de sa petite boutique.
\item Bien que l'héritage ne soit pas grand, la famille ne s'est pas disputée.
\item Conformément au testament de grand-père, la terre a été héritée par mon oncle.
\item Mon père a décidé de laisser la plantation de café à mon frère aîné.
\item Faute de testament, ils se sont disputés longtemps pour l'héritage.
\end{enumerate}
\end{exercice}

\begin{corrige}[Exercice « Entraînement Thème »]
\begin{enumerate}
\item \zh{外婆去世以后，我们继承了她的小店。} (Wàipó qùshì yǐhòu, wǒmen jìchéngle tā de xiǎodiàn.)
\item \zh{虽然遗产不大，但家人没有吵架。} (Suīrán yíchǎn bù dà, dàn jiārén méiyǒu chǎojià.) — \zh{不大} ou \zh{不多} acceptés.
\item \zh{按照爷爷的遗嘱，土地由我的叔叔继承。} (Ànzhào yéye de yízhǔ, tǔdì yóu wǒ de shūshu jìchéng.)
\item \zh{我父亲决定把咖啡园留给了哥哥。} (Wǒ fùqin juédìng bǎ kāfēiyuán liú gěile gēge.) — \zh{咖啡园} (plantation de café) ; \zh{留给了} (laissé à, accompli).
\item \zh{由于没有遗嘱，他们为遗产吵了很久的架。} (Yóuyú méiyǒu yízhǔ, tāmen wèi yíchǎn chǎole hěnjiǔ de jià.) — \zh{由于} (cause formelle) ; \zh{吵了……架} (verbe + durée + objet verbal).
\end{enumerate}
\end{corrige}

\courssub{Grille d'évaluation et entraînement Bac}

\begin{aretenir}[Grille d'évaluation Type Bac — Expression écrite / Traduction (20 pts)]
\begin{center}
\begin{tabularx}{\linewidth}{|X|X|X|}
\hline
\rowcolor{popblueL} Critère & Indicateurs (Version / Thème) & Points \\
\hline
\textbf{Fidélité du sens} (Version) / \textbf{Respect du sens source} (Thème) & Aucune omission, aucun contresens, vocabulaire précis (ex: \zh{外公} vs \zh{爷爷}), temps verbaux justes. & 6 \\
\hline
\textbf{Qualité du français} (Version) / \textbf{Correctness du chinois} (Thème) & Grammaire, syntaxe, orthographe, ponctuation (chinoise : 。，？), ordre des mots (T-L-V-O), mots-outils (\zh{了, 的, 得, 地, 把, 被, 由}). & 6 \\
\hline
\textbf{Style et naturel} (Version) / \textbf{Idiomaticité} (Thème) & Phrases fluides, connecteurs adaptés (\zh{虽然…但…}, \zh{按照…}, \zh{由…}), pas de calque, registres appropriés. & 4 \\
\hline
\textbf{Richesse lexicale \& structures} & Utilisation du vocabulaire de la leçon (\zh{遗产, 遗嘱, 继承, 留下, 把, 由, 按照}), variété des structures (passif, \zh{把}, concession). & 4 \\
\hline
\end{tabularx}
\end{center}
\end{aretenir}

\begin{savaistu}[Héritage : regards croisés Cameroun — Chine]
\begin{itemize}
\item \textbf{Cameroun (droit coutumier / moderne) :} Coexistence du Code civil (héritage légal : conjoint + enfants) et des coutumes (souvent patrilinéaire : transmission aux fils, exclusion fréquente des filles/veuves selon ethnies). Le testament (\zh{遗嘱}) gagne du terrain en zone urbaine (Yaoundé, Douala).
\item \textbf{Chine (Code civil 2021) :} Égalité successorale hommes/femmes (Art. 1127). Héritiers de premier ordre : conjoint, enfants, parents. Testament (\zh{遗嘱}) valide s'il est écrit, signé, daté ; testament oral toléré en urgence. \zh{遗产管理人} (exécuteur testamentaire) prévu par la loi.
\item \textbf{Point culturel :} En Chine, \zh{分家} (fēnjiā, partage du patrimoine familial) est un moment sensible ; la piété filiale (\zh{孝顺} xiàoshun) dicte souvent le soin aux aînés avant l'héritage. Au Cameroun, le \emph{deuil} et la \emph{levée de deuil} rythment la transmission.
\end{itemize}
\end{savaistu}

\begin{experience}[Débat oral guidé : « Héritage et égalité »]
\textbf{Sujet} : « Filles et fils doivent-ils hériter à parts égales ? » \\
\textbf{Consigne} : En groupes de 4. Préparez 3 arguments en chinois (vocabulaire fourni : \zh{平等} égalité, \zh{传统} tradition, \zh{法律} loi, \zh{照顾} prendre soin de, \zh{赡养} subvenir aux besoins des parents âgés). Débat 10 min. Un rapporteur synthétise en français pour la classe. Évaluation : prononciation, vocabulaire, interaction.
\end{experience}

\coursec{Grille d'évaluation et entraînement Bac}

Dans la section précédente, vous avez appris à traduire du chinois vers le français (version) des phrases sur le thème de l'héritage. Vous savez maintenant lire, comprendre et traduire. Il reste une dernière étape, la plus importante pour l'examen : \cle{produire} vous-même un texte écrit en chinois, et savoir exactement comment il sera noté. Cette dernière section vous donne la grille d'évaluation officielle, les critères de réussite, une méthode de gestion du temps, puis deux sujets d'entraînement complets de type Bac, avec corrigés et correction collective d'une copie anonyme.

\begin{definition}[Glossaire clé de l'héritage — rappel pour l'examen]
\begin{center}
\begin{tabularx}{\linewidth}{|X|X|X|X|}
\hline
\rowcolor{popblueL} \textbf{Caractères} & \textbf{Pinyin} & \textbf{Nature} & \textbf{Traduction} \\
\hline
继承 & jìchéng & verbe & hériter \\
\hline
遗产 & yíchǎn & nom & héritage, succession \\
\hline
遗嘱 & yízhǔ & nom & testament \\
\hline
分配 & fēnpèi & verbe & répartir, partager \\
\hline
去世 & qùshì & verbe & décéder, mourir (terme poli) \\
\hline
团结 & tuánjié & adj./verbe & uni, solidaire / s'unir \\
\hline
法律 & fǎlǜ & nom & loi \\
\hline
留下 & liúxià & verbe & laisser (en héritage) \\
\hline
平分 & píngfēn & verbe & partager équitablement \\
\hline
叔叔 & shūshu & nom & oncle (frère cadet du père) \\
\hline
\end{tabularx}
\end{center}
\end{definition}

\courssub{Comment le correcteur note votre production écrite}

Observez d'abord la grille ci-dessous : c'est l'outil que le correcteur a sous les yeux quand il lit votre copie. Elle comporte \cle{cinq critères}, chacun noté sur 4 points, pour un total de 20 points.

\begin{popfigure}
\begin{center}
\begin{tabularx}{\linewidth}{|X|c|X|}
\hline
\rowcolor{popblue} \textbf{\color{white}Critère} & \textbf{\color{white}Points} & \textbf{\color{white}Auto-contrôle de l'élève} \\
\hline
\rowcolor{popblueL} 1. Exactitude des caractères & \textbf{/4} & Aucun caractère inventé, traits complets, écriture lisible \\
\hline
\rowcolor{poporangeL} 2. Grammaire et ordre des mots & \textbf{/4} & Sujet + temps + verbe ; place correcte de 由、把、了 \\
\hline
\rowcolor{popgreenL} 3. Vocabulaire du thème & \textbf{/4} & Au moins 5 mots du glossaire : 继承、遗产、遗嘱、分配… \\
\hline
\rowcolor{poppurpleL} 4. Connecteurs et cohérence & \textbf{/4} & Au moins 3 connecteurs différents : 首先、然后、虽然…但是… \\
\hline
\rowcolor{popgoldL} 5. Consigne et longueur & \textbf{/4} & Sujet respecté, 6 à 10 phrases, réponse complète \\
\hline
\end{tabularx}
\end{center}
\legende{Grille d'évaluation d'une production écrite sur le thème de l'héritage (20 points)}
\end{popfigure}

\begin{definition}[Les cinq critères expliqués]
\begin{enumerate}
\item \textbf{Exactitude des caractères (4 pts)} : chaque caractère doit être complet, avec tous ses traits, et lisible. Un caractère mal formé ou inventé fait perdre des points, même si le sens se devine.
\item \textbf{Grammaire et ordre des mots (4 pts)} : le chinois exige l'ordre Sujet + (complément de temps) + verbe. Le complément de temps se place \emph{avant} le verbe, jamais à la fin comme en français.
\item \textbf{Vocabulaire du thème (4 pts)} : le correcteur compte les mots du glossaire réellement employés : \zh{继承} (jìchéng, hériter), \zh{遗产} (yíchǎn, héritage), \zh{遗嘱} (yízhǔ, testament), \zh{分配} (fēnpèi, partager), \zh{去世} (qùshì, décéder).
\item \textbf{Connecteurs et cohérence (4 pts)} : le texte doit avancer logiquement grâce à des connecteurs variés.
\item \textbf{Respect de la consigne et longueur (4 pts)} : hors sujet ou texte trop court = points perdus, quelle que soit la qualité de la langue.
\end{enumerate}
\end{definition}

\begin{aretenir}[Critères de réussite à mémoriser]
Une copie réussie sur ce thème contient :
\begin{itemize}
\item \textbf{au moins 6 phrases} complètes (sujet + verbe) ;
\item \textbf{au moins 3 connecteurs différents} : \zh{首先} (d'abord), \zh{然后} (ensuite), \zh{最后} (enfin), \zh{虽然…但是…} (bien que… mais…) ;
\item \textbf{au moins 5 mots du glossaire} du thème de l'héritage ;
\item \textbf{zéro confusion} entre caractères proches : \zh{遗} (yí, laisser) et \zh{送} (sòng, offrir) ; \zh{继} (jì, continuer) et \zh{断} (duàn, couper) ; \zh{产} (chǎn, produire) et \zh{立} (lì, établir).
\end{itemize}
\end{aretenir}

\begin{attention}[La qualité graphique compte !]
En expression écrite chinoise, une copie avec des \cle{caractères illisibles ou inventés} perd des points \emph{même si le sens est correct}. Le correcteur ne peut pas deviner : un trait manquant transforme \zh{大} (dà, grand) en \zh{人} (rén, personne), un trait de trop change \zh{土} (tǔ, terre) en \zh{王} (wáng, roi). Écrivez lentement, carrément, en laissant de l'espace entre les caractères. Mieux vaut 6 phrases aux caractères nets que 10 phrases brouillonnes.
\end{attention}

\courssub{Auto-évaluation : la check-list avant de rendre la copie}

Avant de rendre votre copie, cochez mentalement chaque point de cette liste. C'est votre dernier filtre contre les pertes de points évitables.

\begin{methode}[Gérer son temps à l'examen (25 minutes)]
\begin{enumerate}
\item \textbf{5 minutes de plan} : lisez le sujet deux fois, soulignez les mots-clés, notez au brouillon 5 mots du glossaire et 3 connecteurs que vous utiliserez, puis numérotez vos idées (6 à 8 phrases).
\item \textbf{15 minutes de rédaction} : écrivez phrase par phrase, sans ratures, en appliquant l'ordre Sujet + temps + verbe. Ne cherchez jamais un caractère plus de 30 secondes : remplacez le mot par un synonyme simple.
\item \textbf{5 minutes de relecture en trois passages} :
\begin{enumerate}
\item passage 1 : \textbf{caractères} — chaque caractère est-il complet et lisible ?
\item passage 2 : \textbf{grammaire} — le temps est-il avant le verbe ? \zh{了} est-il bien placé ? \zh{由} introduit-il bien l'agent ?
\item passage 3 : \textbf{connecteurs} — y a-t-il au moins 3 connecteurs différents et une conclusion ?
\end{enumerate}
\end{enumerate}
\end{methode}

\courssub{Sujet d'entraînement type Bac n°1 : rédaction}

\begin{exercice}[Sujet type Bac — Rédaction (20 points, 25 minutes)]
Racontez en \textbf{6 à 10 phrases} comment votre famille a partagé l'héritage d'un grand-parent. Vous utiliserez au moins 5 mots du glossaire (\zh{去世、遗产、继承、分配、遗嘱、团结}) et au moins 3 connecteurs différents.
\end{exercice}

\begin{exemplebox}[Sujet corrigé — modèle de réponse en 4 phrases (noyau minimal)]
\zh{我的奶奶前年去世了。她留下了一所房子。按照法律，房子由我的爸爸继承。虽然我们不富有，但是我们很团结。}\\[4pt]
\emph{Wǒ de nǎinai qiánnián qùshì le. Tā liúxià le yì suǒ fángzi. Ànzhào fǎlǜ, fángzi yóu wǒ de bàba jìchéng. Suīrán wǒmen bù fùyǒu, dànshì wǒmen hěn tuánjié.}\\[4pt]
« Ma grand-mère est décédée il y a deux ans. Elle a laissé une maison. Conformément à la loi, la maison a été héritée par mon père. Bien que nous ne soyons pas riches, nous sommes très unis. »\\[6pt]
Pour atteindre 6 à 10 phrases, on développe : \zh{首先，我们全家开了一个会。然后，按照奶奶的遗嘱，土地被分配给了我的叔叔。最后，我们把房子租出去，这笔钱用来付我们的学费。}\\[4pt]
\emph{Shǒuxiān, wǒmen quánjiā kāi le yí gè huì. Ránhòu, ànzhào nǎinai de yízhǔ, tǔdì bèi fēnpèi gěi le wǒ de shūshu. Zuìhòu, wǒmen bǎ fángzi chūqù, zhè bǐ qián yònglái fù wǒmen de xuéfèi.}\\[4pt]
« D'abord, toute notre famille a tenu une réunion. Ensuite, conformément au testament de grand-mère, la terre a été répartie à mon oncle. Enfin, nous avons loué la maison et cet argent sert à payer nos frais de scolarité. »
\end{exemplebox}

\begin{savaistu}[Ce que les correcteurs valorisent]
Les correcteurs récompensent l'usage d'une \cle{structure un peu riche} même dans un texte simple : la voix passive avec \zh{由} (\zh{房子由我的爸爸继承} — « la maison est héritée par mon père »), la construction \zh{把} (\zh{我们把房子租出去} — « nous avons loué la maison »), ou la concession \zh{虽然…但是…}. Une seule de ces structures, correctement employée, distingue immédiatement votre copie. Attention toutefois : une structure riche \emph{fausse} coûte plus cher qu'une structure simple \emph{juste}. Ne tentez à l'écrit que ce que vous maîtrisez.
\end{savaistu}

\courssub{Sujet d'entraînement type Bac n°2 : thème et version}

\begin{exoresolu}[Thème : français → chinois (4 phrases)]
\textbf{Énoncé.} Traduisez en chinois : 1. Mon grand-père est mort l'année dernière. 2. Il a laissé un champ et une maison. 3. Selon son testament, le champ revient à mon oncle. 4. Notre famille est restée très unie.
\tcblower
\textbf{Solution détaillée.}
\begin{enumerate}
\item \zh{我的爷爷去年去世了。} \emph{Wǒ de yéye qùnián qùshì le.} — Le temps \zh{去年} se place \emph{avant} le verbe ; \zh{了} marque l'accomplissement.
\item \zh{他留下了一块地和一所房子。} \emph{Tā liúxià le yí kuài dì hé yì suǒ fángzi.} — Classificateurs : \zh{块} pour la terre, \zh{所} pour la maison.
\item \zh{按照他的遗嘱，地由我的叔叔继承。} \emph{Ànzhào tā de yízhǔ, dì yóu wǒ de shūshu jìchéng.} — « Selon » se dit \zh{按照}, placé en tête de phrase ; « revient à » se rend élégamment par la passive \zh{由…继承}.
\item \zh{我们的家庭一直很团结。} \emph{Wǒmen de jiātíng yìzhí hěn tuánjié.} — \zh{一直} (toujours, en continu) se place avant l'adjectif.
\end{enumerate}
\end{exoresolu}

\begin{exoresolu}[Version : chinois → français (mini-texte de 3 phrases)]
\textbf{Énoncé.} Traduisez en français : \zh{我的外婆去世后，留下了一些钱。按照法律，这笔钱由我的妈妈和她的姐姐平分。虽然钱不多，但是它帮助了我们全家。}
\tcblower
\textbf{Solution détaillée.} « Après le décès de ma grand-mère maternelle, elle a laissé un peu d'argent. Conformément à la loi, cette somme a été partagée en deux parts égales entre ma mère et sa sœur aînée. Bien que cette somme ne soit pas importante, elle a aidé toute notre famille. »\\[4pt]
Points de vigilance : \zh{外婆} = grand-mère \emph{maternelle} (ne pas traduire simplement « grand-mère » si le contexte exige la précision) ; \zh{平分} = partager en parts \emph{égales} ; \zh{虽然…但是…} se traduit par « bien que… » (le « mais » français ne se répète pas dans la même phrase).
\end{exoresolu}

\courssub{Correction collective d'une copie anonyme}

Voici une copie réelle d'élève (anonymisée), reproduite avec ses erreurs. Lisez-la, repérez les fautes, puis comparez avec la correction.

\begin{textebox}[Copie anonyme d'élève — à corriger]
\zh{我的奶奶去世了去年。她留下一所房子和钱。房子我的爸爸继承。我们虽然不富有，很团结。}\\[4pt]
\emph{(Transcription pinyin de l'élève) Wǒ de nǎinai qùshì le qùnián. Tā liúxià yì suǒ fángzi hé qián. Fángzi wǒ de bàba jìchéng. Wǒmen suīrán bù fùyǒu, hěn tuánjié.}\\[4pt]
Intention de traduction : « Ma grand-mère est morte l'année dernière. Elle a laissé une maison et de l'argent. Mon père a hérité de la maison. Bien que nous ne soyons pas riches, nous sommes unis. »
\end{textebox}

\begin{center}
\begin{tabularx}{\linewidth}{|X|X|X|}
\hline
\rowcolor{popblueL} \textbf{Erreur de l'élève} & \textbf{Type d'erreur} & \textbf{Correction} \\
\hline
\zh{去世了去年} & Place du temps : le complément de temps doit précéder le verbe & \zh{去年去世了} \\
\hline
\zh{留下一所房子} & \zh{了} manquant après le verbe (action accomplie) & \zh{留下了一所房子} \\
\hline
\zh{房子我的爸爸继承} & Ordre des mots : sujet et agent confondus ; manque \zh{由} pour la passive & \zh{房子由我的爸爸继承} \\
\hline
\zh{虽然不富有，很团结} & Connecteur incomplet : \zh{虽然} exige \zh{但是} en regard ; sujets répétés & \zh{虽然我们不富有，但是我们很团结} \\
\hline
\end{tabularx}
\end{center}

\begin{corrige}[Copie anonyme]
Texte corrigé : \zh{我的奶奶去年去世了。她留下了一所房子和一些钱。房子由我的爸爸继承。虽然我们不富有，但是我们很团结。}\\[4pt]
« Ma grand-mère est décédée l'année dernière. Elle a laissé une maison et un peu d'argent. La maison a été héritée par mon père. Bien que nous ne soyons pas riches, nous sommes très unis. »\\[4pt]
Bilan selon la grille : caractères 4/4 (écriture correcte), grammaire 1/4 (trois fautes d'ordre des mots), vocabulaire 3/4 (quatre mots du glossaire), connecteurs 1/4 (un seul connecteur, incomplet), consigne 2/4 (seulement 4 phrases). Total : 11/20. La même copie, avec les corrections ci-dessus et deux phrases de plus, atteint 17/20.
\end{corrige}

\begin{exercice}[Entraînement final — à faire seul en 25 minutes]
1. Rédaction : racontez en 6 à 10 phrases le partage de l'héritage d'un grand-parent dans votre famille (réelle ou imaginaire), avec 5 mots du glossaire et 3 connecteurs.\\
2. Thème : traduisez « Mon oncle a écrit un testament. Selon la loi, la maison revient à mes cousins. »\\
3. Auto-évaluez votre copie avec la grille avant de la rendre.
\end{exercice}

Vous disposez désormais de tous les outils pour réussir l'épreuve d'expression écrite sur le thème de l'héritage : un glossaire maîtrisé, des caractères écrits correctement, des connecteurs variés, une méthode de gestion du temps et une grille d'auto-évaluation. Entraînez-vous régulièrement : en chinois, c'est la main qui retient autant que la mémoire.

\newpage

\coursec{Activité d'intégration}

\coursec{Leçon 7 — Expression écrite : l'héritage}

\courssub{Objectifs de l'activité}
Cette activité d'intégration vise à mobiliser l'ensemble des compétences travaillées lors de la leçon 7 (vocabulaire de l'héritage, structures grammaticales \textbf{把} et \textbf{由}, connecteurs logiques, écriture des caractères clés) dans une tâche de production écrite complexe et authentique, calquée sur les attendus de l'épreuve du Baccalauréat (LV2).
Les objectifs spécifiques sont les suivants :
\begin{itemize}
    \item \textbf{Compétence communicative :} Produire un texte narratif cohérent de 8 à 10 phrases en chinois simplifié, intitulé \zh{爷爷留下的房子} (La maison laissée par grand-père), respectant les conventions de l'écrit (ponctuation pleine chasse, paragraphage).
    \item \textbf{Compétence linguistique :} Réemployer correctement le lexique thématique (biens, droit, relations familiales, sentiments), la structure \cle{把} pour focaliser sur la manipulation/disposition des objets, la structure \cle{由} pour exprimer la cause/l'origine ou la passivité formelle, et au minimum quatre connecteurs logiques (\cle{首先}、\cle{然后}、\cle{但是}、\cle{最后} ou équivalents).
    \item \textbf{Compétence graphique :} Écrire correctement les caractères du thème (ordre des traits, radicaux, structure), en portant une attention particulière aux caractères proches (ex: \zh{遗} vs \zh{贵}, \zh{产} vs \zh{生}).
    \item \textbf{Compétence métacognitive / Évaluation :} Utiliser une grille d'évaluation critériée pour corriger le texte d'un autre groupe (correction croisée), justifier les corrections apportées et présenter le meilleur texte à la classe.
    \item \textbf{Compétence socioculturelle :} Ancrer le récit dans le contexte camerounais (coutumes successorales, rôle du patriarche, types de biens : case, champ de cacao, boutique) tout en utilisant la langue cible.
\end{itemize}

\courssub{Situation}
\begin{textebox}[Contexte : La succession de grand-père Mbarga]
\zh{雅温得郊外的一个村庄里，姆巴尔加爷爷 (Mǔbā'ěrjiā Yéye) 于上个月去世，享年八十五岁。} \\
\zh{他一生勤劳，留下了一栋砖房 (砖房 zhuānfáng)、一片可可园 (可可园 kěkěyuán) 和市场上的一间小商店 (小商店 xiǎo shāngdiàn)。} \\
\zh{爷爷生前立有遗嘱 (遗嘱 yízhǔ)，但家里对分配方式 (分配方式 fēnpèi fāngshì) 还有不同意见。} \\
\zh{全家人 (全家人 quánjiārén) 在大舅舅 (大舅舅 dàjiùjiu) 的主持下开会商量。} \\
\zh{最终，大家决定既尊重法律 (法律 fǎlǜ) 又尊重爷爷的心愿 (心愿 xīnyuàn)，妥善处理遗产 (遗产 yíchǎn)，维护家庭和睦 (和睦 hémù)。}
\par
\textbf{Traduction :} Dans un village en périphérie de Yaoundé, Grand-père Mbarga est décédé le mois dernier à l'âge de 85 ans. Il a travaillé dur toute sa vie et a laissé une maison en briques, une plantation de cacao et une petite boutique au marché. Grand-père avait fait un testament de son vivant, mais la famille a des avis divergents sur le mode de répartition. Toute la famille se réunit sous la présidence du grand oncle maternel pour discuter. Finalement, tout le monde décide de respecter à la fois la loi et les souhaits de grand-père, de gérer correctement l'héritage et de préserver l'harmonie familiale.
\end{textebox}

\courssub{Consignes}
\begin{enumerate}
    \item \textbf{Analyse de la situation et collecte lexicale (15 min) :} 
    En groupe de trois, lisez le texte de la situation. Identifiez et listez dans votre cahier : 
    \begin{itemize}
        \item Les 5 à 7 mots-clés du vocabulaire de l'héritage présents dans le texte (ex: \zh{遗产}、\zh{遗嘱}、\zh{分配}、\zh{继承}).
        \item Les personnages impliqués et leurs titres de parenté en chinois (\zh{大舅舅}、\zh{堂哥}、\zh{姑姑} etc.).
        \item Les trois biens laissés par grand-père.
    \end{itemize}
    Vérifiez l'écriture exacte des caractères (radicaux, traits) à l'aide du tableau de la section \textit{Ressources à mobiliser}.

    \item \textbf{Planification de la narration (15 min) :} 
    Construisez le squelette de votre récit \zh{爷爷留下的房子} en 4 étapes logiques (utilisez des mots de liaison en français d'abord) :
    \begin{enumerate}
        \item \textbf{Introduction :} Présentation du décès, du personnage (Grand-père Mbarga) et du lieu (village près de Yaoundé).
        \item \textbf{Énumération des biens :} Description des trois biens (maison, champ, boutique) avec adjectifs de qualité/taille.
        \item \textbf{Processus de répartition :} Conflit ou discussion, intervention de l'oncle, référence à la loi (\zh{法律}) et au testament (\zh{遗嘱}), décision finale. \textbf{Ici, vous devez insérer la structure \cle{把} (ex: \zh{把房子分给了大儿子}) et la structure \cle{由} (ex: \zh{分配由大舅舅负责} / \zh{由遗嘱决定}).}
        \item \textbf{Conclusion sentimentale :} Sentiment de la famille (apaisement, unité, souvenir).
    \end{enumerate}
    Sélectionnez \textbf{au moins 4 connecteurs} dans la liste fournie en Ressources pour baliser ces étapes.

    \item \textbf{Rédaction collaborative du premier jet (25 min) :} 
    Rédigez le texte complet en chinois (caractères + pinyin au brouillon si nécessaire). 
    \textbf{Contraintes obligatoires :}
    \begin{itemize}
        \item 8 à 10 phrases complètes (ponctuation chinoise 。).
        \item Utilisation effective de \textbf{把} (au moins 1 fois) et \textbf{由} (au moins 1 fois).
        \item Utilisation effective de \textbf{4 connecteurs} minimum (ex: \zh{首先}、\cle{接着}、\cle{不过}、\cle{最终}).
        \item Respect de l'ordre des mots chinois (Sujet + Temps + Lieu + Verbe + Objet / Compléments).
        \item Vocabulaire du thème (héritage, famille, biens) prioritaire.
    \end{itemize}
    Un élève écrit au tableau/feuille finale, un autre vérifie la grammaire (structures imposées), le troisième vérifie les caractères (écriture, radicaux) et le pinyin.

    \item \textbf{Correction croisée avec grille d'évaluation (20 min) :} 
    Échangez votre production avec un autre groupe. Chaque groupe corrige le texte reçu en utilisant la \textbf{Grille d'évaluation Bac} (fournie en Ressources). 
    Annotez le texte : 
    \begin{itemize}
        \item \textcolor{popgreen}{Vert} = Correct / Bien formulé.
        \item \textcolor{poporange}{Orange} = Maladresse de style, répétition, connecteur faible.
        \item \textcolor{poppink}{Rouge} = Erreur grammaticale (ordre des mots, mauvaise structure \cle{把}/\cle{由}), erreur de caractère (caractère faux, trait manquant), omission de contrainte.
    \end{itemize}
    Remplissez la grille ensemble et notez 3 points forts et 3 axes de progression.

    \item \textbf{Réécriture finale et présentation orale (15 min) :} 
    Récupérez votre texte original et les annotations. Discutez des corrections proposées, apportez les modifications définitives. 
    Désignez un lecteur par groupe. Le professeur sélectionne le texte le plus abouti (respect des contraintes, richesse lexicale, justesse grammaticale, cohérence narrative) pour une lecture à haute voix devant la classe et affichage au « Mur de la Réussite ».
\end{enumerate}

\courssub{Ressources à mobiliser}
\begin{definition}[Lexique thématique : L'héritage et la famille (HSK 3-4)]
\begin{center}
\begin{tabularx}{\linewidth}{|X|X|X|X|}
\hline \rowcolor{popblueL} 
\textbf{Caractères} & \textbf{Pinyin} & \textbf{Nature} & \textbf{Traduction} \\ \hline
遗产 & yíchǎn & n. & héritage, succession \\ \hline
遗嘱 & yízhǔ & n. & testament \\ \hline
继承 & jìchéng & v. & hériter, succéder à \\ \hline
分配 & fēnpèi & v./n. & répartir, distribution / répartition \\ \hline
律师 & lǜshī & n. & avocat \\ \hline
公证 & gōngzhèng & v./n. & notarié, authentifier / notariat \\ \hline
房产 & fángchǎn & n. & bien immobilier, propriété \\ \hline
田地 & tiándì & n. & terres agricoles, champs \\ \hline
店铺 & diànpù & n. & boutique, magasin, commerce \\ \hline
和睦 & hémù & adj. & harmonieux, en bons termes (famille) \\ \hline
纠纷 & jiūfēn & n. & litige, dispute \\ \hline
心愿 & xīnyuàn & n. & vœu, souhait profond \\ \hline
长子 & zhǎngzǐ & n. & fils aîné \\ \hline
幺妹 & yāomèi & n. & cadette (familier, dernière fille, régional) \\ \hline
主持 & zhǔchí & v. & présider, animer (réunion) \\ \hline
妥善 & tuǒshàn & adv./adj. & correctement, convenablement / approprié \\ \hline
\end{tabularx}
\end{center}
\end{definition}

\begin{propriete}[Structure \cle{把} (bǎ) : Disposition / Manipulation d'un objet connu]
\textbf{Schéma :} \textbf{Sujet + 把 + Objet (connu/spécifique) + Verbe + Complément de résultat / direction / aspect}
\par
\textbf{Emploi :} Met l'accent sur ce qui \textbf{arrive à l'objet} (action aboutie, changement d'état, déplacement). L'objet doit être défini (connu de l'interlocuteur, défini par un démonstratif, un possessif, le contexte). Le verbe ne peut pas être seul ; il nécessite un complément (particule \zh{了}, \zh{完}, \zh{好}, \zh{给某人}, \zh{到...}, \zh{在...}).
\par
\begin{center}
\begin{tikzpicture}[node distance=0.7cm, every node/.style={font=\small, align=center}]
\node (S) {Sujet};
\node[right=of S] (BA) {\zh{把}};
\node[right=of BA] (Obj) {Objet \textbf{connu}};
\node[right=of Obj] (V) {Verbe};
\node[right=of V, align=center] (Comp){Complément \\(résultat / direction / \zh{了}/\zh{好}...)};
\draw[->, popblue, thick] (S) -- (BA);
\draw[->, popblue, thick] (BA) -- (Obj);
\draw[->, popblue, thick] (Obj) -- (V);
\draw[->, popblue, thick] (V) -- (Comp);
\end{tikzpicture}
\end{center}
\textbf{Exemples du thème :}
\begin{itemize}
    \item \zh{大舅舅把遗嘱读给大家听。} (Dàjiùjiu \textbf{bǎ} yízhǔ dú gěi dàjiā tīng.) $\rightarrow$ L'oncle a lu le testament à tout le monde (action aboutie vers un destinataire).
    \item \zh{他们把那片可可园分给了二叔。} (Tāmen \textbf{bǎ} nà piàn kěkěyuán fēn gěi le èrshū.) $\rightarrow$ Ils ont attribué cette plantation de cacao au deuxième oncle (disposition, transfert de possession). \textbf{Attention :} \zh{分给} (fēn gěi) = partager à / attribuer à.
    \item \zh{我们把爷爷的房子修好了。} (Wǒmen \textbf{bǎ} yéye de fángzi xiū hǎo le.) $\rightarrow$ Nous avons réparé la maison de grand-père (changement d'état : cassé $\rightarrow$ réparé).
\end{itemize}
\par
\end{propriete}

\begin{attention}[Pièges francophones avec \cle{把}]
\begin{itemize}
    \item \textbf{Objet indéfini :} *\zh{我把一本书看了} (Faux). $\rightarrow$ \zh{我看了一本书} (Correct). \cle{把} exige un objet spécifique (défini par le contexte, un démonstratif, un possessif).
    \item \textbf{Verbe nu :} *\zh{他把门关} (Faux). $\rightarrow$ \zh{他把门关上了 / 关了} (Correct). Il faut un complément (résultat \zh{上}, aspect \zh{了}).
    \item \textbf{Négation :} \zh{没 (méi)} devant \zh{把} (\zh{他没把门关上}). Pas \zh{不}.
    \item \textbf{Verbes d'état/sentiment :} *\zh{我把汉语喜欢} (Faux). $\rightarrow$ \zh{我喜欢汉语}. \cle{把} s'utilise avec des verbes d'action.
\end{itemize}
\end{attention}


\begin{propriete}[Structure \cle{由} (yóu) : Origine, Cause, Passif formel, « C'est ... qui »]
\textbf{Schéma :} \textbf{Sujet (résultat/événement) + 由 + Agent / Cause / Base + Verbe / Adjectif}
\par
\textbf{Emploi :} Registre plus formel, écrit, administratif ou narratif. Trois usages principaux :
\begin{enumerate}
    \item \textbf{Agent du passif (formel) :} \zh{这次会议由大舅舅主持。} (Zhè cì huìyì \textbf{yóu} dàjiùjiu zhǔchí.) = Cette réunion est présidée par l'oncle. (Remplace \zh{被} dans le style soutenu).
    \item \textbf{Cause / Raison (souvent avec 由于 yóuyú) :} \zh{家庭和睦由相互尊重得来。} (Jiātíng hémù \textbf{yóu} xiānghù zūnzhòng dé lái.) = L'harmonie familiale vient du respect mutuel.
    \item \textbf{Base / Selon (règle, loi, testament) :} \zh{分配方案由法律决定。} (Fēnpèi fāng'àn \textbf{yóu} fǎlǜ juédìng.) = Le plan de répartition est décidé par la loi / selon la loi. \zh{由遗嘱分配} (Réparti selon le testament).
\end{enumerate}
\par
\begin{center}
\begin{tikzpicture}[node distance=0.7cm, every node/.style={font=\small, align=center}]
\node[align=center] (S){Sujet \\(Résultat / Événement)};
\node[right=of S] (YOU) {\zh{由}};
\node[right=of YOU] (Ag) {Agent / Cause / Base};
\node[right=of Ag] (V) {Verbe / Adjectif};
\draw[->, poppurple, thick] (S) -- (YOU);
\draw[->, poppurple, thick] (YOU) -- (Ag);
\draw[->, poppurple, thick] (Ag) -- (V);
\end{tikzpicture}
\end{center}
\end{propriete}

\begin{exemplebox}[Contraste \cle{被} vs \cle{由}]
\begin{center}
\begin{tabularx}{\linewidth}{|X|X|X|}
\hline \rowcolor{poporangeL} \textbf{Phrase} & \textbf{Structure} & \textbf{Nuance / Registre} \\ \hline
\zh{遗产被他们分了。} & \cle{被} (bèi) & Oral, neutre, focus sur l'action subie. \\ \hline
\zh{遗产由律师分配。} & \cle{由} (yóu) & Écrit, formel, focus sur la procédure/l'agent officiel. \\ \hline
\zh{房子被卖了。} & \cle{被} & Constat, parfois négatif (malchance). \\ \hline
\zh{房子由公证处处理。} & \cle{由} & Procédure administrative normale. \\ \hline
\end{tabularx}
\end{center}
\end{exemplebox}


\begin{aretenir}[Connecteurs logiques pour la narration (Niveau Terminale)]
\begin{center}
\begin{tabularx}{\linewidth}{|X|X|X|}
\hline \rowcolor{poppurpleL} \textbf{Fonction} & \textbf{Connecteur (Caractères / Pinyin)} & \textbf{Exemple d'emploi court} \\ \hline
\textbf{Chronologie / Suite} & \zh{首先} (shǒuxiān) — D'abord & \zh{首先，我们通知了亲戚。} \\ \hline
& \zh{接着} (jiēzhe) — Ensuite, puis & \zh{接着，大家坐下来谈。} \\ \hline
& \zh{然后} (ránhòu) — Puis, après & \zh{然后，律师念了遗嘱。} \\ \hline
& \zh{最后} (zuìhòu) — Finalement & \zh{最后，大家签了字。} \\ \hline
\textbf{Contraste / Concession} & \zh{但是} (dànshì) — Mais & \zh{但是，二叔不高兴。} \\ \hline
& \zh{不过} (bùguò) — Toutefois, cependant & \zh{不过，法律公平。} \\ \hline
& \zh{虽然……但是……} (suīrán... dànshì...) — Bien que... & \zh{虽然累，但是心安。} \\ \hline
\textbf{Cause / Conséquence} & \zh{因为……所以……} (yīnwèi... suǒyǐ...) — Comme... donc... & \zh{因为有遗嘱，所以没吵架。} \\ \hline
& \zh{因此} (yīncǐ) — Par conséquent & \zh{因此，分配很快。} \\ \hline
\textbf{Addition / Renforcement} & \zh{而且} (érqiě) — De plus, en plus & \zh{分得公平，而且快。} \\ \hline
& \zh{另外} (lìngwài) — En outre, par ailleurs & \zh{另外，房子要修缮。} \\ \hline
\end{tabularx}
\end{center}
\end{aretenir}

\begin{methode}[Réussir l'écriture des caractères clés du thème]
Pour chaque caractère complexe, appliquez la méthode \textbf{« Radicaux - Structure - Traits - Mémorisation »} :
\begin{enumerate}
    \item \textbf{\zh{遗} (yí) :} Radicale \zh{辶} (chuò, marche). Haut : \zh{贵} (guì, précieux). \textit{Astuce :} Ce qui est \textbf{précieux} (\zh{贵}) et qu'on \textbf{laisse derrière} soi en partant (\zh{辶}) = l'héritage. \textbf{Ordre :} haut \zh{贵} (haut→bas, gauche→droite), puis \zh{辶} (point, horizontal, pli, point final).
    \item \textbf{\zh{产} (chǎn) :} Radicale \zh{生} (shēng, naître/vie). Partie supérieure : \zh{立} (lì, debout) + trait horizontal + \zh{厂} (hǎn, falaise) [mnémotechnique]. \textit{Astuce :} Ce qui \textbf{naît} (\zh{生}) de la production. \textbf{Attention :} Ne pas confondre avec \zh{生} (vie) ni \zh{胜} (victoire, radicale \zh{月}).
    \item \textbf{\zh{嘱} (zhǔ) :} Radicale \zh{口} (kǒu, bouche). Phonétique \zh{主} (zhǔ, maître). \textit{Astuce :} La \textbf{bouche} du \textbf{maître} donne des instructions = testament, enjoindre.
    \item \textbf{\zh{睦} (mù) :} Radicale \zh{目} (mù, œil). Phonétique \zh{莫} (mò). \textit{Astuce :} Regarder (\zh{目}) sans conflit (\zh{莫} $\approx$ non/négation archaïque) = harmonieux.
    \item \textbf{\zh{妥} (tuǒ) :} Radicale \zh{女} (nǚ, femme). Phonétique \zh{叟} (sǒu, vieillard). \textit{Astuce :} Comme une \textbf{femme} sage / un \textbf{vieillard} prudent = approprié, bien réglé.
\end{enumerate}
\end{methode}

\courssub{Proposition de production et corrigé commenté}
\begin{exoresolu}[Production modèle : 爷爷留下的房子]
\textbf{Texte modèle (Caractères + Pinyin) :}
\par
\zh{1. 上个月，住在雅温得郊外村庄的姆巴尔加爷爷去世了，享年八十五岁。} \\
\zh{(Shàng gè yuè, zhù zài Yǎwēndé jiāowài cūnzhuāng de Mǔbā'ěrjiā Yéye qùshì le, xiǎngnián bāshíwǔ suì.)} \\
\zh{2. 首先，爷爷生前很勤劳，留下了三项主要遗产：一栋结实的砖房、一片茂盛的可可园和市场上的一间小商店。} \\
\zh{(Shǒuxiān, yéye shēngqián hěn qínláo, liúxià le sān xiàng zhǔyào yíchǎn: yī dòng jiēshi de zhuānfáng, yī piàn màoshèng de kěkěyuán hé shìchǎng shàng de yī jiān xiǎo shāngdiàn.)} \\
\zh{3. 接着，全家人为了分配遗产，在大舅舅家里开了一次家庭会议。} \\
\zh{(Jiēzhe, quánjiārén wèile fēnpèi yíchǎn, zài dàjiùjiu jiālǐ kāi le yī cì jiātíng huìyì.)} \\
\zh{4. 虽然爷爷留有遗嘱，但是二叔对分配方案有不同意见，一度发生了小小的纠纷。} \\
\zh{(Suīrán yéye liú yǒu yízhǔ, dànshì èrshū duì fēnpèi fāng'àn yǒu bùtóng yìjiàn, yīdù fāshēng le xiǎoxiǎo de jiūfēn.)} \\
\zh{5. 后来，大舅舅把遗嘱拿出来，当着大家的面读给我们听。} \\
\zh{(Hòulái, dàjiùjiu \textbf{bǎ} yízhǔ ná chūlái, dāngzhe dàjiā de miàn dú gěi wǒmen tīng.)} \\
\zh{6. 遗嘱里清楚地写着：房子由长子继承，可可园由二叔管理，商店的收入由全家人平分。} \\
\zh{(Yízhǔ lǐ qīngchǔ de xiězhe: fángzi \textbf{yóu} zhǎngzǐ jìchéng, kěkěyuán \textbf{yóu} èrshū guǎnlǐ, shāngdiàn de shōurù \textbf{yóu} quánjiārén píngfēn.)} \\
\zh{7. 因此，大家都接受了这个安排，没有人再争吵。} \\
\zh{(Yīncǐ, dàjiā dōu jiēshòu le zhège ānpái, méiyǒu rén zài zhēngchǎo.)} \\
\zh{8. 最后，我们把爷爷的照片放在客厅正中，一起吃饭，纪念这位伟大的老人。} \\
\zh{(Zuìhòu, wǒmen \textbf{bǎ} yéye de zhàopiàn fàng zài kètīng zhèngzhōng, yīqǐ chīfàn, jìniàn zhè wèi wěidà de lǎorén.)} \\
\zh{9. 现在，全家人生活得很和睦，这就是爷爷留给我们的最宝贵的遗产。} \\
\zh{(Xiànzài, quánjiārén shēnghuó de hěn hémù, zhè jiùshì yéye liú gěi wǒmen de zuì bǎoguì de yíchǎn.)}

\par
\textbf{Traduction française :}
1. Le mois dernier, Grand-père Mbarga, qui habitait un village en périphérie de Yaoundé, est décédé à l'âge de 85 ans. 
2. D'abord, Grand-père était très travailleur de son vivant, il a laissé trois biens principaux : une maison en briques solide, une plantation de cacao florissante et une petite boutique au marché. 
3. Ensuite, toute la famille, pour répartir l'héritage, a tenu une réunion familiale chez le grand oncle. 
4. Bien que Grand-père ait laissé un testament, le deuxième oncle avait un avis différent sur le plan de répartition, une petite dispute a éclaté un moment. 
5. Plus tard, l'oncle a sorti le testament et l'a lu devant tout le monde. 
6. Le testament stipulait clairement : la maison est héritée par le fils aîné, la plantation de cacao est gérée par le deuxième oncle, les revenus de la boutique sont partagés équitablement par toute la famille. 
7. Par conséquent, tout le monde a accepté cet arrangement, personne ne s'est plus disputé. 
8. Finalement, nous avons mis la photo de Grand-père au milieu du salon, nous avons mangé ensemble, en mémoire de ce grand vieillard. 
9. Maintenant, toute la famille vit en harmonie, c'est ça l'héritage le plus précieux que Grand-père nous a laissé.

\tcblower
\textbf{Commentaire détaillé des choix de langue (Corrigé commenté) :}
\begin{itemize}
    \item \textbf{Respect des contraintes structurelles :}
    \begin{itemize}
        \item \cle{把} (Phrase 5) : \zh{大舅舅把遗嘱拿出来...读给我们听}. Objet \zh{遗嘱} connu/spécifique. Verbe \zh{拿出来} (résultat directionnel) + \zh{读给...听} (complément de direction/bénéficiaire). Structure parfaite.
        \item \cle{把} (Phrase 8) : \zh{我们把爷爷的照片放在客厅正中}. Objet \zh{照片} spécifique. Verbe \zh{放} + complément de lieu \zh{在客厅正中}. Correct.
        \item \cle{由} (Phrase 6) : Trois occurrences successives pour le style formel/juridique du testament. \zh{房子由长子继承} (Agent passif formel), \zh{可可园由二叔管理} (Agent), \zh{收入由全家人平分} (Agent). Montre la maîtrise du registre écrit.
    \end{itemize}
    \item \textbf{Connecteurs utilisés (6 au total, $>$ 4 requis) :}
    \zh{首先} (début énumération), \zh{接着} (suite chronologique), \zh{虽然...但是...} (contraste concession), \zh{后来} (suite temporelle), \zh{因此} (conséquence logique), \zh{最后} (clôture narrative). La progression est fluide.
    \item \textbf{Vocabulaire thématique ciblé :} \zh{遗产} (x3), \zh{遗嘱} (x2), \zh{分配} (v. et n.), \zh{继承}, \zh{管理}, \zh{平分}, \zh{纠纷}, \zh{和睦}, \zh{安排}, \zh{纪念}. Richesse lexicale HSK 3-4 adaptée au sujet.
    \item \textbf{Grammaire / Ordre des mots :}
    \begin{itemize}
        \item Compléments de lieu/temps anteposés : \zh{在大舅舅家里开了一次会议}, \zh{当着大家的面}.
        \item Structure \zh{把} correctement complétée (V + Complément).
        \item Passif formel \zh{由} bien distingué du \zh{被} oral.
        \item Aspects : \zh{了} (accompli Phr 1, 3, 5, 7, 8), \zh{着} (état continu Phr 6 \zh{写着}), \zh{过} (expérience implicite dans \zh{留有}).
    \end{itemize}
    \item \textbf{Contexte camerounais :} Noms propres (\zh{雅温得}、\zh{姆巴尔加}), biens réalistes (\zh{可可园}、\zh{砖房}、\zh{市场小商店}), structure familiale (\zh{大舅舅} autorité morale, \zh{二叔} bénéficiaire, \zh{长子} droit d'aînesse coutumier vs loi). Le texte sonne « vrai ».
    \item \textbf{Écriture des caractères :} Le texte contient les caractères cibles \zh{遗、产、嘱、继、承、分、睦、妥} (dans \zh{妥善} implicite par \zh{安排} ou à ajouter en réécriture). L'élève doit s'assurer de l'ordre des traits (ex: \zh{遗} : \zh{贵} puis \zh{辶}).
\end{itemize}
\end{exoresolu}

\courssub{Grille d'évaluation Bac (Expression écrite — 10 points)}
\begin{center}
\begin{tabularx}{\linewidth}{|X|c|X|}
\hline \rowcolor{popblueL} \textbf{Critère} & \textbf{Points} & \textbf{Indicateurs de réussite} \\ \hline
Respect des consignes & 2 & Titre présent, 8-10 phrases, \cle{把} (1+), \cle{由} (1+), 4 connecteurs+, ponctuation chinoise 。 \\ \hline
Cohérence & 2 & Structure en 4 étapes respectée, progression logique, connecteurs pertinents. \\ \hline
Grammaire (Structures imposées) & 2 & \cle{把} : objet connu + V+complément. \cle{由} : registre formel, place correcte. Pas d'erreur ordre mots (Temps/Lieu). \\ \hline
Lexique & 2 & Vocabulaire thème (héritage, famille, biens) utilisé correctement, variété (pas de répétition \zh{好}、\zh{很}). \\ \hline
Caractères & 1 & Écriture soignée, radicaux respectés, pas de confusion \zh{遗/贵}, \zh{产/生}, \zh{嘱/主}. \\ \hline
Ponctuation / Présentation & 1 & Ponctuation pleine chasse, paragraphage, lisibilité. \\ \hline
\end{tabularx}
\end{center}

\courssub{Bilan de l'activité}
Cette activité d'intégration a permis de valider l'acquisition des objectifs de la leçon 7 dans une situation de communication proche de l'examen du Baccalauréat (épreuve d'expression écrite et de traduction).
\begin{itemize}
    \item \textbf{Sur le plan linguistique :} Les élèves ont dû faire la preuve qu'ils savent discriminer l'usage de \cle{把} (action concrète sur objet connu, style narratif/quotidien) et de \cle{由} (relation logique, cause, passif formel, style administratif/juridique). L'obligation d'utiliser 4 connecteurs a forcé la structuration du discours (cohérence textuelle), point faible fréquent chez les apprenants francophones qui juxtaposent souvent des phrases simples (SVO) sans articulation logique.
    \item \textbf{Sur le plan graphique :} La réécriture finale après correction croisée a focalisé l'attention sur la justesse du trait (ex: le point final de \zh{辶} dans \zh{遗}、\zh{过}, la distinction \zh{产}/\zh{生}/\zh{胜}). L'échange de copies a transformé l'erreur en objet d'apprentissage collectif.
    \item \textbf{Sur le plan culturel :} L'ancrage camerounais (Mbarga, cacao, grand oncle médiateur, droit coutumier vs droit écrit) a donné du sens à la tâche. Les élèves ont compris que le chinois n'est pas une langue abstraite mais un outil pour dire \textit{leur} réalité. La conclusion du texte modèle (\zh{和睦就是最宝贵的遗产}) ancre la valeur universelle de la paix familiale au-delà des biens matériels.
    \item \textbf{Pistes de remédiation (si difficultés observées) :}
    \begin{itemize}
        \item Si \cle{把} est mal placé (objet indéfini) : Revoir la \textbf{Propriété} avec exercices de transformation \zh{我吃了苹果} $\rightarrow$ \zh{我把苹果吃了} (seulement si pomme connue).
        \item Si \cle{由} est confondu avec \zh{从} (provenance spatiale) : Contraster \zh{我从北京来} (espace) vs \zh{这个决定由我做} (agent/logique).
        \item Si connecteurs absents : Travail sur « Texte à trous » avec connecteurs imposés avant production libre.
    \end{itemize}
\end{itemize}
La grille d'évaluation utilisée (critères : Respect consignes / Cohérence / Grammaire structures imposées / Lexique / Caractères / Ponctuation) est directement calquée sur les attendus du Bac LV2. Les élèves repartent avec un modèle de texte \textbf{annoté et validé} dans leur portfolio, ressource précieuse pour la révision terminale.

\newpage

\coursec{Méthodes et exercices résolus}

\coursec{Leçon 7 — Expression écrite : l'héritage}

\courssub{Texte support : l'héritage du grand-père}

\begin{textebox}[Modèle de rédaction commenté]
\zh{我的爷爷去年去世了。他给我们留下了一座房子和一块地。房子在雅温得，地在农村。爸爸和姑姑把财产分好了：爸爸继承了房子，姑姑继承了地。我觉得分遗产的时候，家庭和睦最重要。}\\
(Wǒ de yéye qùnián qùshì le. Tā gěi wǒmen liúxià le yí zuò fángzi hé yí kuài dì. Fángzi zài Yǎwēndé, dì zài nóngcūn. Bàba hé gūgu bǎ cáichǎn fēn hǎo le : bàba jìchéng le fángzi, gūgu jìchéng le dì. Wǒ juéde fēn yíchǎn de shíhou, jiātíng hémù zuì zhòngyào.)\\
« Mon grand-père est décédé l'année dernière. Il nous a laissé une maison et un terrain. La maison est à Yaoundé, le terrain à la campagne. Papa et ma tante ont bien partagé les biens : papa a hérité de la maison, ma tante du terrain. Je pense qu'au moment de partager un héritage, l'harmonie familiale est le plus important. »
\end{textebox}

\begin{aretenir}[Analyse du modèle]
\begin{itemize}
\item \textbf{Chronologie claire} : décès (\zh{去世了}) $\rightarrow$ biens (\zh{留下}) $\rightarrow$ localisation $\rightarrow$ partage (\zh{把...分好}) $\rightarrow$ avis (\zh{我觉得}).
\item \textbf{Classificateurs obligatoires} : \zh{一座房子} (bâtiment), \zh{一块地} (terrain).
\item \textbf{Construction \zh{把}} pour le partage : \zh{把财产分好} (objet défini \zh{财产} déplacé avant le verbe + résultat \zh{好}).
\item \textbf{Complément de temps} \zh{分遗产的时候} placé \emph{avant} le verbe \zh{est} (omniprésent en chinois).
\end{itemize}
\end{aretenir}

\courssub{Lexique et glossaire du thème}

\begin{definition}[Vocabulaire essentiel — Héritage et succession]
\begin{center}
\begin{tabularx}{\linewidth}{|X|X|X|X|}\hline
\rowcolor{popblueL} Caractères & Pinyin & Nature & Traduction \\ \hline
\zh{遗产} & yíchǎn & nom & héritage, succession \\ \hline
\zh{继承} & jìchéng & verbe & hériter (de), succéder à \\ \hline
\zh{遗嘱} & yízhǔ & nom & testament \\ \hline
\zh{留给} & liú gěi & verbe & laisser à, léguer à \\ \hline
\zh{分给} & fēn gěi & verbe & distribuer à, partager à \\ \hline
\zh{分成} & fēn chéng & verbe & diviser en, partager en \\ \hline
\zh{去世} & qùshì & verbe & décéder, mourir (terme respectueux) \\ \hline
\zh{留下} & liú xià & verbe & laisser (derrière soi) \\ \hline
\zh{财产} & cáichǎn & nom & biens, patrimoine, propriété \\ \hline
\zh{和睦} & hémù & adj. & harmonieux, uni (famille, relations) \\ \hline
\zh{平分} & píng fēn & verbe & partager équitablement, diviser en parts égales \\ \hline
\zh{吵架} & chǎo jià & verbe & se disputer, se quereller \\ \hline
\zh{榜样} & bǎng yàng & nom & exemple, modèle \\ \hline
\zh{按照} & àn zhào & prép. & selon, conformément à, d'après \\ \hline
\zh{大儿子} & dà érzi & nom & fils aîné \\ \hline
\zh{姑姑} & gū gu & nom & tante (sœur du père) \\ \hline
\zh{奶奶} & nǎi nai & nom & grand-mère paternelle \\ \hline
\zh{爷爷} & yé ye & nom & grand-père paternel \\ \hline
\zh{座} & zuò & classif. & classificateur : maisons, bâtiments, montagnes \\ \hline
\zh{块} & kuài & classif. & classificateur : terrains, morceaux, pâtés \\ \hline
\zh{份} & fèn & classif. & classificateur : documents, parts, cadeaux, journaux \\ \hline
\end{tabularx}
\end{center}
\end{definition}

\courssub{Méthode 1 : Rédiger un texte simple sur l'héritage}

\begin{methode}[Comment construire un texte sur l'héritage]
Pour produire un texte simple et correct sur le thème de l'héritage (\zh{遗产} yíchǎn), suis ces étapes :
\begin{enumerate}
\item \textbf{Présenter la situation} : qui est décédé, quand. Exemple : \zh{我爷爷去年去世了。} (Wǒ yéye qùnián qùshì le.) « Mon grand-père est décédé l'année dernière. »
\item \textbf{Énumérer les biens} avec des classificateurs corrects : \zh{一座房子} (yí zuò fángzi, une maison), \zh{一块地} (yí kuài dì, un terrain), \zh{一些钱} (yìxiē qián, de l'argent), \zh{一份遗嘱} (yí fèn yízhǔ, un testament).
\item \textbf{Décrire le partage} avec les verbes du thème : \zh{继承} (jìchéng, hériter), \zh{分给} (fēn gěi, répartir à), \zh{留给} (liú gěi, laisser à), \zh{平分} (píng fēn, partager équitablement).
\item \textbf{Donner un avis ou une conclusion} avec un connecteur : \zh{所以} (suǒyǐ, donc), \zh{但是} (dànshì, mais), \zh{我觉得} (wǒ juéde, je pense que).
\item \textbf{Relire} : vérifier les tons, l'ordre \cle{Sujet + Temps + Lieu + Verbe + Objet}, la place du complément de temps \emph{avant} le verbe, et les mots-outils (\zh{了}, \zh{的}, \zh{把}).
\end{enumerate}
Structure type à retenir : \cle{Sujet + 把 + Objet + Verbe + Résultat / Complément} (construction en \zh{把} pour le partage d'un objet défini).
\end{methode}

\begin{exoresolu}[Rédaction guidée : l'héritage du grand-père]
Énoncé : Rédige un texte de 5 à 6 phrases en chinois : ton grand-père est décédé ; il a laissé une maison à Yaoundé et un terrain au village ; ton père et ta tante ont partagé les biens ; donne ton avis sur l'importance de la famille.
\tcblower
\textbf{Étape 1 — Plan.} Décès $\rightarrow$ biens $\rightarrow$ localisation $\rightarrow$ partage $\rightarrow$ avis.\\
\textbf{Étape 2 — Choix des structures.} Passé accompli avec \zh{了} ; construction \zh{把} pour le partage (\zh{把财产分好}) ; avis avec \zh{我觉得}.\\
\textbf{Étape 3 — Rédaction.}\\
\zh{我的爷爷去年去世了。他给我们留下了一座房子和一块地。房子在雅温得，地在农村。爸爸和姑姑把财产分好了：爸爸继承了房子，姑姑继承了地。我觉得分遗产的时候，家庭和睦最重要。}\\
(Wǒ de yéye qùnián qùshì le. Tā gěi wǒmen liúxià le yí zuò fángzi hé yí kuài dì. Fángzi zài Yǎwēndé, dì zài nóngcūn. Bàba hé gūgu bǎ cáichǎn fēn hǎo le : bàba jìchéng le fángzi, gūgu jìchéng le dì. Wǒ juéde fēn yíchǎn de shíhou, jiātíng hémù zuì zhòngyào.)\\
« Mon grand-père est décédé l'année dernière. Il nous a laissé une maison et un terrain. La maison est à Yaoundé, le terrain à la campagne. Papa et ma tante ont bien partagé les biens : papa a hérité de la maison, ma tante du terrain. Je pense qu'au moment de partager un héritage, l'harmonie familiale est le plus important. »\\
\textbf{Justifications.} \zh{去世} est le terme respectueux pour « décéder » ; \zh{把财产分好} utilise la construction \zh{把} (objet déplacé avant le verbe + résultat) ; le complément de temps \zh{分遗产的时候} précède le verbe ; \zh{最} marque le superlatif.
\end{exoresolu}

\courssub{Méthode 2 : Écrire correctement les caractères du thème}

\begin{methode}[Traits, radicaux et caractères proches]
\begin{enumerate}
\item \textbf{Identifier le radical (clé)} : il donne l'indice sémantique. \zh{财} (cái, richesse) porte la clé \zh{贝} (coquillage $\rightarrow$ argent) ; \zh{继} (jì, continuer) porte la clé \zh{纟} (soie, fil continu) ; \zh{留} (liú, laisser/rester) porte la clé \zh{田} (champ, lieu où l'on reste).
\item \textbf{Respecter l'ordre des traits} : de haut en bas, de gauche à droite, l'horizontal avant le vertical, l'extérieur avant l'intérieur, on ferme la boîte en dernier.
\item \textbf{Distinguer les caractères proches (pièges visuels)} :
\begin{itemize}
\item \zh{遗} (yí, laisser/transmettre, clé \zh{辶} de la marche) $\neq$ \zh{贵} (guì, cher/précieux, clé \zh{贝}).
\item \zh{继} (jì, continuer, clé \zh{纟}) $\neq$ \zh{断} (duàn, couper, clé \zh{斤}).
\item \zh{份} (fèn, part/classificateur, clé \zh{亻} homme) $\neq$ \zh{分} (fēn, diviser/partager, clé \zh{刀} sabre/division).
\end{itemize}
\item \textbf{Mémoriser par mots entiers (chunks)} : apprendre \zh{遗产} (héritage), \zh{继承} (hériter), \zh{遗嘱} (testament), \zh{财产} (patrimoine) comme des blocs indissociables, pas caractère par caractère.
\end{enumerate}
\end{methode}

\begin{exoresolu}[Caractères proches : choisir la bonne graphie]
Énoncé : Complète avec le bon caractère : (\zh{遗} / \zh{贵}) ; (\zh{份} / \zh{分}).
\begin{enumerate}
\item \zh{爷爷的\_\_产是一座房子。}
\item \zh{这件礼物太\_\_重了！}
\item \zh{他把地\_\_给了两个儿子。}
\item \zh{每\_\_文件都很重要。}
\end{enumerate}
\tcblower
1. \zh{遗产} (yíchǎn) : « héritage ». \zh{遗} = laisser après soi (clé \zh{辶}) ; \zh{贵} = cher, ne convient pas. « L'héritage de grand-père est une maison. »\\
2. \zh{贵重} (guìzhòng) : « précieux ». Ici c'est bien \zh{贵} (cher, précieux, clé \zh{贝}). « Ce cadeau est trop précieux ! »\\
3. \zh{分给} (fēn gěi) : « répartir à », verbe d'action $\rightarrow$ \zh{分} (diviser, clé \zh{刀}). « Il a réparti le terrain entre ses deux fils. »\\
4. \zh{每份文件} (měi fèn wénjiàn) : \zh{份} est le classificateur des documents et des parts ; il prend le radical \zh{亻} (homme) à gauche. « Chaque document est important. »\\
\textbf{Réflexe.} \zh{分} est un verbe (diviser) ; \zh{份} est un nom/classificateur (une part). La présence du radical \zh{亻} signale le classificateur.
\end{exoresolu}

\courssub{Structures et connecteurs du texte}

\begin{propriete}[Connecteurs logiques indispensables]
\begin{center}
\begin{tabularx}{\linewidth}{|X|X|X|}\hline
\rowcolor{popblueL} Connecteur & Pinyin & Emploi et exemple \\ \hline
\zh{先……然后……} & xiān... ránhòu... & D'abord... ensuite... (chronologie) \\ \hline
\zh{因为……所以……} & yīnwèi... suǒyǐ... & Cause $\rightarrow$ Conséquence (les deux termes coexistent en chinois) \\ \hline
\zh{虽然……但是……} & suīrán... dànshì... & Concession : Bien que... mais... \\ \hline
\zh{按照} & ànzhào & Selon (la loi, la tradition, le testament) — placé en tête ou avant le verbe \\ \hline
\zh{如果……就……} & rúguǒ... jiù... & Hypothèse : Si... alors... \\ \hline
\end{tabularx}
\end{center}
\textbf{Règle d'or :} En chinois, on utilise souvent \emph{les deux} éléments de la paire (\zh{因为} et \zh{所以} ensemble ; \zh{虽然} et \zh{但是} ensemble), contrairement au français où « parce que » et « donc » ne se cumulent pas. \zh{按照} + Nom fonctionne comme un complément de mode/référence.
\end{propriete}

\begin{exoresolu}[Relier les phrases avec connecteurs]
Énoncé : Relie les phrases avec le connecteur indiqué. Traduis le résultat en français.
\begin{enumerate}
\item \zh{爷爷留下了遗嘱。孩子们没有吵架。} (\zh{因为……所以……})
\item \zh{遗产不多。家人很和睦。} (\zh{虽然……但是……})
\end{enumerate}
\tcblower
1. \zh{因为爷爷留下了遗嘱，所以孩子们没有吵架。} (Yīnwèi yéye liúxià le yízhǔ, suǒyǐ háizimen méiyǒu chǎojià.)\\
« Parce que grand-père avait laissé un testament, les enfants ne se sont pas disputés. » (Les deux connecteurs coexistent en chinois ; en français, on ne garde que « parce que » ou « donc »).\\
2. \zh{虽然遗产不多，但是家人很和睦。} (Suīrán yíchǎn bù duō, dànshì jiārén hěn hémù.)\\
« Bien que l'héritage soit modeste, la famille est très unie. »\\
\textbf{Justifications.} Négation du passé : \zh{没有} + verbe (\zh{没有吵架}), jamais \zh{不} + verbe pour le passé. \zh{不多} = « peu important / pas beaucoup » ; adjectif \zh{和睦} précédé de \zh{很} (adverbe de degré obligatoire pour l'adjectif attribut).
\end{exoresolu}

\courssub{Méthode 3 : Thème — traduire du français vers le chinois}

\begin{methode}[Traduire une phrase sur l'héritage (français $\rightarrow$ chinois)]
\begin{enumerate}
\item \textbf{Repérer le sujet et le verbe principal}, puis placer les compléments de temps et de lieu \emph{avant} le verbe. Ex. : « Mon oncle a laissé un testament hier » $\rightarrow$ \zh{我叔叔昨天留下了遗嘱。} (Temps \zh{昨天} avant le verbe).
\item \textbf{Traduire « hériter de » par \zh{继承}} (verbe transitif direct) : \zh{她继承了母亲的房子。} (Tā jìchéng le mǔqīn de fángzi.) « Elle a hérité de la maison de sa mère. »
\item \textbf{Choisir le bon classificateur} (obligatoire après nombre/démonstratif) :
\begin{itemize}
\item Maison / Bâtiment : \zh{座} (zuò) $\rightarrow$ \zh{一座房子}
\item Terrain / Morceau : \zh{块} (kuài) $\rightarrow$ \zh{一块地}
\item Testament / Document / Part : \zh{份} (fèn) $\rightarrow$ \zh{一份遗嘱}, \zh{三份}
\item Voiture : \zh{辆} (liàng) $\rightarrow$ \zh{一辆车}
\end{itemize}
\item \textbf{Marquer l'accompli} avec \zh{了} après le verbe (ou après le complément de résultat) si l'action est réalisée dans le passé.
\item \textbf{Ne pas calquer le français} : « selon la tradition » = \zh{按照传统} (ànzhào chuántǒng), jamais mot à mot. « Laisser à » = \zh{留给} (liú gěi) + personne + objet.
\end{enumerate}
\end{methode}

\begin{exoresolu}[Thème : trois phrases type Bac]
Énoncé : Traduis en chinois :
\begin{enumerate}
\item « Ma grand-mère a laissé un terrain à ses enfants. »
\item « Selon la tradition, le fils aîné hérite de la maison. »
\item « Ils ont partagé l'argent en trois parts. »
\end{enumerate}
\tcblower
1. \zh{我奶奶留给孩子们一块地。} (Wǒ nǎinai liú gěi háizimen yí kuài dì.) Verbe \zh{留给} (laisser à) + classificateur \zh{块} pour \zh{地}.\\
2. \zh{按照传统，大儿子继承房子。} (Ànzhào chuántǒng, dà érzi jìchéng fángzi.) « Selon » = \zh{按照} en tête ; « fils aîné » = \zh{大儿子}. (Présent de vérité générale : pas de \zh{了}).\\
3. \zh{他们把钱分成了三份。} (Tāmen bǎ qián fēn chéng le sān fèn.) Construction \zh{把} (objet \zh{钱} défini) + \zh{分成} (diviser en) + \zh{份} (parts) + \zh{了} (accompli).\\
\textbf{Conclusion.} En thème, la réussite tient à trois points : ordre des mots (temps/lieu avant le verbe), classificateur juste, construction \zh{把} quand l'objet défini subit une transformation.
\end{exoresolu}

\courssub{Méthode 4 : Version — traduire du chinois vers le français}

\begin{methode}[Comprendre et traduire un texte chinois sur l'héritage]
\begin{enumerate}
\item \textbf{Lire deux fois} : une fois pour le sens global, une fois en soulignant les mots du thème (\zh{遗产}, \zh{继承}, \zh{遗嘱}, \zh{分}, \zh{把}, \zh{了}).
\item \textbf{Découper la phrase} : Sujet | Compléments (temps, lieu, \zh{按照}...) | Verbe (+ \zh{了}) | Objet. Traduire dans l'ordre français : Sujet – Verbe – Objet, puis compléments.
\item \textbf{Traduire les mots-outils} :
\begin{itemize}
\item \zh{把} : ne se traduit pas (signale l'objet déplacé avant le verbe).
\item \zh{了} : action accomplie $\rightarrow$ passé composé (ou plus-que-parfait selon contexte).
\item \zh{的} : possession (\zh{他的房子} = sa maison) ou relative (\zh{他留下的遗产} = l'héritage qu'il a laissé).
\end{itemize}
\item \textbf{Rendre en français naturel} : \zh{家庭和睦} $\rightarrow$ « l'harmonie familiale » (pas « la famille en paix ») ; \zh{平分} $\rightarrow$ « partager équitablement » ; \zh{榜样} $\rightarrow$ « un exemple / un modèle ».
\end{enumerate}
\end{methode}

\begin{exoresolu}[Version : texte court type examen]
Énoncé : Traduis en français :
\zh{老人去世以前写了一份遗嘱。按照遗嘱，他的两个孩子平分了遗产：儿子得到房子，女儿得到钱。他们没有吵架，村里人都说这是一个好榜样。}
\tcblower
\textbf{Découpage.} \\
Phrase 1 : \zh{老人} (sujet) + \zh{去世以前} (temps : avant de mourir) + \zh{写了} (verbe accompli) + \zh{一份遗嘱} (objet + classif.). \\
Phrase 2 : \zh{按照遗嘱} (selon le testament) + \zh{他的两个孩子} (sujet) + \zh{平分了} (verbe accompli) + \zh{遗产} (objet) ; \zh{儿子得到房子，女儿得到钱} (deux clauses coordonnées : obtenir/recevoir). \\
Phrase 3 : \zh{他们没有吵架} (négation passée) ; \zh{村里人都说这是一个好榜样} (jugement). \\
\textbf{Traduction.} « Avant de mourir, le vieil homme avait écrit un testament. Selon le testament, ses deux enfants ont partagé l'héritage en parts égales : le fils a reçu la maison, la fille a reçu l'argent. Ils ne se sont pas disputés, et tous les gens du village disent que c'est un bon exemple. » \\
\textbf{Points de vigilance.} \zh{去世以前} se traduit « avant de mourir » (complément de temps antéposé en chinois, postposé en français) ; \zh{平分} = « partager équitablement » ; ne pas traduire \zh{的} systématiquement par « de » quand il structure une relative : \zh{他留下的遗产} = « l'héritage qu'il a laissé ».
\end{exoresolu}

\courssub{Grille d'évaluation et entraînement Bac}

\begin{aretenir}[Grille d'évaluation — Expression écrite (20 pts)]
\begin{center}
\begin{tabularx}{\linewidth}{|X|X|X|}\hline
\rowcolor{popblueL} Critère & Indicateurs & Points \\ \hline
\textbf{Contenu / Respect du sujet} & Idées claires, situation décrite, biens énumérés, avis donné, longueur (5-6 phrases min.) & 4 \\ \hline
\textbf{Vocabulaire} & Mots du thème utilisés (\zh{遗产, 继承, 遗嘱, 留给, 分给, 和睦}), classificateurs corrects (\zh{座, 块, 份}), variété lexicale & 4 \\ \hline
\textbf{Grammaire / Structures} & Ordre Sujet-Temps-Lieu-Verbe-Objet, \zh{把} correct, \zh{了} accompli, négation \zh{没有}, connecteurs (\zh{因为/所以}, \zh{虽然/但是}) & 4 \\ \hline
\textbf{Caractères / Écriture} & Graphies exactes, pas de confusion (\zh{遗/贵}, \zh{份/分}, \zh{继/断}), traits dans l'ordre, radicaux respectés & 3 \\ \hline
\textbf{Cohérence / Connecteurs} & Enchaînement logique, ponctuation chinoise (，。), paragraphage si nécessaire & 3 \\ \hline
\textbf{Présentation / Soin} & Lisibilité, pinyin non requis en production finale mais caractères soignés & 2 \\ \hline
\end{tabularx}
\end{center}
\end{aretenir}

\begin{exercice}[Entraînement Bac — Sujet complet]
\textbf{Sujet :} « Ton oncle est décédé l'année dernière. Il a laissé un testament. Selon le testament, ta mère et ton oncle (frère cadet) ont partagé l'héritage équitablement : ta mère a eu la maison, ton oncle a eu l'argent. Écris un message à ton ami chinois pour lui raconter cela et exprimer ton opinion sur l'importance du testament. (60-80 caractères) »\\
\textbf{Consignes :} Utilise au moins 5 mots du lexique, 2 connecteurs, la construction \zh{把} une fois, et soigne les classificateurs.
\end{exercice}

\begin{corrige}[Exercice Entraînement Bac]
\textbf{Proposition de rédaction :}\\
\zh{我叔叔去年去世了。他留下了一份遗嘱。按照遗嘱，妈妈和叔叔平分了遗产：妈妈继承了房子，叔叔得到钱。我觉得立遗嘱很有必要，能让家人不吵架，和睦相处。}\\
(Wǒ shūshu qùnián qùshì le. Tā liúxià le yí fèn yízhǔ. Ànzhào yízhǔ, māma hé shūshu píng fēn le yíchǎn : māma jìchéng le fángzi, shūshu dédào qián. Wǒ juéde lì yízhǔ hěn yǒu bìyào, néng ràng jiārén bù chǎojià, hémù xiāngchǔ.)\\
« Mon oncle est décédé l'an dernier. Il a laissé un testament. Selon le testament, maman et l'oncle ont partagé l'héritage équitablement : maman a hérité de la maison, l'oncle a reçu l'argent. Je pense qu'écrire un testament est très nécessaire, cela permet à la famille de ne pas se disputer et de vivre en harmonie. »\\
\textbf{Analyse :} 6 phrases. Lexique : \zh{去世, 遗嘱, 按照, 平分, 遗产, 继承, 和睦} (7 mots). Connecteurs : \zh{按照}, \zh{因为/所以} (implicite dans \zh{能让...}), structure \zh{把} non forcée ici car \zh{平分} prend objet direct, mais \zh{把遗产平分了} possible. Classificateurs : \zh{份} (遗嘱), \zh{座} (sous-entendu pour maison), pas de classificateur pour \zh{钱} (masse). Caractères soignés.
\end{corrige}

\begin{attention}[Les 5 erreurs qui coûtent des points au Bac]
\begin{enumerate}
\item \textbf{Ordre des mots (calque français)} : placer le complément de temps \emph{après} le verbe. Faux : \zh{他去世了去年}. Juste : \zh{他去年去世了。} \cle{Règle : Temps / Lieu avant le Verbe.}
\item \textbf{Caractères confondus (pièges visuels)} : écrire \zh{分} à la place de \zh{份} (une part / classif.), \zh{贵} à la place de \zh{遗} (héritage). \cle{Vérifie le radical : \zh{亻} pour \zh{份}, \zh{辶} pour \zh{遗}.}
\item \textbf{Classificateur faux ou oublié} : une maison = \zh{一座房子} (pas \zh{一个}), un terrain = \zh{一块地}, un testament = \zh{一份遗嘱}. \cle{Nom dénombré sans classif. = faute grave.}
\item \textbf{Négation du passé} : dire \zh{不吵架} pour « ils ne se sont pas disputés ». \cle{Passé négatif = \zh{没有} + V : \zh{他们没有吵架。}}
\item \textbf{Mots-outils mal employés} : oublier \zh{了} pour une action accomplie (\zh{他继承了房子}), ou croire que \zh{把} se traduit. \cle{\zh{把} ne se traduit pas ; il signale : Objet Défini + Transformation.}
\end{enumerate}
\end{attention}

\newpage

\coursec{Exercices}

\courssub{Je vérifie mes connaissances}

\begin{exercice}[QCM : le vocabulaire de l'héritage]
Pour chaque question, choisis la bonne réponse (a, b ou c).
\begin{enumerate}
\item « hériter » se dit : \quad a) \zh{继承} \quad b) \zh{继续} \quad c) \zh{纪念}
\item « testament » se dit : \quad a) \zh{法律} \quad b) \zh{遗嘱} \quad c) \zh{财产}
\item \zh{去世} signifie : \quad a) « partir en voyage » \quad b) « déménager » \quad c) « décéder »
\item « bien, patrimoine » se dit : \quad a) \zh{财产} \quad b) \zh{材料} \quad c) \zh{家庭}
\item Dans la phrase \zh{房子由我爸爸继承。}, \zh{由} introduit : \quad a) le temps \quad b) l'agent (celui qui fait l'action) \quad c) le lieu
\end{enumerate}
\end{exercice}

\begin{exercice}[Vrai ou faux ? Justifie ta réponse]
Dis si chaque affirmation est vraie (V) ou fausse (F), puis justifie en une phrase en français.
\begin{enumerate}
\item Le caractère \zh{遗} contient le radical de la marche \zh{辶}.
\item \zh{遗嘱} et \zh{属于} se prononcent exactement de la même façon.
\item Dans la structure \zh{虽然……但是……}, les deux mots peuvent être utilisés ensemble dans la même phrase en chinois.
\item \zh{外公} désigne le grand-père paternel.
\item Pour dire « il y a deux ans », on peut dire \zh{两年前}.
\end{enumerate}
\end{exercice}

\begin{exercice}[Trouve le mot du lexique]
Complète chaque définition avec le mot chinois qui convient (caractères + pinyin).
\begin{enumerate}
\item Ce qu'une personne laisse après sa mort (maison, champ, argent) : \ldots\ldots\ldots
\item Le document écrit qui dit à qui reviennent les biens : \ldots\ldots\ldots
\item Le frère de la mère : \ldots\ldots\ldots
\item L'ensemble des règles écrites d'un pays : \ldots\ldots\ldots
\item Recevoir les biens d'un parent décédé : \ldots\ldots\ldots
\end{enumerate}
\end{exercice}

\begin{exercice}[Associe le caractère à son pinyin]
Relie chaque caractère de la colonne A à son pinyin de la colonne B.
\begin{center}
\begin{tabularx}{\linewidth}{|X|X|}
\hline
\rowcolor{popblueL} A : Caractères & B : Pinyin \\
\hline
\zh{遗} & chǎn \\
\hline
\zh{嘱} & jì \\
\hline
\zh{财} & yí \\
\hline
\zh{产} & zhǔ \\
\hline
\zh{继} & cái \\
\hline
\end{tabularx}
\end{center}
\end{exercice}

\courssub{Je m'entraîne}

\begin{exercice}[Traduis en français]
Traduis les phrases suivantes en français correct.
\begin{enumerate}
\item \zh{我的爷爷前年去世了。}\par \emph{Wǒ de yéye qiánnián qùshì le.}
\item \zh{他留下了一所房子和一块地。}\par \emph{Tā liúxià le yì suǒ fángzi hé yí kuài dì.}
\item \zh{按照他的遗嘱，房子由我爸爸继承。}\par \emph{Ànzhào tā de yízhǔ, fángzi yóu wǒ bàba jìchéng.}
\item \zh{虽然我们很难过，但是我们家还是很团结。}\par \emph{Suīrán wǒmen hěn nánguò, dànshì wǒmen jiā háishi hěn tuánjié.}
\end{enumerate}
\end{exercice}

\begin{exercice}[Transforme les phrases]
\begin{enumerate}
\item Réécris la phrase \zh{爸爸继承了房子。} (\emph{Bàba jìchéng le fángzi.}) avec la structure \zh{把}.
\item Réécris la même phrase avec la structure \zh{由} (l'agent après \zh{由}).
\item Transforme \zh{妈妈把钱分给了孩子们。} en phrase avec \zh{被}.
\end{enumerate}
\end{exercice}

\begin{exercice}[Texte à trous : les connecteurs]
Complète le texte suivant avec les connecteurs : \zh{首先}、\zh{然后}、\zh{因为……所以……}、\zh{虽然……但是……}、\zh{最后}.
\begin{textebox}[Le partage de l'héritage]
我的舅舅上个月去世了。\ldots\ldots\ldots 我们全家都很难过，\ldots\ldots\ldots 我们必须处理他的财产。\ldots\ldots\ldots{}，我们看了他的遗嘱；\ldots\ldots\ldots{}，我们开了一个家庭会议；\ldots\ldots\ldots{}，大家同意：房子给我的妈妈，钱给我的表哥。\ldots\ldots\ldots 事情不多，\ldots\ldots\ldots 我们谈了很长时间。\\
\emph{Wǒ de jiùjiu shàng ge yuè qùshì le. \ldots wǒmen quán jiā dōu hěn nánguò, \ldots wǒmen bìxū chǔlǐ tā de cáichǎn. \ldots, wǒmen kàn le tā de yízhǔ ; \ldots, wǒmen kāi le yí ge jiātíng huìyì ; \ldots, dàjiā tóngyì : fángzi gěi wǒ de māma, qián gěi wǒ de biǎogē. \ldots shìqing bù duō, \ldots wǒmen tán le hěn cháng shíjiān.}\\
« Mon oncle est décédé le mois dernier. \ldots toute notre famille était très triste, \ldots nous devons régler ses biens. \ldots, nous avons lu son testament ; \ldots, nous avons tenu une réunion de famille ; \ldots, tout le monde a accepté : la maison pour ma mère, l'argent pour mon cousin. \ldots l'affaire était simple, \ldots nous avons longtemps discuté. »
\end{textebox}
\end{exercice}

\begin{exercice}[Discrimination des caractères]
Choisis le bon caractère pour compléter chaque phrase : \zh{遗}/\zh{遣}, \zh{嘱}/\zh{属}, \zh{财}/\zh{材}.
\begin{enumerate}
\item 爷爷留下了很多 \zh{}\ldots\zh{} 产。(\emph{yéye liúxià le hěn duō \ldots chǎn.})
\item 他的 \zh{}\ldots\zh{} 嘱写得很清楚。(\emph{tā de \ldots zhǔ xiě de hěn qīngchu.})
\item 这所房子 \zh{}\ldots\zh{} 于我的爸爸。(\emph{zhè suǒ fángzi \ldots yú wǒ de bàba.})
\item 这些木 \zh{}\ldots\zh{} 是从市场上买的。(\emph{zhèxiē mù \ldots shì cóng shìchǎng shang mǎi de.})
\item 他的 \zh{}\ldots\zh{} 体安葬在村子里。(\emph{tā de \ldots tǐ ānzàng zài cūnzi li.})
\end{enumerate}
\end{exercice}

\begin{exercice}[Remets les mots dans l'ordre]
Remets les mots dans le bon ordre pour faire une phrase correcte, puis traduis-la en français.
\begin{enumerate}
\item \zh{继承 / 房子 / 由 / 的 / 我 / 叔叔}
\item \zh{去世 / 去年 / 外公 / 了 / 的 / 我}
\item \zh{遗嘱 / 按照 / 财产 / 分了 / 他的 / 大家}
\end{enumerate}
\end{exercice}

\courssub{Je me prépare au Bac}

\begin{exercice}[Compréhension de l'écrit et questions de langue]
Lis le texte suivant, puis réponds aux questions.
\begin{textebox}[L'héritage de grand-père]
我的爷爷是杜阿拉的一个农民。他工作了四十年，留下了一所房子、两块地和一些钱。去年他去世了，全家都很难过。按照他的遗嘱，房子由我爸爸继承，地由我的两个叔叔分，钱留给我的奶奶。虽然大家开始有一点不同的意见，但是开了一个家庭会议以后，每个人都同意了。爷爷常说：「财产不重要，家庭团结最重要。」我们都记住了这句话。\\
\emph{Wǒ de yéye shì Dù'ālā de yí ge nóngmín. Tā gōngzuò le sìshí nián, liúxià le yì suǒ fángzi, liǎng kuài dì hé yìxiē qián. Qùnián tā qùshì le, quán jiā dōu hěn nánguò. Ànzhào tā de yízhǔ, fángzi yóu wǒ bàba jìchéng, dì yóu wǒ de liǎng ge shūshu fēn, qián liú gěi wǒ de nǎinai. Suīrán dàjiā kāishǐ yǒu yìdiǎn bùtóng de yìjian, dànshì kāi le yí ge jiātíng huìyì yǐhòu, měi ge rén dōu tóngyì le. Yéye cháng shuō : « Cáichǎn bù zhòngyào, jiātíng tuánjié zuì zhòngyào. » Wǒmen dōu jìzhu le zhè jù huà.}\\
« Mon grand-père était un paysan de Douala. Il a travaillé quarante ans et a laissé une maison, deux champs et un peu d'argent. L'année dernière, il est décédé et toute la famille était très triste. Selon son testament, la maison revient à mon père, les champs sont partagés entre mes deux oncles, l'argent est laissé à ma grand-mère. Bien qu'au début chacun ait eu des avis un peu différents, après une réunion de famille, tout le monde a accepté. Grand-père disait souvent : "Les biens ne sont pas importants, l'union de la famille est le plus important." Nous avons tous retenu cette phrase. »
\end{textebox}
\textbf{A. Questions de compréhension} (réponds en français)
\begin{enumerate}
\item Quel était le métier du grand-père et où vivait-il ?
\item Quels biens le grand-père a-t-il laissés ?
\item Qui hérite de la maison ? À qui revient l'argent ?
\item Comment la famille a-t-elle résolu son désaccord ?
\item Quelle leçon le grand-père voulait-il transmettre ?
\end{enumerate}
\textbf{B. Questions de langue}
\begin{enumerate}
\item Releve dans le texte une phrase avec la structure \zh{由} et explique son rôle.
\item Releve la phrase avec \zh{虽然……但是……} et traduis-la.
\item Donne le pinyin et le sens des mots \zh{遗嘱}, \zh{继承}, \zh{团结}.
\end{enumerate}
\end{exercice}

\begin{exercice}[Thème et version type Bac]
\textbf{A. Thème.} Traduis en chinois (caractères + pinyin) :
\begin{enumerate}
\item Mon oncle est décédé il y a deux ans.
\item Il a laissé une maison et un champ.
\item Selon son testament, mon père hérite de la maison.
\item Bien que nous soyons tristes, notre famille reste unie.
\end{enumerate}
\textbf{B. Version.} Traduis en français le texte suivant :
\begin{textebox}[Version]
我的外公去年去世了。他留下了一些土地和钱。按照法律，这些财产由我的妈妈和舅舅继承。\\
\emph{Wǒ de wàigōng qùnián qùshì le. Tā liúxià le yìxiē tǔdì hé qián. Ànzhào fǎlǜ, zhèxiē cáichǎn yóu wǒ de māma hé jiùjiu jìchéng.}
\end{textebox}
\end{exercice}

\begin{exercice}[Expression écrite type Bac]
\textbf{Sujet :} Raconte le partage d'un héritage dans une famille que tu connais (réelle ou imaginaire).
\begin{itemize}
\item Longueur attendue : \textbf{6 à 10 phrases} en chinois (environ 80 à 120 caractères).
\item Tu dois utiliser \textbf{au moins trois} des mots suivants : \zh{去世}、\zh{留下}、\zh{遗嘱}、\zh{继承}、\zh{财产}.
\item Tu dois utiliser \textbf{au moins deux} connecteurs : \zh{首先}、\zh{然后}、\zh{最后}、\zh{虽然……但是……}、\zh{因为……所以……}.
\item Soigne l'écriture des caractères du thème : \zh{遗}、\zh{嘱}、\zh{财}、\zh{承}、\zh{法}.
\end{itemize}
\textbf{Entraînement à l'écriture :} avant de rédiger, recopie les caractères \zh{遗}、\zh{嘱}、\zh{财}、\zh{承}、\zh{法} en respectant l'ordre des traits (de haut en bas, de gauche à droite, l'extérieur avant l'intérieur), puis écris chaque caractère trois fois de mémoire.
\end{exercice}

\newpage

\coursec{Corrigés des exercices}

\begin{corrige}[Exercice 1 — QCM : le vocabulaire de l'héritage]
\begin{enumerate}
\item \textbf{a)} \zh{继承} (\emph{jìchéng}) = « hériter ». Attention aux homophones proches : \zh{继续} (\emph{jìxù}) signifie « continuer », \zh{纪念} (\emph{jìniàn}) signifie « commémorer ». Les trois mots commencent par la syllabe \emph{jì}, mais le second caractère change tout le sens.
\item \textbf{b)} \zh{遗嘱} (\emph{yízhǔ}) = « testament ». \zh{法律} (\emph{fǎlǜ}) = « la loi », \zh{财产} (\emph{cáichǎn}) = « les biens, le patrimoine ».
\item \textbf{c)} \zh{去世} (\emph{qùshì}) = « décéder » (formule respectueuse, littéralement « quitter le monde »). On évite \zh{死} (\emph{sǐ}, « mourir »), trop direct, quand on parle d'un proche.
\item \textbf{a)} \zh{财产} (\emph{cáichǎn}) = « bien, patrimoine ». \zh{材料} (\emph{cáiliào}) = « matériau », \zh{家庭} (\emph{jiātíng}) = « famille ».
\item \textbf{b)} Dans \zh{房子由我爸爸继承。} (\emph{Fángzi yóu wǒ bàba jìchéng.}), \zh{由} introduit \textbf{l'agent}, c'est-à-dire celui qui fait l'action : « la maison est héritée par mon père ». Structure : Complément (ce qui est transmis) + \zh{由} + agent + verbe.
\end{enumerate}
\end{corrige}

\begin{corrige}[Exercice 2 — Vrai ou faux ?]
\begin{enumerate}
\item \textbf{Vrai.} \zh{遗} contient bien la clé de la marche \zh{辶} (\emph{chuò}), qui entoure la partie \zh{贵}. On écrit d'abord l'intérieur (\zh{贵}), puis la clé \zh{辶} en dernier : pour les caractères à clé enveloppante, l'intérieur s'écrit avant l'extérieur.
\item \textbf{Faux.} \zh{遗嘱} se prononce \emph{yízhǔ} (2\textsuperscript{e} ton + 3\textsuperscript{e} ton), tandis que \zh{属于} se prononce \emph{shǔyú}. Seul le caractère \zh{属} est commun, mais sa prononciation change selon le mot : \emph{zhǔ} dans \zh{遗嘱}, \emph{shǔ} dans \zh{属于}. C'est un cas classique de caractère polyphonique.
\item \textbf{Vrai.} Contrairement au français (« bien que » et « mais » ne se combinent pas), le chinois emploie volontiers les deux ensemble : \zh{虽然……但是……} (\emph{suīrán\ldots dànshì\ldots}) = « bien que\ldots, (mais)\ldots ». Exemple : \zh{虽然我们很难过，但是我们家还是很团结。} (\emph{Suīrán wǒmen hěn nánguò, dànshì wǒmen jiā háishi hěn tuánjié.}) « Bien que nous soyons très tristes, notre famille reste très unie. »
\item \textbf{Faux.} \zh{外公} (\emph{wàigōng}) désigne le grand-père \textbf{maternel}. Le grand-père paternel se dit \zh{爷爷} (\emph{yéye}). En chinois, la parenté distingue toujours le côté paternel du côté maternel.
\item \textbf{Vrai.} \zh{两年前} (\emph{liǎng nián qián}) = « il y a deux ans ». Structure : durée + \zh{前}. Attention : devant un classificateur ou dans une durée, « deux » se dit \zh{两} et non \zh{二}.
\end{enumerate}
\end{corrige}

\begin{corrige}[Exercice 3 — Trouve le mot du lexique]
\begin{enumerate}
\item \zh{财产} (\emph{cáichǎn}) — nom — « biens, patrimoine ».
\item \zh{遗嘱} (\emph{yízhǔ}) — nom — « testament ».
\item \zh{舅舅} (\emph{jiùjiu}) — nom — « oncle maternel » (frère de la mère).
\item \zh{法律} (\emph{fǎlǜ}) — nom — « loi, législation ».
\item \zh{继承} (\emph{jìchéng}) — verbe — « hériter ».
\end{enumerate}
\textbf{Point de vigilance :} ne pas confondre \zh{舅舅} (frère de la mère) avec \zh{叔叔} (\emph{shūshu}, frère cadet du père) ni avec \zh{伯伯} (\emph{bóbo}, frère aîné du père). Le chinois possède un mot précis pour chaque lien de parenté : c'est une exigence du lexique familial au programme de Terminale.
\end{corrige}

\begin{corrige}[Exercice 4 — Associe le caractère à son pinyin]
\begin{center}
\begin{tabularx}{\linewidth}{|X|X|X|}
\hline
\rowcolor{popblueL} Caractère & Pinyin & Sens \\
\hline
\zh{遗} & yí & laisser, léguer \\
\hline
\zh{嘱} & zhǔ & recommander, ordonner \\
\hline
\zh{财} & cái & richesse, biens \\
\hline
\zh{产} & chǎn & produire, propriété \\
\hline
\zh{继} & jì & continuer, succéder \\
\hline
\end{tabularx}
\end{center}
\textbf{Astuce de mémorisation :} \zh{财} contient la clé de la coquille \zh{贝} (symbole ancien de la monnaie) ; \zh{继} contient la clé de la soie \zh{纟} (idée de continuité, de fil qui ne se rompt pas) ; \zh{嘱} contient la clé de la bouche \zh{口} (on fait des recommandations orales). Identifier la clé aide à deviner le sens et à classer le caractère dans un dictionnaire.
\end{corrige}

\begin{corrige}[Exercice 5 — Traduis en français]
\begin{enumerate}
\item \zh{我的爷爷前年去世了。} (\emph{Wǒ de yéye qiánnián qùshì le.}) — « Mon grand-père (paternel) est décédé il y a deux ans. » \zh{前年} (\emph{qiánnián}) = « l'avant-dernière année » ; \zh{了} marque l'accomplissement.
\item \zh{他留下了一所房子和一块地。} (\emph{Tā liúxià le yì suǒ fángzi hé yí kuài dì.}) — « Il a laissé une maison et un champ. » Classificateurs : \zh{所} pour les bâtiments (maison, école), \zh{块} pour les terrains et les morceaux.
\item \zh{按照他的遗嘱，房子由我爸爸继承。} (\emph{Ànzhào tā de yízhǔ, fángzi yóu wǒ bàba jìchéng.}) — « Selon son testament, la maison est héritée par mon père (revient à mon père). » \zh{按照} + nom = « selon, conformément à » ; \zh{由} + agent = tournure passive de l'agent.
\item \zh{虽然我们很难过，但是我们家还是很团结。} (\emph{Suīrán wǒmen hěn nánguò, dànshì wǒmen jiā háishi hěn tuánjié.}) — « Bien que nous soyons très tristes, notre famille reste très unie. » \zh{还是} (\emph{háishi}) = « quand même, tout de même ».
\end{enumerate}
\end{corrige}

\begin{corrige}[Exercice 6 — Transforme les phrases]
\begin{enumerate}
\item Structure \zh{把} : \zh{爸爸把房子继承了。} (\emph{Bàba bǎ fángzi jìchéng le.}) — « Papa a hérité de la maison. » Schéma : Sujet + \zh{把} + complément + verbe + \zh{了}. Le complément (\zh{房子}) passe avant le verbe.
\item Structure \zh{由} : \zh{房子由爸爸继承。} (\emph{Fángzi yóu bàba jìchéng.}) — « La maison est héritée par papa. » Schéma : Complément + \zh{由} + agent + verbe. Le complément devient thème en tête de phrase.
\item Structure \zh{被} : \zh{钱被妈妈分给了孩子们。} (\emph{Qián bèi māma fēn gěi le háizi men.}) — « L'argent a été partagé par maman entre les enfants. » Schéma : Complément + \zh{被} + agent + verbe + autres éléments.
\end{enumerate}
\textbf{Point de vigilance :} dans les trois structures, le complément est \textbf{défini} (connu des interlocuteurs) ; on ne peut pas dire \zh{把一所房子} dans ce contexte sans précision. Ne jamais laisser le verbe seul après \zh{把} : il faut un élément après lui (ici \zh{了}). La structure \zh{被} marque souvent une valeur subie ; \zh{由} est neutre et fréquent dans les textes administratifs ou familiaux, ce qui explique son emploi pour l'héritage.
\end{corrige}

\begin{corrige}[Exercice 7 — Texte à trous : les connecteurs]
Texte complété :

我的舅舅上个月去世了。\textbf{因为}我们全家都很难过，\textbf{所以}我们必须处理他的财产。\textbf{首先}，我们看了他的遗嘱；\textbf{然后}，我们开了一个家庭会议；\textbf{最后}，大家同意：房子给我的妈妈，钱给我的表哥。\textbf{虽然}事情不多，\textbf{但是}我们谈了很长时间。

\emph{Wǒ de jiùjiu shàng ge yuè qùshì le. Yīnwèi wǒmen quán jiā dōu hěn nánguò, suǒyǐ wǒmen bìxū chǔlǐ tā de cáichǎn. Shǒuxiān, wǒmen kàn le tā de yízhǔ ; ránhòu, wǒmen kāi le yí ge jiātíng huìyì ; zuìhòu, dàjiā tóngyì : fángzi gěi wǒ de māma, qián gěi wǒ de biǎogē. Suīrán shìqing bù duō, dànshì wǒmen tán le hěn cháng shíjiān.}

« Mon oncle (maternel) est décédé le mois dernier. Comme toute notre famille était très triste, nous avons dû régler ses biens. D'abord, nous avons lu son testament ; ensuite, nous avons tenu une réunion de famille ; enfin, tout le monde a accepté : la maison pour ma mère, l'argent pour mon cousin. Bien que l'affaire fût simple, nous avons longtemps discuté. »

\textbf{Justification :} \zh{因为……所以……} exprime la cause (la tristesse n'empêche pas l'obligation de régler les biens) ; \zh{首先}、\zh{然后}、\zh{最后} ordonnent les étapes chronologiques ; \zh{虽然……但是……} exprime la concession finale (affaire simple, mais longue discussion). En français, on ne traduit qu'un seul des deux connecteurs de chaque paire (« comme » seul, « bien que » seul).
\end{corrige}

\begin{corrige}[Exercice 8 — Discrimination des caractères]
\begin{enumerate}
\item 爷爷留下了很多 \textbf{\zh{财}} 产。(\emph{Yéye liúxià le hěn duō cáichǎn.}) — « Grand-père a laissé beaucoup de biens. » \zh{财} (clé \zh{贝}, la richesse), pas \zh{材} (clé \zh{木}, le bois/matériau).
\item 他的 \textbf{\zh{遗}} 嘱写得很清楚。(\emph{Tā de yízhǔ xiě de hěn qīngchu.}) — « Son testament est écrit très clairement. » \zh{遗} (clé \zh{辶}) = léguer ; \zh{遣} (\emph{qiǎn}) = envoyer, renvoyer.
\item 这所房子 \textbf{\zh{属}} 于我的爸爸。(\emph{Zhè suǒ fángzi shǔyú wǒ de bàba.}) — « Cette maison appartient à mon père. » \zh{属} = appartenir ; \zh{嘱} = recommander.
\item 这些木 \textbf{\zh{材}} 是从市场上买的。(\emph{Zhèxiē mùcái shì cóng shìchǎng shang mǎi de.}) — « Ce bois (de construction) a été acheté au marché. » \zh{材} avec la clé du bois \zh{木}.
\item 他的 \textbf{\zh{遗}} 体安葬在村子里。(\emph{Tā de yítǐ ānzàng zài cūnzi li.}) — « Sa dépouille est enterrée au village. » \zh{遗体} = « corps (d'un défunt) ».
\end{enumerate}
\textbf{Point de vigilance :} \zh{遗} et \zh{遣} se ressemblent (même clé enveloppante \zh{辶}) mais diffèrent par leur partie intérieure : \zh{贵} dans \zh{遗}, un composant différent dans \zh{遣} ; \zh{嘱} et \zh{属} partagent le phonétique \zh{属}, mais \zh{嘱} ajoute la clé de la bouche \zh{口} à gauche. Méthode : repérer d'abord la clé (le sens), puis le phonétique (le son).
\end{corrige}

\begin{corrige}[Exercice 9 — Remets les mots dans l'ordre]
\begin{enumerate}
\item \zh{房子由我的叔叔继承。} (\emph{Fángzi yóu wǒ de shūshu jìchéng.}) — « La maison est héritée par mon oncle (paternel). » Ordre : complément-thème + \zh{由} + agent + verbe.
\item \zh{我的外公去年去世了。} (\emph{Wǒ de wàigōng qùnián qùshì le.}) — « Mon grand-père maternel est décédé l'année dernière. » Le complément de temps \zh{去年} se place après le sujet, avant le verbe.
\item \zh{大家按照他的遗嘱分了财产。} (\emph{Dàjiā ànzhào tā de yízhǔ fēn le cáichǎn.}) — « Tout le monde a partagé les biens selon son testament. » Le complément prépositionnel \zh{按照他的遗嘱} se place avant le verbe \zh{分}.
\end{enumerate}
\textbf{Rappel d'ordre des mots :} Sujet + (temps) + (complément prépositionnel) + verbe + complément. En chinois, tout ce qui encadre l'action (temps, manière, moyen) se place \textbf{avant} le verbe, jamais après comme en français.
\end{corrige}

\begin{corrige}[Exercice 10 — Compréhension de l'écrit et questions de langue]
\textbf{A. Questions de compréhension}
\begin{enumerate}
\item Le grand-père était \textbf{paysan (agriculteur)} et vivait à \textbf{Douala} (\zh{杜阿拉}, \emph{Dù'ālā}).
\item Il a laissé \textbf{une maison, deux champs et un peu d'argent} (\zh{一所房子、两块地和一些钱}).
\item La maison revient au \textbf{père} du narrateur (\zh{房子由我爸爸继承}) ; l'argent est laissé à la \textbf{grand-mère} (\zh{钱留给我的奶奶}).
\item La famille a tenu \textbf{une réunion de famille} (\zh{开了一个家庭会议}), après laquelle chacun a accepté.
\item Il voulait transmettre que \textbf{les biens ne sont pas importants, c'est l'union de la famille qui compte le plus} (\zh{财产不重要，家庭团结最重要。}).
\end{enumerate}
\textbf{B. Questions de langue}
\begin{enumerate}
\item \zh{房子由我爸爸继承。} (\emph{Fángzi yóu wǒ bàba jìchéng.}) — \zh{由} introduit l'\textbf{agent} de l'action : la maison (thème, en tête) est héritée par le père. C'est une tournure passive neutre, très fréquente dans les textes administratifs ou familiaux.
\item \zh{虽然大家开始有一点不同的意见，但是开了一个家庭会议以后，每个人都同意了。} (\emph{Suīrán dàjiā kāishǐ yǒu yìdiǎn bùtóng de yìjian, dànshì kāi le yí ge jiātíng huìyì yǐhòu, měi ge rén dōu tóngyì le.}) — « Bien qu'au début chacun ait eu des avis un peu différents, après avoir tenu une réunion de famille, tout le monde a accepté. »
\item \zh{遗嘱} \emph{yízhǔ} : « testament » ; \zh{继承} \emph{jìchéng} : « hériter » ; \zh{团结} \emph{tuánjié} : « uni, solidaire ; l'union ».
\end{enumerate}
\textbf{Barème indicatif (sur 10) :} A : 5 pts (1 pt par réponse correcte et complète) ; B : 5 pts (relevé + explication 2 pts, traduction 1,5 pt, pinyin et sens 1,5 pt).
\end{corrige}

\begin{corrige}[Exercice 11 — Thème et version type Bac]
\textbf{A. Thème (propositions de correction)}
\begin{enumerate}
\item \zh{我的舅舅两年前去世了。} (\emph{Wǒ de jiùjiu liǎng nián qián qùshì le.}) — « il y a deux ans » = \zh{两年前}, placé avant le verbe ; « oncle » ici : \zh{舅舅} (maternel) ou \zh{叔叔} (paternel) sont acceptés.
\item \zh{他留下了一所房子和一块地。} (\emph{Tā liúxià le yì suǒ fángzi hé yí kuài dì.}) — classificateurs obligatoires : \zh{所} pour la maison, \zh{块} pour le champ.
\item \zh{按照他的遗嘱，房子由我爸爸继承。} (\emph{Ànzhào tā de yízhǔ, fángzi yóu wǒ bàba jìchéng.}) — ou, avec la structure \zh{把} : \zh{按照他的遗嘱，我爸爸把房子继承了。}
\item \zh{虽然我们很难过，但是我们的家庭还是很团结。} (\emph{Suīrán wǒmen hěn nánguò, dànshì wǒmen de jiātíng háishi hěn tuánjié.}) — ne pas oublier \zh{但是} dans la seconde proposition.
\end{enumerate}
\textbf{B. Version}

\zh{我的外公去年去世了。他留下了一些地和钱。按照法律，这些财产由我的妈妈和舅舅继承。} (\emph{Wǒ de wàigōng qùnián qùshì le. Tā liúxià le yìxiē dì hé qián. Ànzhào fǎlǜ, zhèxiē cáichǎn yóu wǒ de māma hé jiùjiu jìchéng.})

« Mon grand-père maternel est décédé l'année dernière. Il a laissé des terres et de l'argent. Selon la loi, ces biens sont hérités par ma mère et mon oncle (reviennent à ma mère et à mon oncle). »

\textbf{Barème indicatif (sur 10) :} Thème : 6 pts (1,5 pt par phrase : vocabulaire 0,5, structure 0,5, tons/classificateurs 0,5) ; Version : 4 pts (fidélité 2 pts, qualité du français 2 pts).

\textbf{Points de vigilance :} place du complément de temps avant le verbe ; classificateurs \zh{所} et \zh{块} ; \zh{两} et non \zh{二} devant \zh{年} dans \zh{两年前} ; « selon » se traduit ici par \zh{按照} (et non par une préposition française calquée) ; dans la version, \zh{由} se rend naturellement par « revenir à » ou par une passive française.
\end{corrige}

\begin{corrige}[Exercice 12 — Expression écrite type Bac]
\textbf{Proposition de rédaction modèle}

去年，我的外公去世了，我们全家都很难过。他留下了一所房子、一块地和一些钱。按照他的遗嘱，房子由我的妈妈继承，地由我的舅舅继承，钱留给我的外婆。首先，我们开了一个家庭会议；然后，大家按照遗嘱分了财产；最后，每个人都同意了。虽然我们失去了外公，但是我们的家庭还是很团结。我觉得财产不重要，家庭的爱最重要。

\emph{Qùnián, wǒ de wàigōng qùshì le, wǒmen quán jiā dōu hěn nánguò. Tā liúxià le yì suǒ fángzi, yí kuài dì hé yìxiē qián. Ànzhào tā de yízhǔ, fángzi yóu wǒ de māma jìchéng, dì yóu wǒ de jiùjiu jìchéng, qián liú gěi wǒ de wàipó. Shǒuxiān, wǒmen kāi le yí ge jiātíng huìyì ; ránhòu, dàjiā ànzhào yízhǔ fēn le cáichǎn ; zuìhòu, měi ge rén dōu tóngyì le. Suīrán wǒmen shīqù le wàigōng, dànshì wǒmen de jiātíng háishi hěn tuánjié. Wǒ juéde cáichǎn bù zhòngyào, jiātíng de ài zuì zhòngyào.}

« L'année dernière, mon grand-père maternel est décédé et toute notre famille était très triste. Il a laissé une maison, un champ et un peu d'argent. Selon son testament, la maison revient à ma mère, le champ à mon oncle, l'argent est laissé à ma grand-mère. D'abord, nous avons tenu une réunion de famille ; ensuite, chacun a partagé les biens selon le testament ; enfin, tout le monde a accepté. Bien que nous ayons perdu grand-père, notre famille reste très unie. Je pense que les biens ne sont pas importants : l'amour de la famille est le plus important. »

\textbf{Respect des contraintes :} 8 phrases (environ 110 caractères) ; mots du thème utilisés : \zh{去世}、\zh{留下}、\zh{遗嘱}、\zh{继承}、\zh{财产} (5 sur 3 exigés) ; connecteurs : \zh{首先}、\zh{然后}、\zh{最后}、\zh{虽然……但是……} (4 sur 2 exigés).

\textbf{Barème indicatif (sur 20) :} respect du sujet et de la longueur 4 pts ; mots du lexique imposés 4 pts ; connecteurs 4 pts ; correction grammaticale (ordre des mots, \zh{了}, classificateurs) 4 pts ; écriture des caractères 2 pts ; cohérence et touche personnelle 2 pts.

\textbf{Points de vigilance :}
\begin{itemize}
\item Ordre des mots : le temps (\zh{去年}) avant le verbe ; jamais \zh{去世了去年}.
\item \zh{由} + agent : \zh{房子由我的妈妈继承}, pas d'inversion.
\item Classificateurs : \zh{一所房子}、\zh{一块地}.
\item Caractères : \zh{遗} (intérieur \zh{贵} d'abord, clé \zh{辶} en dernier) ; \zh{嘱} (clé \zh{口} à gauche) ; \zh{财} (clé \zh{贝} à gauche) ; \zh{承} (caractère symétrique, trait vertical central d'abord) ; \zh{法} (clé de l'eau \zh{氵} à gauche, trois traits de haut en bas).
\item Conclure par une phrase d'opinion (\zh{我觉得……}) valorise la copie.
\end{itemize}
\end{corrige}

\newpage

\coursec{Fiche bilan}

\begin{popfigure}
\begin{tikzpicture}[
  mindmap,
  concept color=popblue,
  level 1/.append style={level distance=3.8cm, sibling angle=360/7},
  level 2/.append style={level distance=2.5cm},
  every node/.style={font=\small, align=center},
  concept/.append style={minimum size=2.2cm, inner sep=4pt},
  grow cyclic
]
\node[concept, align=center] (center){\textbf{Leçon 7}\\\textbf{L'héritage}\\\zh{传承 / 遗产}}
  child[concept color=popblue] {node[concept, align=center]{\textbf{Vocabulaire}\\\zh{重点词汇}}}
  child[concept color=poporange] {node[concept, align=center]{\textbf{Caractères}\\\zh{重点汉字}}}
  child[concept color=popgreen] {node[concept, align=center]{\textbf{Connecteurs}\\\zh{连接词}}}
  child[concept color=poppurple] {node[concept, align=center]{\textbf{Structures}\\\zh{语法结构}}}
  child[concept color=poppink] {node[concept, align=center]{\textbf{Modèle texte}\\\zh{范文结构}}}
  child[concept color=popgold] {node[concept, align=center]{\textbf{Thème/Version}\\\zh{翻译练习}}}
  child[concept color=popmuted] {node[concept, align=center]{\textbf{Culture}\\\zh{中喀对比}}};
\end{tikzpicture}
\legende{Carte mentale de la leçon 7 : Expression écrite — L'héritage}
\end{popfigure}

\begin{bilan}

\end{bilan}

\courssub{Lexique clé du thème « Héritage / Transmission » (HSK 3–4)}

\begin{center}
\begin{tabularx}{\linewidth}{|X|X|X|X|}
\hline
\rowcolor{popblueL} \textbf{Caractères} & \textbf{Pinyin} & \textbf{Nature} & \textbf{Traduction} \\ \hline
传承 & chuánchéng & verbe & transmettre, perpétuer (héritage) \\ \hline
遗产 & yíchǎn & nom & héritage, patrimoine \\ \hline
非物质文化遗产 & fēiwùzhì wénhuà yíchǎn & nom & patrimoine culturel immatériel (PCI) \\ \hline
传统 & chuántǒng & nom/adj & tradition, traditionnel \\ \hline
文化 & wénhuà & nom & culture \\ \hline
一代 & yídài & nom & une génération \\ \hline
传递 & chuándì & verbe & transmettre, passer (relais, info) \\ \hline
保护 & bǎohù & verbe/nom & protéger, protection \\ \hline
口述 & kǒushù & adj/verbe & oral, transmis oralement \\ \hline
文字 & wénzì & nom & écriture, caractères, texte écrit \\ \hline
祖先 & zǔxiān & nom & ancêtres \\ \hline
后代 & hòudài & nom & descendants, générations futures \\ \hline
珍贵 & zhēnguì & adj & précieux \\ \hline
使命 & shǐmìng & nom & mission \\ \hline
责任 & zérèn & nom & responsabilité \\ \hline
\end{tabularx}
\end{center}

\courssub{Règles de grammaire et connecteurs à connaître par cœur}

\begin{center}
\begin{tabularx}{\linewidth}{|X|X|X|}
\hline
\rowcolor{poporangeL} \textbf{Notion / Fonction} & \textbf{Structure / Schéma} & \textbf{Exemple (Caractères + Pinyin + Traduction)} \\ \hline
\textbf{Connecteurs : Énumération} & 首先…，其次…，最后… / 总之 & \zh{首先，我们要保护非遗；其次，要传给后代；总之，这是责任。} (Shǒuxiān… qícì… zuìhòu… / zǒngzhī…) — D'abord…, ensuite…, enfin… / En résumé… \\ \hline
\textbf{Connecteurs : Addition} & 此外 / 而且 / 另外 & \zh{此外，学校也开设了相关课程。} (Cǐwài…) — De plus / En outre… \\ \hline
\textbf{Connecteurs : Cause/Effet} & 因为…所以… / 由于…因此… & \zh{因为非遗珍贵，所以我们必须保护它。} (Yīnwèi… suǒyǐ… / Yóuyú… yīncǐ…) — Parce que…, donc… \\ \hline
\textbf{把字句 (Action sur objet)} & Sujet + \cle{把} + Objet + Verbe + Complément de résultat/direction & \zh{我们要把传统文化传承下去。} (Wǒmen yào \textbf{bǎ} chuántǒng wénhuà chuánchéng xiàqù.) — Nous devons transmettre la culture traditionnelle. \\ \hline
\textbf{被字句 (Passif formel)} & Sujet (patient) + \cle{被} + Agent + Verbe + Complément & \zh{这个非遗项目被列入了国家名录。} (Zhège fēiyí xiàngmù \textbf{bèi} lièrùle guójiā mínglù.) — Ce projet PCI a été inscrit au registre national. \\ \hline
\textbf{Causatif : 让 / 使 / 令} & Sujet + \cle{让/使/令} + Objet + Verbe/Adjectif & \zh{这件事让年轻人了解了历史。} (Zhè jiàn shì \textbf{ràng} niánqīngrén liǎojiěle lìshǐ.) — Cette affaire a fait que les jeunes connaissent l'histoire. \\ \hline
\textbf{Comparatif 比} & A + \cle{比} + B + Adj (pas de 很) & \zh{现在比过去更重视非遗保护。} (Xiànzài \textbf{bǐ} guòqù gèng zhòngshì fēiyí bǎohù.) — Maintenant on attache plus d'importance qu'avant. \\ \hline
\textbf{Structure « 由…组成/构成 »} & Sujet + \cle{由} + Élément + 组成/构成 & \zh{喀麦隆的文化由多个民族的传统组成。} (Kāmàilóng de wénhuà \textbf{yóu} duōgè mínzú de chuántǒng zǔchéng.) — La culture camerounaise est composée des traditions de multiples ethnies. \\ \hline
\end{tabularx}
\end{center}

\courssub{Méthodes express}

\begin{itemize}
\item \textbf{Rédaction guidée (150–200 caractères) :}
      \begin{enumerate}
      \item \textbf{Analyser} le sujet : mots-clés (héritage, Cameroun, Chine, jeunesse, responsabilité).
      \item \textbf{Plan en 3 parties} : Introduction (accroche + définition + thèse) → Développement (2 arguments : ex. rôle de la famille/école + exemples concrets Ndolé/Opéra de Pékin) → Conclusion (bilan + ouverture sur l'avenir).
      \item \textbf{Mobiliser} : 5 connecteurs minimum, 2 structures \cle{把}/\cle{被}, vocabulaire précis (PCI, 传承, 非遗).
      \item \textbf{Écrire} : phrases courtes, ordre SVO strict, vérification des tons sur le brouillon.
      \item \textbf{Relire} : accords classificateurs, aspect \cle{了}/\cle{过}, ponctuation chinoise pleine chasse.
      \end{enumerate}
\item \textbf{Thème (FR → CN) :}
      \begin{enumerate}
      \item Segmenter la phrase française (Sujet / Verbe / COD / Compléments).
      \item Identifier la structure cible (把, 被, 比, 让, 由).
      \item Traduire les mots-clés (dictionnaire mental).
      \item Assembler dans l'ordre chinois (Sujet + Temps/Lieu + Manière + Verbe + Complément).
      \item Vérifier : aspect, classificateurs, connecteurs.
      \end{enumerate}
\item \textbf{Version (CN → FR) :}
      \begin{enumerate}
      \item Segmenter en groupes syntaxiques (repérer \cle{的}, \cle{了}, \cle{把}, \cle{被}, verbes).
      \item Traduire le noyau de phrase (Sujet + Verbe).
      \item Replacer les compléments (ordre logique français).
      \item Ajuster le temps/aspect (passé/présent/futur, accompli/durée).
      \item Polir le français (style naturel, pas de calque).
      \end{enumerate}
\end{itemize}



\begin{aretenir}[Checklist avant le Bac — Leçon 7 : L'héritage]
$\square$ Je connais par cœur les 15 mots du lexique (caractères, pinyin tons corrects, nature, sens). \\
$\square$ J'écris correctement les 10 caractères clés (ordre des traits, radicaux : \zh{辶} dans 传/遗, \zh{贝} dans 产/贵, \zh{人} dans 代/传). \\
$\square$ Je distingue \zh{传承} (transmettre un héritage abstrait) de \zh{传递} (passer un objet/info concret). \\
$\square$ J'utilise 5 connecteurs logiques pour structurer un paragraphe (\zh{首先/其次/此外/因此/总之}). \\
$\square$ Je construis une phrase \cle{把} correcte : Sujet + 把 + Objet connu + V + Complément (ex. \zh{把故事讲给孩子听}). \\
$\square$ Je construis une phrase \cle{被} correcte : Patient + 被 + Agent + V + Complément (ex. \zh{非遗被年轻人保护}). \\
$\square$ J'emploie \cle{让/使} pour exprimer la causalité dans le contexte éducatif (ex. \zh{让后代了解历史}). \\
$\square$ Je rédige une introduction type : Accroche (citation/proverbe) + Définition (PCI) + Thèse + Plan annoncé. \\
$\square$ Je développe deux arguments distincts avec exemples Cameroun (chefferies, ndolé, danses) et Chine (opéra, calligraphie, fêtes). \\
$\square$ Je conclus par un bilan + une ouverture (rôle de l'école, numérique, responsabilité jeunesse). \\
$\square$ Je réussis un thème de 5 phrases FR→CN sans erreur d'ordre des mots ni d'aspect verbal. \\
$\square$ Je réussis une version CN→FR fidèle, fluide, sans calque syntaxique.
\end{aretenir}

\findecours

\end{document}
